\documentclass[12pt,a4paper]{report}
% Impostazioni lingua e codifica
\usepackage[utf8]{inputenc}
\usepackage[T1]{fontenc}
\usepackage{lmodern}
\usepackage[italian]{babel}
\usepackage{microtype}  % ✅ Migliora kerning, spaziatura e leggibilità automatica

% Matematica e simboli
\usepackage{amsmath, amssymb, amsthm}
\usepackage{mathtools}
\usepackage{siunitx}    % ✅ Formattazione coerente di unità e numeri
\sisetup{locale = IT, separate-uncertainty = true, inter-unit-product = \ensuremath{{}\cdot{}}}

% Figure, tabelle e liste
\usepackage{graphicx}
\usepackage{caption}
\usepackage{subcaption}
\usepackage{booktabs}   % ✅ Tabelle professionali
\usepackage{array}
\usepackage{enumitem}   % ✅ Liste compatte e senza spazi vuoti errati
\setlist[itemize,enumerate]{leftmargin=*, noitemsep, topsep=2pt plus 1pt minus 1pt}

% Impostazioni pagina e collegamenti
\usepackage{geometry}
\geometry{margin=1.5cm, headheight=15pt, footskip=20pt}
\usepackage{hyperref}
\hypersetup{
  colorlinks=true, linkcolor=blue, filecolor=magenta, urlcolor=cyan,
  pdfborder={0 0 0}, pdfauthor={Dispensa Universitaria - Enrico Desideri},
  pdftitle={Biochimica}
}

% Ambienti personalizzati
\newtheorem{definizione}{Definizione}[section]
\newtheorem{teorema}{Teorema}[section]
\newtheorem{proposizione}{Proposizione}[section]
\newtheorem{esempio}{Esempio}[section]
\newtheorem{nota}{Nota}[section]

\begin{document}

\title{\textbf{Biochimica}}
\author{}
\date{}
\maketitle
\phantomsection
\tableofcontents
\clearpage

% =========================================================
% CAPITOLO 1: Organizzazione della Materia
% =========================================================
\chapter{Organizzazione della Materia}
\section{La materia e gli elementi chimici}
La \textbf{materia} è definita come qualunque entità che occupi uno spazio e possieda una massa. Essa è costituita da \textbf{elementi chimici} o da combinazioni di essi. Attualmente sono noti 118 elementi, di cui 92 presenti naturalmente sulla Terra; i restanti sono stati sintetizzati artificialmente in laboratorio. Ciascun elemento è identificato da un simbolo composto da una o due lettere (ad esempio, \(\mathrm{O}\) per l'ossigeno, \(\mathrm{Na}\) per il sodio).
Nel corpo umano sono presenti meno di 30 dei 92 elementi naturali, distribuiti in proporzioni molto diverse. Quattro elementi --- ossigeno (\(\mathrm{O}\)), carbonio (\(\mathrm{C}\)), idrogeno (\(\mathrm{H}\)) e azoto (\(\mathrm{N}\)) --- rappresentano oltre il \SI{96}{\percent} del peso corporeo. La maggior parte dell'ossigeno è contenuta nell'acqua (\(\mathrm{H_2O}\)), che costituisce il 50--70\% della massa corporea totale.
Altri elementi presenti in quantità significative includono:
\begin{itemize}
\item Calcio (\(\mathrm{Ca}\)): \SI{1,5}{\percent}
\item Fosforo (\(\mathrm{P}\)): \SI{1}{\percent}
\item Potassio (\(\mathrm{K}\)): \SI{0,35}{\percent}
\item Zolfo (\(\mathrm{S}\)): \SI{0,25}{\percent}
\item Sodio (\(\mathrm{Na}\)): \SI{0,2}{\percent}
\item Cloro (\(\mathrm{Cl}\)): \SI{0,2}{\percent}
\item Ferro (\(\mathrm{Fe}\)): \SI{0,005}{\percent}
\end{itemize}
Il restante \SI{0,2}{\percent} è costituito da \textbf{elementi in tracce}, principalmente metalli come rame (\(\mathrm{Cu}\)), zinco (\(\mathrm{Zn}\)) e manganese (\(\mathrm{Mn}\)), essenziali per funzioni enzimatiche e strutturali.
Nella tavola periodica, gli elementi sono organizzati in \textbf{gruppi} (colonne) e \textbf{periodi} (righe). Elementi appartenenti allo stesso gruppo condividono proprietà chimiche simili, poiché presentano lo stesso numero di elettroni di valenza.

\section{Struttura dell'atomo}
L'\textbf{atomo} è la più piccola unità di un elemento che ne conserva le proprietà chimiche. Esso può esistere isolato o legarsi ad altri atomi, uguali o diversi, per formare molecole.
Ogni atomo è costituito da:
\begin{itemize}
\item \textbf{Nucleo centrale}, contenente:
\begin{itemize}
    \item \textit{Protoni} (\(p^+\)): particelle con carica elettrica positiva;
    \item \textit{Neutroni} (\(n^0\)): particelle neutre, di massa simile ai protoni.
\end{itemize}
\item \textbf{Elettroni} (\(e^-\)): particelle con carica negativa, di massa circa 2000 volte inferiore a quella di protoni e neutroni, che si muovono in regioni dello spazio dette \textit{orbitali elettronici}.
\end{itemize}
Gli orbitali sono organizzati in livelli energetici progressivamente più esterni. Il primo orbitale (più vicino al nucleo) può contenere al massimo 2 elettroni, il secondo fino a 8, il terzo fino a 18, e così via.
\begin{definizione}[Numero atomico]
Il \textbf{numero atomico} (\(Z\)) di un elemento è il numero di protoni presenti nel nucleo di ciascun suo atomo. Poiché in un atomo neutro il numero di elettroni eguaglia quello dei protoni, il numero atomico determina anche la configurazione elettronica e, di conseguenza, le proprietà chimiche dell'elemento.
\end{definizione}
Ad esempio:
\begin{itemize}
\item Idrogeno (\(\mathrm{H}\)): \(Z = 1\) (1 protone, 1 elettrone);
\item Ossigeno (\(\mathrm{O}\)): \(Z = 8\) (8 protoni, 8 elettroni);
\item Uranio (\(\mathrm{U}\)): \(Z = 92\) (elemento naturale con il più alto numero atomico).
\end{itemize}

\section{Isotopi e massa atomica}
Tutti gli atomi di uno stesso elemento condividono il numero atomico, ma possono differire per il numero di neutroni nel nucleo. La somma di protoni e neutroni è detta \textbf{numero di massa} (\(A\)).
\begin{definizione}[Isotopi]
Gli \textbf{isotopi} sono atomi dello stesso elemento (stesso \(Z\)) con diverso numero di massa (\(A\)), ovvero con differente numero di neutroni.
\end{definizione}
Un esempio classico è rappresentato dagli isotopi dell'idrogeno:
\begin{itemize}
\item \textbf{Prozio} (\(^1\mathrm{H}\)): 1 protone, 0 neutroni (isotopo più abbondante);
\item \textbf{Deuterio} (\(^2\mathrm{H}\)): 1 protone, 1 neutrone (stabile);
\item \textbf{Trizio} (\(^3\mathrm{H}\)): 1 protone, 2 neutroni (radioisotopo instabile).
\end{itemize}
La maggior parte degli isotopi è stabile; alcuni, detti \textbf{radioisotopi}, sono instabili e decadono emettendo radiazioni. Gli isotopi di un elemento condividono le stesse proprietà chimiche, poiché la configurazione elettronica --- determinante per le reazioni chimiche --- è identica.
Nella tavola periodica, per ciascun elemento è riportata la \textbf{massa atomica relativa}, espressa in \textbf{unità di massa atomica} (u.m.a.) o \textbf{dalton} (\(\mathrm{Da}\)). Un dalton è definito come \(1/12\) della massa di un atomo di carbonio-12 (\(^{12}\mathrm{C}\)). La massa atomica indicata è la media ponderata delle masse di tutti gli isotopi naturali, in base alla loro abbondanza relativa.

\section{Ioni e ionizzazione}
Un atomo neutro può acquisire o perdere uno o più elettroni, trasformandosi in uno \textbf{ione} carico elettricamente.
\begin{itemize}
\item \textbf{Catione}: ione con carica positiva, formato per perdita di elettroni (es. \(\mathrm{Ca^{2+}}\): calcio che ha perso 2 elettroni);
\item \textbf{Anione}: ione con carica negativa, formato per acquisizione di elettroni (es. \(\mathrm{Cl^-}\): cloro che ha acquisito 1 elettrone).
\end{itemize}
La ionizzazione è un processo fondamentale in molte reazioni biochimiche, come la formazione di sali e il trasporto ionico attraverso le membrane cellulari.

\section{Legami chimici}
Gli atomi si legano tra loro per formare molecole, raggiungendo configurazioni elettroniche più stabili (a minore energia potenziale). La formazione o rottura di legami chimici comporta sempre scambi di energia.
\subsection{Legame covalente}
Nel \textbf{legame covalente}, due atomi condividono una o più coppie di elettroni. È il legame più forte tra quelli rilevanti in biologia.
\begin{definizione}[Valenza]
La \textbf{valenza} di un atomo è il numero di elettroni che esso può mettere in condivisione per formare legami covalenti.
\end{definizione}
Esempi:
\begin{itemize}
\item Idrogeno (\(\mathrm{H}\)): valenza 1;
\item Ossigeno (\(\mathrm{O}\)): valenza 2;
\item Carbonio (\(\mathrm{C}\)): valenza 4.
\end{itemize}
A seconda del numero di coppie condivise, si distinguono:
\begin{itemize}
\item \textbf{Legame singolo}: condivisione di 1 coppia di elettroni;
\item \textbf{Legame doppio}: condivisione di 2 coppie;
\item \textbf{Legame triplo}: condivisione di 3 coppie.
\end{itemize}
Maggiore è il numero di elettroni condivisi, maggiore è l'\textbf{energia di legame} (espressa in \(\mathrm{kJ\,mol^{-1}}\) o \(\mathrm{kcal\,mol^{-1}}\)), ovvero l'energia necessaria per rompere il legame.
\subsubsection{Polarità del legame covalente}
La distribuzione degli elettroni condivisi dipende dall'\textbf{elettronegatività}, ovvero la capacità di un atomo di attrarre elettroni.
\begin{itemize}
\item \textbf{Legame covalente non polare}: gli atomi hanno elettronegatività simile; gli elettroni sono equamente condivisi (es. \(\mathrm{O_2}\), \(\mathrm{CH_4}\)).
\item \textbf{Legame covalente polare}: differenza di elettronegatività \(0 < \Delta e < 1.9\); l'atomo più elettronegativo acquisisce una parziale carica negativa (\(\delta^-\)), l'altro una parziale carica positiva (\(\delta^+\)). Esempio: legame \(\mathrm{O-H}\) nell'acqua.
\end{itemize}
\begin{definizione}[Molecola polare]
Una molecola è \textbf{polare} quando presenta una separazione permanente di carica dovuta a legami polari e a una geometria asimmetrica. Al contrario, una molecola \textbf{apolare} ha distribuzione simmetrica degli elettroni.
\end{definizione}
La polarità influenza la solubilità: le molecole polari sono \textbf{idrofiliche} (si sciolgono in acqua), mentre quelle apolari sono \textbf{idrofobiche} (es. lipidi, oli).
\subsection{Legame ionico}
Quando la differenza di elettronegatività tra due atomi è elevata (\(\Delta e > 1.9\)), un atomo trasferisce completamente uno o più elettroni all'altro. Si formano ioni di carica opposta che si attraggono elettrostaticamente, dando origine al \textbf{legame ionico}.
Esempio: nel cloruro di sodio (\(\mathrm{NaCl}\)), il sodio cede un elettrone al cloro, formando \(\mathrm{Na^+}\) e \(\mathrm{Cl^-}\) che si organizzano in un \textit{reticolo cristallino}. La carica complessiva del composto è sempre nulla.
\subsection{Legame idrogeno}
Il \textbf{legame idrogeno} (o ponte a idrogeno) è un'interazione debole (10--20 volte meno energetica di un legame covalente) che si instaura tra molecole polari contenenti idrogeno legato a un atomo altamente elettronegativo (ossigeno, azoto, fluoro).
Si forma quando la parziale carica positiva dell'idrogeno (\(\delta^+\)) interagisce con la parziale carica negativa (\(\delta^-\)) di un'altra molecola. Nell'acqua, ogni molecola può formare fino a 4 legami idrogeno, creando una rete dinamica responsabile di proprietà uniche come coesione, adesione e alto calore specifico.
Nonostante la debolezza individuale, la grande quantità di legami idrogeno nelle macromolecole biologiche (proteine, DNA) conferisce stabilità strutturale e, al contempo, flessibilità dinamica, essenziale per la funzione biologica.

\section{Molecole, gruppi funzionali e composti organici}
Le molecole sono aggregati di atomi legati covalentemente. Quando contengono atomi di due o più elementi diversi, sono dette \textbf{composti}.
\subsection{Gruppi funzionali}
I \textbf{gruppi funzionali} sono specifici raggruppamenti di atomi che conferiscono caratteristiche chimiche definite alle molecole in cui sono presenti. Nei composti biologici, i più comuni sono:
\begin{itemize}
\item \textbf{Carbossilico} (\(-\mathrm{COOH}\)): acido, presente negli amminoacidi;
\item \textbf{Aminico} (\(-\mathrm{NH_2}\)): basico, presente negli amminoacidi;
\item \textbf{Idrossilico} (\(-\mathrm{OH}\)): polare, presente in zuccheri e alcoli;
\item \textbf{Fosfato} (\(-\mathrm{PO_4^{2-}}\)): coinvolto nel trasferimento di energia (ATP) e nella fosforilazione;
\item \textbf{Carbonilico} (\(>\mathrm{C=O}\)): presente in aldeidi e chetoni (es. monosaccaridi);
\item \textbf{Sulfidrilico} (\(-\mathrm{SH}\)): importante per la struttura terziaria delle proteine (ponti disolfuro nella cisteina).
\end{itemize}
\subsection{Composti organici e scheletro carbonioso}
I \textbf{composti organici} sono definiti come molecole contenenti atomi di carbonio (\(\mathrm{C}\)), con eccezione di ossidi del carbonio e carbonati. Le quattro classi di macromolecole biologiche --- carboidrati, lipidi, proteine e acidi nucleici --- sono tutte composti organici.
Il carbonio, con valenza 4, forma legami covalenti stabili con sé stesso e con altri elementi, dando origine a:
\begin{itemize}
\item Catene lineari;
\item Strutture ramificate;
\item Anelli ciclici.
\end{itemize}
Questa versatilità è alla base della diversità strutturale delle molecole biologiche.
\subsection{Notazioni molecolari}
Per rappresentare le molecole si utilizzano diverse convenzioni:
\begin{itemize}
\item \textbf{Formula molecolare (bruta)}: indica il tipo e il numero di atomi (es. glucosio: \(\mathrm{C_6H_{12}O_6}\)). Permette il calcolo del peso molecolare:
\[
\text{PM}_{\text{glucosio}} = (12 \times 6) + (1 \times 12) + (16 \times 6) = 180\,\mathrm{Da}
\]
\item \textbf{Formula di struttura}: mostra come gli atomi sono collegati tra loro, fornendo informazioni sulla connettività e sulla geometria.
\end{itemize}
\begin{definizione}[Isomeri]
Gli \textbf{isomeri} sono molecole con la stessa formula molecolare ma diversa formula di struttura o diversa disposizione spaziale degli atomi.
\end{definizione}
\subsubsection{Isomeria strutturale}
Gli isomeri strutturali differiscono per il numero o il tipo di legami tra gli atomi:
\begin{itemize}
\item \textbf{Isomeri di catena}: differenze nello scheletro carbonioso (es. leucina e isoleucina);
\item \textbf{Isomeri di posizione}: diversa posizione di gruppi funzionali o legami multipli (es. acidi grassi insaturi).
\end{itemize}
\subsubsection{Stereoisomeria}
Gli stereoisomeri hanno gli stessi legami ma diversa disposizione tridimensionale:
\begin{itemize}
\item \textbf{Isomeri geometrici (cis-trans)}: in presenza di doppi legami \(\mathrm{C=C}\), i sostituenti possono trovarsi dallo stesso lato (\textit{cis}) o da lati opposti (\textit{trans}). Il doppio legame impedisce la rotazione, "bloccando" la conformazione.
\item \textbf{Isomeri ottici (enantiomeri)}: molecole che sono immagini speculari non sovrapponibili. Si verificano quando un atomo di carbonio (detto \textit{centro chirale}) è legato a quattro gruppi diversi. Gli enantiomeri sono indicati come \textit{D} (destrogiro) e \textit{L} (levogiro).
\end{itemize}
\begin{nota}[Chiralità in biologia]
In sistemi biologici, gli enantiomeri non sono intercambiabili. Nelle proteine, tutti gli amminoacidi sono nella forma \textit{L}; negli zuccheri metabolici, prevale la forma \textit{D} (es. \textit{D}-glucosio). Molti enzimi mostrano specificità stereochimica, riconoscendo solo uno dei due enantiomeri.
\end{nota}

\section{Reazioni chimiche ed energia}
Una \textbf{reazione chimica} comporta la rottura di legami nei reagenti e la formazione di nuovi legami nei prodotti. L'insieme di tutte le reazioni chimiche in un organismo costituisce il \textbf{metabolismo}.
\subsection{Conservazione della massa e dell'energia}
\begin{itemize}
\item \textbf{Legge di conservazione della massa}: il numero totale di atomi di ciascun elemento è identico tra reagenti e prodotti.
\item \textbf{Legge di conservazione dell'energia}: l'energia non può essere creata né distrutta, ma solo trasformata da una forma all'altra.
\end{itemize}
L'energia immagazzinata nei legami chimici è \textbf{energia potenziale chimica}, che può essere convertita in \textbf{energia cinetica} (movimento) o \textbf{energia termica} (calore).
\subsection{Unità di misura dell'energia}
\begin{itemize}
\item \textbf{Joule} (\(\mathrm{J}\)): unità SI; \(1\,\mathrm{J} = 1\,\mathrm{N \cdot m}\);
\item \textbf{Caloria} (\(\mathrm{cal}\)): calore necessario per innalzare di \(1^\circ\mathrm{C}\) la temperatura di \(1\,\mathrm{g}\) di acqua;
\item Relazione: \(1\,\mathrm{cal} = 4.184\,\mathrm{J}\); \(1\,\mathrm{kcal} = 1000\,\mathrm{cal} = 4.184\,\mathrm{kJ}\).
\end{itemize}
Un adulto consuma circa \(2000\,\mathrm{kcal}\) al giorno, equivalenti a \(8.4\,\mathrm{MJ}\).
\subsection{Reazioni esoergoniche ed endoergoniche}
\begin{itemize}
\item \textbf{Reazioni esoergoniche}: l'energia dei prodotti è inferiore a quella dei reagenti; rilasciano energia (es. degradazione di glucosio o lipidi).
\item \textbf{Reazioni endoergoniche}: l'energia dei prodotti è superiore; richiedono apporto energetico (es. sintesi di ATP, biosintesi di macromolecole).
\end{itemize}
L'\textbf{ATP} (adenosina trifosfato) funge da ``moneta energetica'': la sua idrolisi (\(\mathrm{ATP \rightarrow ADP + P_i}\)) rilascia circa \(30\,\mathrm{kJ/mol}\), utilizzati per lavoro cellulare (contrazione muscolare, trasporto attivo, sintesi).
\subsection{Classificazione delle reazioni biochimiche}
\begin{enumerate}
\item \textbf{Sintesi (anabolismo)}: due molecole si combinano per formarne una più complessa, spesso con rilascio di acqua (\textit{reazione di condensazione}). Esempio: formazione del legame peptidico tra amminoacidi. Reazioni endoergoniche.
\item \textbf{Degradazione (catabolismo)}: una molecola complessa viene scissa in unità più semplici, spesso con consumo di acqua (\textit{reazione di idrolisi}). Esempio: idrolisi dell'ATP. Reazioni esoergoniche.
\item \textbf{Fosforilazione/defosforilazione}: trasferimento di un gruppo fosfato (\(\mathrm{PO_4^{2-}}\)) da o verso una molecola. Meccanismo chiave nella regolazione enzimatica e nell'attivazione metabolica (es. fosforilazione del glucosio nella glicolisi).
\item \textbf{Ossidoriduzioni (redox)}: trasferimento di elettroni da una molecola (che si \textit{ossida}) a un'altra (che si \textit{riduce}). Spesso coinvolgono il trasferimento di atomi di idrogeno (\textit{deidrogenazioni}). Sono fondamentali nella respirazione cellulare e nella fotosintesi.
\end{enumerate}
\subsection{Reversibilità delle reazioni}
Le reazioni possono essere:
\begin{itemize}
\item \textbf{Irreversibili}: procedono in una sola direzione;
\item \textbf{Reversibili}: indicate con doppia freccia (\(\rightleftharpoons\)); le due direzioni avvengono simultaneamente, ma una può essere favorita in base a concentrazione dei reagenti, pH, temperatura, ecc.
\end{itemize}

\section{L'acqua: proprietà e ruolo biologico}
L'acqua (\(\mathrm{H_2O}\)) è il composto inorganico più abbondante negli organismi viventi, costituendo oltre il 60\% della massa corporea (con variazioni tissutali: sangue 90\%, muscolo 75\%, osso 25\%, tessuto adiposo 5\%). Circa due terzi dell'acqua corporea sono intracellulari.
\subsection{Funzioni biologiche dell'acqua}
\begin{itemize}
\item Trasporto di nutrienti e ossigeno (tramite il sangue);
\item Eliminazione di prodotti di scarto (urea, sali) tramite le urine;
\item Termoregolazione mediante sudorazione;
\item Solvente per reazioni biochimiche;
\item Reagente (idrolisi) o prodotto (condensazione) in reazioni metaboliche.
\end{itemize}
Una perdita di acqua pari al 2\% della massa corporea compromette già le performance fisiche e cognitive.
\subsection{Polarità e capacità solvente}
La geometria angolare e la differenza di elettronegatività tra ossigeno e idrogeno conferiscono all'acqua una forte polarità: l'ossigeno presenta parziale carica negativa (\(\delta^-\)), gli idrogeni parziale carica positiva (\(\delta^+\)).
Ciò permette all'acqua di:
\begin{itemize}
\item Formare fino a 4 legami idrogeno per molecola, creando una rete dinamica;
\item Sciogliere sostanze polari o ioniche (\textbf{idrofiliche}), come zuccheri e sali;
\item Escludere molecole apolari (\textbf{idrofobiche}), come lipidi, che tendono ad aggregarsi per minimizzare il contatto con l'acqua (\textit{effetto idrofobico}).
\end{itemize}
\begin{definizione}[Molecole anfipatiche]
Molecole che possiedono sia una regione polare/idrofilica sia una regione apolare/idrofobica sono dette \textbf{anfipatiche} (es. fosfolipidi, acidi biliari). In ambiente acquoso, si organizzano spontaneamente in strutture come micelle o doppi strati lipidici, fondamentali per le membrane biologiche.
\end{definizione}
\subsection{Stati fisici dell'acqua}
I cambiamenti di stato riflettono variazioni nella rete di legami idrogeno:
\begin{itemize}
\item \textbf{Ghiaccio}: molecole fisse, ciascuna forma 4 legami idrogeno stabili;
\item \textbf{Acqua liquida}: legami idrogeno in continuo formarsi e rompersi;
\item \textbf{Vapore}: assenza di legami idrogeno significativi.
\end{itemize}

\section{Acidi, basi e regolazione del pH}
\subsection{Definizioni di acido e base}
\begin{itemize}
\item \textbf{Acido}: composto che, in soluzione acquosa, dona uno ione idrogeno (\(\mathrm{H^+}\), protone). Esempio: acido lattico \(\rightleftharpoons\) lattato + \(\mathrm{H^+}\).
\item \textbf{Base}: composto che accetta uno ione \(\mathrm{H^+}\). Esempio: \(\mathrm{NH_3 + H^+ \rightleftharpoons NH_4^+}\).
\end{itemize}
Acidi e basi reagiscono tra loro formando sali e acqua (reazione di neutralizzazione): 
\[
\mathrm{HCl + NaOH \rightarrow NaCl + H_2O}
\]
\subsection{Autoionizzazione dell'acqua}
L'acqua può comportarsi sia da acido che da base, autoionizzandosi:
\[
\mathrm{H_2O + H_2O \rightleftharpoons H_3O^+ + OH^-}
\]
In forma semplificata:
\[
\mathrm{H_2O \rightleftharpoons H^+ + OH^-}
\]
\subsection{Scala del pH}
L'acidità o basicità di una soluzione è quantificata dal \textbf{pH}, definito come:
\[
\mathrm{pH} = -\log_{10}[\mathrm{H^+}]
\]
dove \([\mathrm{H^+}]\) è la concentrazione molare di ioni idrogeno.
\begin{itemize}
\item \(\mathrm{pH} = 7\): soluzione neutra (\([\mathrm{H^+}] = [\mathrm{OH^-}]\), acqua pura);
\item \(\mathrm{pH} < 7\): soluzione acida;
\item \(\mathrm{pH} > 7\): soluzione basica/alcalina.
\end{itemize}
La scala è logaritmica: una soluzione a pH 6 ha una concentrazione di \(\mathrm{H^+}\) dieci volte maggiore di una a pH 7.
\subsection{Importanza fisiologica del pH}
Il pH ematico è strettamente regolato tra 7.35 e 7.45. Valori esterni a 7.0--7.8 possono essere letali. Durante esercizio intenso, la produzione di acido lattico può abbassare il pH muscolare fino a ~6.5, compromettendo la contrattilità.
\subsection{Sistemi tampone}
Per contrastare variazioni di pH, l'organismo utilizza \textbf{sistemi tampone}, costituiti da una coppia acido debole/base coniugata:
\[
\mathrm{HA \rightleftharpoons H^+ + A^-}
\]
\begin{itemize}
\item In presenza di eccesso di \(\mathrm{H^+}\): la base coniugata (\(\mathrm{A^-}\)) li lega, riformando l'acido debole (\(\mathrm{HA}\));
\item In presenza di eccesso di \(\mathrm{OH^-}\): l'acido debole (\(\mathrm{HA}\)) cede \(\mathrm{H^+}\), neutralizzando gli ioni idrossido.
\end{itemize}
\textbf{Principali sistemi tampone fisiologici}:
\begin{itemize}
\item \textbf{Acido carbonico/bicarbonato} (\(\mathrm{H_2CO_3/HCO_3^-}\)): principale tampone extracellulare (sangue);
\item \textbf{Fosfati} (\(\mathrm{H_2PO_4^-/HPO_4^{2-}}\)): tampone intracellulare e urinario;
\item \textbf{Carnosina} e \textbf{fosfocreatina}: tamponi intracellulari nel muscolo scheletrico;
\item \textbf{Emoglobina}: tampona gli \(\mathrm{H^+}\) nei globuli rossi durante il trasporto di \(\mathrm{CO_2}\).
\end{itemize}
Questi sistemi cooperano per mantenere l'omeostasi acido-base anche in condizioni metaboliche estreme.
% =========================================================
% Fine Capitolo 1
% =========================================================

% =========================================================
% CAPITOLO 2: La cellula eucariotica
% =========================================================
\chapter{La cellula eucariotica}
\section{Introduzione alla cellula}
La \textbf{cellula} è la più piccola unità funzionale in grado di mantenere le caratteristiche del tessuto di appartenenza. Gli organismi viventi possono essere classificati in:
\begin{itemize}
\item \textbf{Unicellulari}: costituiti da una singola cellula (es. batteri, lieviti);
\item \textbf{Multicellulari}: composti da numerose cellule specializzate (es. organismi animali e vegetali).
\end{itemize}
Il corpo umano è costituito da circa \(3 \times 10^{13}\) cellule, appartenenti a circa 200 tipi diversi, che si organizzano secondo una gerarchia strutturale:
\[
\text{Cellule} \rightarrow \text{Tessuti} \rightarrow \text{Organi} \rightarrow \text{Sistemi} \rightarrow \text{Organismo}
\]
Nonostante la diversità morfologica e funzionale, le cellule condividono componenti strutturali di base, adattate alle specifiche esigenze fisiologiche.

\section{Procarioti ed eucarioti}
\begin{definizione}[Cellula eucariotica]
Una cellula \textbf{eucariotica} è caratterizzata dalla presenza di un nucleo delimitato da membrana, contenente il materiale genetico organizzato in cromosomi lineari, e di organelli membranosi specializzati.
\end{definizione}
La distinzione fondamentale tra i due domini cellulari è sintetizzata nella Tabella~\ref{tab:procarioti-eucarioti}.
\begin{table}[htbp]
\centering
\caption{Confronto tra cellule procariotiche ed eucariotiche}
\label{tab:procarioti-eucarioti}
\renewcommand{\arraystretch}{1.2}
\begin{tabular}{@{} p{0.45\textwidth} p{0.45\textwidth} @{}}
\toprule
\textbf{Procarioti} & \textbf{Eucarioti} \\
\midrule
DNA circolare libero nel citoplasma & DNA lineare racchiuso nel nucleo \\
Assenza di organelli membranosi & Presenza di organelli membranosi \\
Dimensioni tipiche: 1--10 µm & Dimensioni tipiche: 10--100 µm \\
Organismi: batteri & Organismi: animali, piante, funghi \\
\bottomrule
\end{tabular}
\end{table}

\section{Composizione chimica della cellula}
La maggior parte delle cellule presenta un contenuto idrico superiore al 70\%. Il \textbf{peso secco} (privato dell'acqua) è costituito principalmente da macromolecole biologiche, come riportato in Tabella~\ref{tab:composizione-cellula}.
\begin{table}[htbp]
\centering
\caption{Composizione percentuale del peso secco cellulare}
\label{tab:composizione-cellula}
\renewcommand{\arraystretch}{1.2}
\begin{tabular}{@{} l c @{}}
\toprule
\textbf{Componente} & \textbf{Peso secco (\%)} \\
\midrule
Proteine & 75 \\
Lipidi & 12,5 \\
Glucidi & 5 \\
Acidi nucleici (DNA, RNA) & 5 \\
Composti inorganici / sali minerali & 2,5 \\
\bottomrule
\end{tabular}
\end{table}
Le biomolecole si organizzano in \textbf{macromolecole} e successivamente in \textbf{complessi sopramolecolari}, che costituiscono le strutture funzionali della cellula.

\section{La membrana plasmatica}
La \textbf{membrana plasmatica} è una struttura dinamica che svolge funzioni essenziali:
\begin{itemize}
\item \textbf{Separazione}: delimita il citoplasma dall'ambiente extracellulare;
\item \textbf{Filtro selettivo}: regola il passaggio di sostanze in entrata e in uscita;
\item \textbf{Comunicazione}: ospita recettori per segnali ormonali e molecole di adesione;
\item \textbf{Interazione}: media i contatti fisici con altre cellule e con la matrice extracellulare.
\end{itemize}
\subsection{Struttura e composizione}
La membrana è costituita principalmente da un \textbf{doppio strato fosfolipidico}, in cui le code idrofobiche sono orientate verso l'interno e le teste polari verso l'ambiente acquoso. A questo si associano:
\begin{itemize}
\item \textbf{Colesterolo}: modula la fluidità di membrana in risposta alle variazioni termiche;
\item \textbf{Glicolipidi}: coinvolti nel riconoscimento cellulare;
\item \textbf{Proteine di membrana}, distinte in:
\begin{itemize}
    \item \textit{Integrali} (o intrinseche): attraversano completamente il doppio strato;
    \item \textit{Periferiche} (o estrinseche): associate a una delle due facce della membrana.
\end{itemize}
\end{itemize}

\section{Trasporto attraverso le membrane}
Il passaggio di molecole attraverso la membrana avviene tramite meccanismi distinti, classificabili in base alla direzione del gradiente e al consumo energetico.
\subsection{Trasporto passivo}
Nel \textbf{trasporto passivo}, le molecole si muovono \textit{secondo gradiente} di concentrazione (da zone a concentrazione maggiore a zone a concentrazione minore), senza dispendio energetico.
\begin{itemize}
\item \textbf{Diffusione semplice}: molecole apolari o gas (\(\mathrm{O_2}\), \(\mathrm{CO_2}\), \(\mathrm{N_2}\)) attraversano direttamente il doppio strato lipidico;
\item \textbf{Diffusione facilitata}: molecole polari, cariche o di grandi dimensioni (es. glucosio) richiedono proteine di trasporto specifiche (es. trasportatori GLUT).
\end{itemize}
\subsection{Trasporto attivo}
Nel \textbf{trasporto attivo}, le molecole vengono spostate \textit{contro gradiente} di concentrazione, richiedendo energia (generalmente sotto forma di ATP) e proteine carrier specializzate.
\begin{definizione}[Tipologie di trasporto mediato]
\begin{itemize}
\item \textbf{Uniporto}: trasporto di una singola specie molecolare;
\item \textbf{Simporto}: trasporto contemporaneo di due specie nella stessa direzione;
\item \textbf{Antiporto}: trasporto contemporaneo di due specie in direzioni opposte.
\end{itemize}
\end{definizione}
\subsection{Proteine recettoriali}
Alcune proteine di membrana fungono da \textbf{recettori}, capaci di legare specifici ligandi extracellulari (es. ormoni) e trasdurre il segnale all'interno della cellula, attivando cascate di risposta metabolica.
\begin{esempio}[Segnalazione insulinica]
L'insulina, secreta dal pancreas, si lega al proprio recettore di membrana sulle cellule muscolari, innescando una cascata di fosforilazioni che promuove l'ingresso di glucosio e la sintesi di glicogeno.
\end{esempio}

\section{Organelli cellulari}
Oltre alla membrana plasmatica, le cellule eucariotiche presentano un sistema di membrane interne che delimitano \textbf{organelli} specializzati, ciascuno con composizione lipidica e proteica specifica.
\subsection{Il nucleo}
\begin{itemize}
\item Contiene il materiale genetico (DNA lineare) organizzato in \textbf{cromosomi}, complessati con proteine basiche dette \textbf{istoni};
\item È delimitato da un \textbf{involucro nucleare} a doppia membrana, perforato da \textbf{pori nucleari} che regolano il traffico nucleo-citoplasmatico;
\item Le cellule umane somatiche possiedono 46 cromosomi (23 coppie);
\item Eccezioni: i globuli rossi maturi sono anucleati; le fibre muscolari scheletriche sono multinucleate.
\end{itemize}
\subsection{Il citoplasma e il citoscheletro}
Il \textbf{citoplasma} comprende tutto il contenuto cellulare escluso il nucleo ed è costituito da:
\begin{itemize}
\item \textbf{Citosol}: fase acquosa contenente ioni, metaboliti e macromolecole;
\item \textbf{Organelli}: strutture membranose specializzate;
\item \textbf{Citoscheletro}: rete dinamica di filamenti proteici che conferisce forma, stabilità e capacità di movimento alla cellula.
\end{itemize}
\begin{table}[htbp]
\centering
\caption{Componenti del citoscheletro}
\label{tab:citoscheletro}
\renewcommand{\arraystretch}{1.2}
\begin{tabular}{@{} p{0.25\textwidth} p{0.35\textwidth} p{0.35\textwidth} @{}}
\toprule
\textbf{Tipo} & \textbf{Composizione} & \textbf{Funzioni principali} \\
\midrule
Microfilamenti & Actina & Supporto meccanico, contrazione muscolare, motilità cellulare \\
Filamenti intermedi & Proteine fibrose (es. cheratine, vimentina) & Ancoraggio di organelli, resistenza meccanica \\
Microtubuli & Tubulina & Determinazione della forma, trasporto intracellulare, fuso mitotico \\
\bottomrule
\end{tabular}
\end{table}
\subsection{Reticolo endoplasmatico}
Il \textbf{reticolo endoplasmatico} (RE) è una rete di membrane continue con l'involucro nucleare, distinta in due domini funzionali:
\begin{itemize}
\item \textbf{RE rugoso}: presenta ribosomi associati alla superficie citosolica; è sede di sintesi e modificazione post-traduzionale delle proteine destinate a secrezione o a organelli membranosi;
\item \textbf{RE liscio}: privo di ribosomi; coinvolto nella biosintesi lipidica, nel metabolismo dei carboidrati e nella detossificazione. Nel muscolo scheletrico, il RE liscio specializzato (\textbf{reticolo sarcoplasmatico}) regola il rilascio di ioni \(\mathrm{Ca^{2+}}\) necessari per la contrazione.
\end{itemize}
\subsection{Apparato di Golgi}
L'\textbf{apparato di Golgi} è costituito da una serie di cisterne membranose impilate, organizzate in due polarità:
\begin{itemize}
\item \textbf{Lato \textit{cis}}: riceve vescicole di trasporto dal RE rugoso;
\item \textbf{Lato \textit{trans}}: rilascia vescicole mature verso la membrana plasmatica, i lisosomi o altri compartimenti.
\end{itemize}
Funzioni principali: smistamento e indirizzamento delle proteine neosintetizzate, modificazioni post-traduzionali (glicosilazione, solfatazione), sintesi di carboidrati complessi e glicolipidi.
\subsection{Mitocondri}
I \textbf{mitocondri} sono gli organelli deputati alla produzione di energia cellulare sotto forma di ATP, tramite la fosforilazione ossidativa. Struttura a doppia membrana (esterna permeabile, interna ricca di creste e cardiolipina), contenente DNA mitocondriale circolare e ribosomi propri. Ereditati per via materna.
\subsection{Lisosomi}
I \textbf{lisosomi} sono vescicole membranose contenenti \textbf{enzimi idrolitici} attivi a pH acido (\(\mathrm{pH} \approx 4.5\text{--}5.0\)), deputati alla degradazione di macromolecole endocitate o di componenti cellulari danneggiati (autofagia). La membrana è protetta da un rivestimento glicoproteico interno.
% =========================================================
% Fine Capitolo 2
% =========================================================

% =========================================================
% CAPITOLO 3: Amminoacidi
% =========================================================
\chapter{Amminoacidi}
\section{Struttura generale e proprietà chimiche}
Gli \textbf{amminoacidi} sono molecole organiche caratterizzate dalla presenza di almeno due gruppi funzionali specifici: un \textbf{gruppo carbossilico} (\(-\mathrm{COOH}\)) e un \textbf{gruppo amminico} (\(-\mathrm{NH_2}\)). Negli amminoacidi costituenti le proteine, entrambi i gruppi sono legati allo stesso atomo di carbonio, denominato \(\alpha\)-carbonio (\(\alpha\)-C). Per questa ragione vengono definiti \(\alpha\)-amminoacidi.
Oltre al gruppo carbossilico e amminico, il \(\alpha\)-C è sempre legato a un atomo di idrogeno (\(\mathrm{H}\)) e a un gruppo variabile indicato con \(\mathrm{R}\), detto \textbf{catena laterale}. La composizione, la dimensione e la carica del gruppo \(\mathrm{R}\) differiscono per ciascun amminoacido e ne determinano le proprietà fisico-chimiche distintive, come la solubilità in acqua o la reattività.

\section{Chiralità e stereoisomeria}
Con l'eccezione della \textbf{glicina}, il cui gruppo \(\mathrm{R}\) è costituito da un singolo atomo di idrogeno, il \(\alpha\)-C di tutti gli amminoacidi è legato a quattro gruppi differenti. Di conseguenza, il \(\alpha\)-C rappresenta un \textbf{centro chirale} e i quattro sostituenti possono disporsi nello spazio in due configurazioni differenti, generando due stereoisomeri noti come \textbf{enantiomeri}. Gli enantiomeri sono immagini speculari non sovrapponibili l'una dell'altra.
Le due configurazioni sono designate con la notazione \textbf{D} (destrogiro) e \textbf{L} (levogiro), convenzionalmente definite in riferimento alla configurazione della gliceraldeide. Nei sistemi biologici, \textbf{tutti gli amminoacidi incorporati nelle proteine appartengono alla serie L}. Gli amminoacidi della serie D sono rari e si ritrovano esclusivamente in alcuni peptidi batterici o in contesti metabolici specifici. Questa omogeneità stereochimica è dovuta alla specificità degli enzimi biosintetici, che catalizzano reazioni stereospecifiche producendo esclusivamente L-amminoacidi.

\section{Ionizzazione e forma zwitterionica}
I gruppi carbossilico e amminico, nonché alcune catene laterali, sono \textbf{ionizzabili} e il loro stato di protonazione dipende dal pH dell'ambiente:
\begin{itemize}
\item \textbf{A pH acido}: entrambi i gruppi sono protonati (\(-\mathrm{COOH}\) e \(-\mathrm{NH_3^+}\)); la molecola presenta una carica netta positiva.
\item \textbf{A pH fisiologico} (\(\approx 7.4\)): il gruppo carbossilico perde un protone (\(-\mathrm{COO^-}\)), mentre il gruppo amminico rimane protonato (\(-\mathrm{NH_3^+}\)). La molecola possiede simultaneamente una carica positiva e una negativa, risultando elettricamente neutra. Questa forma è detta \textbf{zwitterione}. Fanno eccezione gli amminoacidi con catene laterali ionizzabili (es. glutammato, lisina).
\item \textbf{A pH basico}: anche il gruppo amminico perde un protone (\(-\mathrm{NH_2}\)); la molecola assume una carica netta negativa.
\end{itemize}

\section{Classificazione in base alla catena laterale}
Le proteine sono costituite da 20 amminoacidi standard (più la selenocisteina, considerata il 21°), identificabili tramite abbreviazioni di tre lettere o simboli di una lettera. Sulla base delle proprietà chimiche del gruppo \(\mathrm{R}\), vengono classificati in cinque categorie principali.

\subsection{Amminoacidi non polari (idrofobici)}
Presentano catene laterali apolari, insolubili in acqua, e tendono a localizzarsi all'interno delle proteine globulari per minimizzare il contatto con il solvente acquoso, stabilizzando la struttura tramite \textit{interazioni idrofobiche}.
\begin{itemize}
\item \textbf{Glicina (Gly, G)}: il più semplice (\(\mathrm{R} = \mathrm{H}\)). Privo di chiralità, meno idrofobico del gruppo, non partecipa significativamente alle interazioni idrofobiche.
\item \textbf{Alanina (Ala, A), Valina (Val, V), Leucina (Leu, L), Isoleucina (Ile, I)}: catene completamente idrocarburiche. Val, Leu e Ile sono noti come \textbf{BCAA} (amminoacidi a catena ramificata). Sono essenziali e possono essere ossidati dal muscolo scheletrico a scopo energetico. Leucina e isoleucina sono isomeri strutturali.
\item \textbf{Metionina (Met, M)}: contiene zolfo legato tramite un ponte tioetere (\(-\mathrm{S-CH_3}\)), non polare.
\item \textbf{Prolina (Pro, P)}: presenta un gruppo amminico secondario (imminico) che forma un anello ciclico rigido con il \(\alpha\)-C. Riduce la flessibilità conformazionale delle catene proteiche in cui è inserita.
\end{itemize}

\subsection{Amminoacidi aromatici}
Contengono anelli aromatici nelle catene laterali, conferendo caratteristiche generalmente non polari.
\begin{itemize}
\item \textbf{Fenilalanina (Phe, F)}: completamente non polare.
\item \textbf{Tirosina (Tyr, Y)}: possiede un gruppo ossidrilico (\(-\mathrm{OH}\)) fenolico che può formare legami idrogeno e subire \textbf{fosforilazione} (modificazione post-traduzionale regolatoria).
\item \textbf{Triptofano (Trp, W)}: contiene un gruppo amminico secondario nell'anello indolico, capace di formare legami idrogeno.
\end{itemize}

\subsection{Amminoacidi polari non carichi}
Le catene laterali contengono gruppi funzionali capaci di formare legami idrogeno con l'acqua, rendendoli solubili. A pH fisiologico non presentano carica netta.
\begin{itemize}
\item \textbf{Serina (Ser, S) e Treonina (Thr, T)}: polarità dovuta al gruppo ossidrilico (\(-\mathrm{OH}\)). La treonina differisce per la presenza di un gruppo metile. I gruppi \(-\mathrm{OH}\) di Ser, Thr e Tyr sono siti privilegiati per la \textbf{fosforilazione}, meccanismo chiave nella regolazione enzimatica e nella trasduzione del segnale.
\item \textbf{Cisteina (Cys, C)}: contiene un gruppo sulfidrilico (\(-\mathrm{SH}\)) di modesta polarità. Due cisteine possono ossidarsi formando un \textbf{ponte disolfuro} (\(-\mathrm{S-S-}\)), legame covalente fondamentale per la stabilizzazione della struttura terziaria e quaternaria delle proteine.
\item \textbf{Asparagina (Asn, N) e Glutammina (Gln, Q)}: contengono gruppi ammidici (\(-\mathrm{CONH_2}\)) che partecipano a reti di legami idrogeno. Sono le ammidi corrispondenti di acido aspartico e acido glutammico.
\end{itemize}

\subsection{Amminoacidi acidi (carichi negativamente)}
Presentano un gruppo carbossilico aggiuntivo nella catena laterale, ionizzabile a pH fisiologico. Sono pertanto carichi negativamente e fortemente idrofilici.
\begin{itemize}
\item \textbf{Acido aspartico / Aspartato (Asp, D)}
\item \textbf{Acido glutammico / Glutammato (Glu, E)}
\end{itemize}
Oltre al ruolo strutturale nelle proteine, questi amminoacidi svolgono funzioni metaboliche cruciali e partecipano alla trasduzione del segnale nervoso.

\subsection{Amminoacidi basici (carichi positivamente)}
Contengono gruppi basici che rimangono protonati a pH 7, conferendo carica positiva e massima idrofilicità.
\begin{itemize}
\item \textbf{Lisina (Lys, K)}: catena lineare con gruppo amminico terminale.
\item \textbf{Arginina (Arg, R)}: contiene il gruppo \textbf{guanidinico}, fortemente basico.
\item \textbf{Istidina (His, H)}: possiede un anello imidazolico. Il suo pKa (\(\approx 6.0\)) lo rende parzialmente protonato a pH fisiologico, fungendo da tampone intracellulare e partecipando attivamente ai siti catalitici enzimatici.
\end{itemize}

\section{Classificazione nutrizionale e destino metabolico}
\subsection{Essenziali e non essenziali}
In base alla capacità biosintetica dell'organismo:
\begin{itemize}
\item \textbf{Non essenziali}: sintetizzati \textit{de novo} dall'organismo.
\item \textbf{Essenziali}: non sintetizzabili, devono essere introdotti con la dieta. I BCAA rientrano in questa categoria e sono tra i più abbondanti nelle proteine muscolari.
\item \textbf{Condizionalmente essenziali}: in condizioni fisiologiche estreme (sviluppo, stress, patologie), la richiesta supera la capacità sintetica, rendendoli di fatto indispensabili dalla dieta.
\end{itemize}

\subsection{Destino metabolico}
In carenza di carboidrati e lipidi, gli amminoacidi possono essere catabolizzati a scopo energetico:
\begin{itemize}
\item \textbf{Glucogenici}: la degradazione produce intermedi utilizzabili per la gluconeogenesi (es. Asp, Glu, Asn, Gln, His, Pro, Arg, Gly, Ala, Ser, Cys, Met, Val).
\item \textbf{Chetogenici}: la degradazione produce acetil-CoA, precursore della chetogenesi (Leu, Lys).
\item \textbf{Entrambi}: Phe, Tyr, Trp, Ile, Thr.
\end{itemize}
Il metabolismo degli amminoacidi (tranne i BCAA) avviene prevalentemente nel fegato, che li trasforma in substrati utilizzabili da altri tessuti. I BCAA sono invece ossidati direttamente nel muscolo scheletrico.

\section{Amminoacidi speciali e non proteici}
Oltre ai 20-21 amminoacidi proteinogenici, esistono derivati e modificazioni post-traduzionali con ruoli fisiologici distinti:
\begin{itemize}
\item \textbf{Selenocisteina (Sec, U)}: 21° amminoacido proteinogenico. Analogo della cisteina con selenio al posto dello zolfo. Presente nelle \textbf{selenoproteine} (es. glutatione perossidasi) con funzioni antiossidanti.
\item \textbf{Idrossilisina e Idrossiprolina}: derivano da lisina e prolina tramite idrossilazione post-traduzionale. Abbondanti nel collagene e nell'elastina; il gruppo \(-\mathrm{OH}\) favorisce la formazione di legami idrogeno e la glicosilazione.
\item \textbf{L-Carnitina}: sintetizzata da lisina e metionina. Trasporta acidi grassi a catena lunga nella matrice mitocondriale per la \(\beta\)-ossidazione. La carenza provoca miopatie e accumulo lipidico.
\item \textbf{GABA} (\(\gamma\)-aminobutirrico): derivato del glutammato, principale neurotrasmettitore inibitorio del SNC. Modula ansia, sonno e stimola il rilascio di GH.
\item \textbf{Citrullina e Ornitina}: intermedi del ciclo dell'urea. Derivano dall'arginina; la citrullina è sottoprodotto della sintesi di ossido nitrico (NO), potente vasodilatatore. Utilizzate in ambito sportivo per migliorare il flusso ematico e ridurre la fatica.
\item \textbf{\(\beta\)-alanina}: isomero dell'alanina con gruppo amminico in posizione \(\beta\). Precursore della carnosina (dipeptide con istidina), importante sistema tampone intramuscolare contro l'acidosi da esercizio. Componente del Coenzima A.
\item \textbf{Taurina}: amminoacido solfonico (\(-\mathrm{SO_3H}\)). Abbondante nel muscolo, componente dei sali biliari e regolatore del volume cellulare. Presente negli energy drink, sebbene i benefici ergogenici siano dibattuti.
\end{itemize}
% =========================================================
% Fine Capitolo 3
% =========================================================

% =========================================================
% CAPITOLO 4: Le Proteine
% =========================================================
\chapter{Le Proteine}
\section{Introduzione e legame peptidico}
Le \textbf{proteine} sono macromolecole biologiche di dimensioni altamente variabili, comprese tra pochi kilodalton (kDa) e alcuni megadalton (MDa). Sono mediamente più estese di lipidi e carboidrati e le loro unità monomeriche sono i 20 amminoacidi proteinogenici.
Gli amminoacidi all'interno di una proteina sono uniti tramite il \textbf{legame peptidico}, un legame covalente che si forma tra il gruppo carbossilico (\(-\mathrm{COOH}\)) di un amminoacido e il gruppo amminico (\(-\mathrm{NH_2}\)) del successivo. La reazione libera una molecola di acqua e rappresenta quindi una \textbf{reazione di condensazione}. Il legame peptidico non coinvolge le catene laterali (\(\mathrm{R}\)).
Due amminoacidi uniti formano un \textbf{dipeptide}; tre formano un \textbf{tripeptide}. Per catene più lunghe si utilizzano le seguenti convenzioni:
\begin{itemize}
\item \(< 20\) amminoacidi: \textbf{oligopeptide}
\item \(20\text{--}100\) amminoacidi: \textbf{polipeptide}
\item \(> 100\) amminoacidi: \textbf{proteina}
\end{itemize}
Nella rappresentazione lineare della sequenza, si adotta la convenzione di scrivere il gruppo amminico libero a sinistra e il gruppo carbossilico libero a destra. Il primo residuo è definito \textbf{N-terminale}, l'ultimo \textbf{C-terminale}. Poiché solo le estremità mantengono i gruppi funzionali liberi, le proprietà idrofiliche/idrofobiche della molecola dipendono principalmente dalla natura delle catene laterali.

\section{Struttura delle proteine}
La conformazione proteica viene descritta a quattro livelli gerarchici.
\subsection{Struttura primaria}
Rappresenta la sequenza lineare di amminoacidi, determinata univocamente dall'informazione genetica contenuta nel DNA. Il numero e l'ordine dei residui definiscono la funzione della proteina; anche una singola sostituzione (mutazione) può comprometterne l'attività o renderla patologica. La sequenza primaria contiene le informazioni necessarie per il ripiegamento tridimensionale, sebbene la previsione computazionale della struttura finale rimanga complessa.

\subsection{Struttura secondaria}
Corrisponde ai ripiegamenti locali ricorrenti della catena polipeptidica, stabilizzati da \textbf{legami idrogeno} regolari tra i gruppi ammidici e carbonilici dello scheletro peptidico. Le due conformazioni più comuni sono:
\begin{itemize}
\item \textbf{\(\alpha\)-elica}: struttura a spirale stabilizzata da legami idrogeno intra-catena.
\item \textbf{Foglietto \(\beta\) ripiegato}: disposizione planare di segmenti polipeptidici adiacenti, collegati da legami idrogeno inter-catena o intra-catena.
\end{itemize}
Combinazioni di \(\alpha\)-eliche e foglietti \(\beta\) formano \textbf{motivi strutturali} (strutture supersecondarie).

\subsection{Struttura terziaria}
È la conformazione tridimensionale complessiva della singola catena polipeptidica, ovvero la disposizione spaziale di tutti gli atomi. A differenza della struttura secondaria, la stabilità terziaria è garantita da interazioni tra catene laterali:
\begin{itemize}
\item \textbf{Ponti disolfuro} (\(-\mathrm{S-S-}\)): legami covalenti tra residui di cisteina, molto stabili.
\item \textbf{Interazioni idrofobiche}: i residui apolari si ripiegano verso l'interno della proteina per evitare il contatto con l'acqua, mentre i residui polari si dispongono sulla superficie.
\item \textbf{Legami idrogeno} e \textbf{interazioni ioniche} tra gruppi carichi.
\end{itemize}
Una proteina funzionale assume la sua \textbf{conformazione nativa}; la perdita di tale struttura è definita \textbf{denaturazione}.
Le catene polipeptidiche lunghe (\(>200\) residui) si organizzano spesso in \textbf{domini strutturali}: unità globulari compatte, indipendenti dal punto di vista strutturale e funzionale, collegate da regioni flessibili.

\subsection{Struttura quaternaria}
Presente nelle proteine multimeriche, descrive l'associazione di due o più catene polipeptidiche (\textbf{subunità}). Le subunità possono essere identiche (\textbf{omomultimeriche}, es. omodimero) o diverse (\textbf{eteromultimeriche}, es. eterodimero). Le interazioni tra subunità sono dello stesso tipo di quelle che stabilizzano la struttura terziaria. Esempi noti sono l'\textbf{emoglobina} (tetramero: \(2\alpha + 2\beta\)) e la \textbf{miosina} (esamero: 2 catene pesanti + 4 leggere).

\section{Classificazione e proteine coniugate}
\subsection{Forma e solubilità}
\begin{itemize}
\item \textbf{Proteine fibrose}: una dimensione prevale sulle altre. Presentano spesso un unico tipo di struttura secondaria ripetuta, sono insolubili in acqua e svolgono ruoli strutturali o di supporto (es. cheratine, collageno, elastina, seta).
\item \textbf{Proteine globulari}: le tre dimensioni sono comparabili, assumendo forma compatta. Sono generalmente solubili e svolgono funzioni dinamiche (enzimi, trasportatori, ormoni, anticorpi).
\end{itemize}

\subsection{Proteine coniugate}
Molte proteine contengono \textbf{gruppi prostetici} non amminoacidici, essenziali per la funzione. Possono essere legati covalentemente (es. carboidrati nelle glicoproteine, fosfati nelle fosfoproteine), associati strettamente in tasche idrofobiche (es. gruppo \textbf{eme} nell'emoglobina) o legati tramite interazioni ioniche (metalli nelle metalloproteine).

\section{Modifiche post-traduzionali}
Dopo la sintesi, le proteine possono subire modificazioni chimiche che ne regolano attività, localizzazione o stabilità:
\begin{itemize}
\item \textbf{Fosforilazione}: aggiunta di un gruppo fosfato (da ATP) a residui di Serina, Treonina o Tirosina. Modifica conformazione e attività; è un meccanismo reversibile di regolazione enzimatica (es. attivazione della glicogeno fosforilasi in risposta al glucagone).
\item \textbf{Acetilazione}: aggiunta di un gruppo acetile (\(-\mathrm{COCH_3}\)) a residui di Lisina.
\item \textbf{Glicosilazione}: aggiunta di catene oligosaccaridiche. La N-glicosilazione avviene sull'Asparagina, la O-glicosilazione su Serina/Treonina. Cruciale per il targeting e la funzione delle proteine di membrana.
\end{itemize}

\section{Sintesi, turnover e degradazione}
\subsection{Sintesi proteica}
L'informazione fluisce dal DNA all'mRNA (\textbf{trascrizione}), che viene esportato nel citosol e tradotto dai ribosomi in catena polipeptidica (\textbf{traduzione}). Il codice genetico è letto in triplette (\textbf{codoni}); 61 codoni specificano amminoacidi (codice degenerato), 3 sono codoni di stop. La traduzione inizia generalmente dal codone AUG (Metionina). Il folding avviene spontaneamente o con l'assistenza di \textbf{chaperoni}. Il processo è energeticamente costoso e regolato finemente, ad esempio dalla proteina \textbf{mTOR}, attivata in condizioni di abbondanza nutrizionale e promotrice dell'ipertrofia muscolare.

\subsection{Turnover e degradazione}
Le proteine sono soggette a continuo \textbf{turnover} (sintesi + degradazione), che consuma fino al 20\% del metabolismo basale. Ogni proteina ha un'\textbf{emivita} variabile (minuti per enzimi epatici, settimane per proteine muscolari). La degradazione avviene principalmente tramite:
\begin{itemize}
\item \textbf{Lisosomi}: degradano proteine extracellulari (endocitosi) o intracellulari/organelli (autofagia) in condizioni di digiuno o stress. Processo generalmente non selettivo, ottimale a pH acido.
\item \textbf{Proteasoma}: complesso citosolico che degrada proteine intracellulari marcate con \textbf{ubiquitina} (ubiquitinazione). Processo ATP-dipendente e selettivo.
\item \textbf{Digestione proteica}: scissione delle proteine alimentari in amminoacidi assorbibili tramite proteasi gastrointestinali.
\end{itemize}

\section{Funzioni biologiche}
Le proteine svolgono una vasta gamma di ruoli fisiologici:
\begin{enumerate}
\item \textbf{Strutturale e contrattile}: Actina e miosina (contrazione muscolare), collageno e cheratina (supporto tissutale), distrofina (ancoraggio miofibrille al sarcolemma).
\item \textbf{Trasporto e deposito}: Proteine di membrana (canali, trasportatori GLUT), emoglobina (trasporto \(\mathrm{O_2}\) tramite gruppo eme).
\item \textbf{Ormonale}: Insulina, glucagone, GH (regolazione metabolica e crescita).
\item \textbf{Segnalazione}: Recettori di membrana e cascate di chinasi intracellulari che trasducono stimoli esterni in risposte cellulari.
\item \textbf{Immunitaria}: Immunoglobuline (anticorpi) che riconoscono e neutralizzano antigeni.
\item \textbf{Regolazione}: Fattori di trascrizione (es. PGC-1\(\alpha\)) che modulano l'espressione genica in risposta a stimoli come l'esercizio.
\item \textbf{Catalitica}: Enzimi (es. esochinasi) che accelerano le reazioni metaboliche abbassando l'energia di attivazione.
\end{enumerate}

\section{Stabilità e denaturazione}
La funzione proteica è strettamente legata alla conformazione nativa. Fattori come variazioni di pH, aumento di temperatura o stress ossidativo possono rompere i legami deboli che stabilizzano la struttura, portando a \textbf{denaturazione} (spesso irreversibile, come nella cottura dell'uovo). Una proteina denaturata perde la sua funzione e può risultare tossica, venendo quindi degradata.
Durante l'esercizio intenso, l'ambiente intracellulare diventa acido (\(\mathrm{pH} \approx 6.5\)) e ipertermico (\(\approx 40^\circ\mathrm{C}\)). Per contrastare questi effetti, le cellule esprimono \textbf{Heat Shock Proteins (HSPs)}: chaperoni molecolari che riconoscono le proteine misfolded, ne facilitano il refolding e promuovono una risposta adattativa che aumenta la resilienza cellulare a stress futuri.
% =========================================================
% Fine Capitolo 4
% =========================================================

% =========================================================
% CAPITOLO 5: Trasporto e utilizzo dell'ossigeno
% =========================================================
\chapter{Trasporto e utilizzo dell'ossigeno}
\section{Introduzione: emoglobina e mioglobina}
L'\textbf{ossigeno} (\(\mathrm{O_2}\)) è l'elemento più abbondante nel nostro organismo, principalmente incorporato nell'acqua. Costituisce inoltre il 21\% dell'atmosfera terrestre. L'\(\mathrm{O_2}\) inspirato nei polmoni deve essere trasportato ai tessuti per sostenere il metabolismo ossidativo mitocondriale.
Nel corpo umano, due proteine specializzate legano reversibilmente l'ossigeno:
\begin{itemize}
\item \textbf{Emoglobina (Hb)}: proteina tetramerica presente negli eritrociti (globuli rossi), deputata al trasporto di \(\mathrm{O_2}\) dai polmoni ai tessuti e di \(\mathrm{CO_2}\) dai tessuti ai polmoni.
\item \textbf{Mioglobina (Mb)}: proteina monomerica localizzata nelle cellule muscolari, funge da riserva di \(\mathrm{O_2}\) rilasciandolo in condizioni di ipossia per mantenere attivo il metabolismo aerobico.
\end{itemize}

\section{Struttura delle emoproteine}
Sia l'emoglobina che la mioglobina sono \textbf{emoproteine}, ovvero contengono un gruppo prostetico \textbf{eme} legato non covalentemente alla catena polipeptidica. L'eme contiene un atomo di ferro nello stato di ossidazione ferroso (\(\mathrm{Fe^{2+}}\)), responsabile del legame reversibile con l'\(\mathrm{O_2}\).
\subsection{Il gruppo eme}
L'eme è una protoporfirina IX, struttura planare derivante dalla porfirina, caratterizzata da quattro anelli pirrolici con atomi di azoto all'apice. Al centro della protoporfirina è coordinato un singolo atomo di \(\mathrm{Fe^{2+}}\).
I legami di coordinazione del ferro sono sei:
\begin{itemize}
\item Quattro legami giacciono sul piano dell'anello porfirinico e coinvolgono gli atomi di azoto dei pirroli;
\item Il quinto legame è occupato dall'azoto dell'anello imidazolico di un residuo di \textbf{istidina prossimale};
\item Il sesto legame, perpendicolare al piano, è disponibile per il legame con l'\(\mathrm{O_2}\).
\end{itemize}
Un secondo residuo di \textbf{istidina distale} stabilizza il legame \(\mathrm{Fe-O_2}\) tramite un legame idrogeno, favorendo il posizionamento planare del ferro rispetto all'eme. Per prevenire l'ossidazione del ferro a \(\mathrm{Fe^{3+}}\) (che impedirebbe il legame con \(\mathrm{O_2}\)), il gruppo eme è alloggiato in una tasca idrofobica della proteina.

\subsection{Differenze strutturali tra Hb e Mb}
\begin{table}[htbp]
\centering
\caption{Confronto tra emoglobina e mioglobina}
\label{tab:hb-mb}
\renewcommand{\arraystretch}{1.2}
\begin{tabular}{@{} p{0.45\textwidth} p{0.45\textwidth} @{}}
\toprule
\textbf{Emoglobina} & \textbf{Mioglobina} \\
\midrule
Tetramero (\(2\alpha + 2\beta\) nell'adulto) & Monomero (singola catena) \\
4 gruppi eme, lega fino a 4 \(\mathrm{O_2}\) & 1 gruppo eme, lega 1 \(\mathrm{O_2}\) \\
Curva di saturazione sigmoide (cooperatività) & Curva di saturazione iperbolica \\
Bassa affinità a bassa \(p\mathrm{O_2}\) & Alta affinità anche a bassa \(p\mathrm{O_2}\) \\
Trasporto sistemico di \(\mathrm{O_2}\) & Riserva intramuscolare di \(\mathrm{O_2}\) \\
\bottomrule
\end{tabular}
\end{table}

\section{Legame dell'ossigeno e cooperatività}
La mioglobina lega l'\(\mathrm{O_2}\) con alta affinità anche a pressioni parziali (\(p\mathrm{O_2}\)) relativamente basse, mostrando una curva di saturazione \textbf{iperbolica}. Ciò la rende ideale per immagazzinare \(\mathrm{O_2}\) nel muscolo.
L'emoglobina, invece, presenta una curva di saturazione \textbf{sigmoide}, riflettente il legame cooperativo tra le sue quattro subunità. Questo comportamento è spiegato dal modello degli stati conformazionali:
\begin{itemize}
\item \textbf{Stato T (teso)}: conformazione a bassa affinità per \(\mathrm{O_2}\), predominante in condizioni di bassa \(p\mathrm{O_2}\) (tessuti);
\item \textbf{Stato R (rilassato)}: conformazione ad alta affinità, favorita dal legame dell'\(\mathrm{O_2}\) e predominante nei polmoni.
\end{itemize}
Il legame dell'\(\mathrm{O_2}\) a una subunità induce un cambiamento conformazionale che facilita il legame alle subunità adiacenti (\textbf{cooperatività omotropica}). Nei polmoni (alta \(p\mathrm{O_2}\)), l'Hb si satura rapidamente; nei tessuti (bassa \(p\mathrm{O_2}\)), il rilascio è favorito dalla transizione allo stato T.

\section{Regolazione allosterica dell'affinità per l'ossigeno}
L'affinità dell'emoglobina per l'\(\mathrm{O_2}\) è modulata da fattori allosterici eterotropici che agiscono su siti distinti dal sito di legame dell'\(\mathrm{O_2}\):
\subsection{Effetto Bohr (pH)}
L'aumento della concentrazione di ioni \(\mathrm{H^+}\) (diminuzione del pH) riduce l'affinità dell'Hb per l'\(\mathrm{O_2}\), favorendone il rilascio nei tessuti metabolicamente attivi. Durante l'esercizio intenso, la produzione di acido lattico e di \(\mathrm{CO_2}\) abbassa il pH tissutale:
\[
\mathrm{CO_2 + H_2O \rightleftharpoons H_2CO_3 \rightleftharpoons HCO_3^- + H^+}
\]
Gli ioni \(\mathrm{H^+}\) protonano residui di istidina, stabilizzando lo stato T e promuovendo il rilascio di \(\mathrm{O_2}\). Nei polmoni, l'eliminazione di \(\mathrm{CO_2}\) e l'aumento del pH favoriscono il ricarico di \(\mathrm{O_2}\).

\subsection{Anidride carbonica (\(\mathrm{CO_2}\))}
Oltre all'effetto indiretto tramite il pH, la \(\mathrm{CO_2}\) lega direttamente l'N-terminale delle catene dell'Hb formando \textbf{carbammati}, stabilizzando lo stato T e riducendo l'affinità per l'\(\mathrm{O_2}\). Circa il 20\% della \(\mathrm{CO_2}\) prodotta dai tessuti è trasportata dall'Hb sotto questa forma.
\begin{nota}[Monossido di carbonio]
Il monossido di carbonio (\(\mathrm{CO}\)) lega il ferro dell'eme con affinità ~200 volte superiore all'\(\mathrm{O_2}\), formando carbossiemoglobina stabile e impedendo il trasporto di ossigeno (ipossia anemica).
\end{nota}

\subsection{Temperatura}
L'aumento della temperatura riduce l'affinità dell'Hb per l'\(\mathrm{O_2}\). Durante l'esercizio intenso, la temperatura muscolare può raggiungere i 40°C, favorendo il rilascio di \(\mathrm{O_2}\) a livello capillare per sostenere il metabolismo ossidativo.

\subsection{2,3-Bisfosfoglicerato (2,3-BPG)}
Il \textbf{2,3-BPG} è un metabolita glicolitico prodotto dagli eritrociti che lega esclusivamente la deossiemoglobina in una cavità centrale tra le quattro subunità, stabilizzando lo stato T e riducendo l'affinità per l'\(\mathrm{O_2}\).
\begin{itemize}
\item In condizioni di ipossia (alta quota), la sintesi di 2,3-BPG aumenta, facilitando il rilascio tissutale di \(\mathrm{O_2}\) (adattamento a breve termine).
\item L'emoglobina fetale (\(\mathrm{HbF} = \alpha_2\gamma_2\)) ha minore affinità per il 2,3-BPG rispetto all'Hb adulta, garantendo un'affinità maggiore per l'\(\mathrm{O_2}\) e permettendo il trasferimento di ossigeno dal sangue materno a quello fetale.
\end{itemize}

\section{Adattamenti all'ipossia e considerazioni cliniche}
\subsection{Adattamenti fisiologici}
Oltre all'aumento di 2,3-BPG, l'esposizione prolungata all'ipossia stimola la produzione renale di \textbf{eritropoietina (EPO)}, un ormone glicoproteico che promuove l'eritropoiesi nel midollo osseo, aumentando la capacità di trasporto di \(\mathrm{O_2}\) del sangue (adattamento a lungo termine).

\subsection{Beta-talassemia (anemia mediterranea)}
La \textbf{beta-talassemia} è una malattia genetica causata da mutazioni che riducono o aboliscono la sintesi delle catene \(\beta\) dell'emoglobina. Il risultato è la produzione di globuli rossi difettosi, anemia emolitica e necessità di trasfusioni frequenti. La comprensione della regolazione dell'Hb ha implicazioni terapeutiche per queste patologie.

\section{Sintesi funzionale}
\begin{itemize}
\item Emoglobina e mioglobina sono emoproteine essenziali per il trasporto e lo stoccaggio dell'\(\mathrm{O_2}\).
\item La struttura tetramerica dell'Hb e la cooperatività omotropica consentono un rilascio efficiente di \(\mathrm{O_2}\) nei tessuti.
\item Fattori allosterici eterotropici (pH, \(\mathrm{CO_2}\), temperatura, 2,3-BPG) modulano finemente l'affinità dell'Hb per l'\(\mathrm{O_2}\) in risposta alle esigenze metaboliche.
\item Adattamenti fisiologici (2,3-BPG, EPO) e patologie (talassemia) illustrano l'importanza clinica della regolazione del trasporto di ossigeno.
\end{itemize}
% =========================================================
% Fine Capitolo 5
% =========================================================

% =========================================================
% CAPITOLO 6: Enzimi
% =========================================================
\chapter{Enzimi}
\section{Introduzione e ruolo catalitico}
Gli \textbf{enzimi} sono proteine (con rare eccezioni rappresentate da RNA catalitici detti \textbf{ribozimi}) che svolgono la funzione di \textbf{catalizzatori biologici}. La loro presenza è indispensabile per accelerare le reazioni chimiche cellulari che, in condizioni fisiologiche, avverrebbero a velocità incompatibili con la vita.
Gli enzimi presentano tre caratteristiche fondamentali:
\begin{itemize}
\item \textbf{Non vengono consumati} durante la reazione e possono quindi catalizzare numerosi cicli catalitici.
\item \textbf{Non alterano l'equilibrio} della reazione: non influenzano la direzione termodinamica, ma accelerano il raggiungimento dell'equilibrio.
\item \textbf{Aumentano la velocità} sia della reazione diretta che di quella inversa, fino a un fattore di \(10^{10}\) volte o superiore.
\end{itemize}
Un esempio classico è la reazione catalizzata dalla lattato deidrogenasi (LDH):
\[
\text{Piruvato} + \text{NADH} \rightleftharpoons \text{Lattato} + \text{NAD}^+
\]

\section{Energia di attivazione e meccanismo d'azione}
Una reazione esoergonica (\(\Delta G < 0\)) è termodinamicamente favorita, ma ciò non garantisce che avvenga rapidamente. Per procedere, i reagenti devono superare una barriera energetica, raggiungendo uno \textbf{stato di transizione} ad alta energia. L'energia necessaria per raggiungere questo stato è definita \textbf{energia di attivazione} (\(E_a\)).
È possibile immaginare l'\(E_a\) come l'energia necessaria per spingere un masso (reagenti) sulla sommità di una collina (stato di transizione), prima che rotoli a valle (prodotti). Più alta è la \(E_a\), minore è la probabilità che i reagenti si orientino casualmente in modo favorevole alla reazione, risultando in una cinetica estremamente lenta.
La funzione esclusiva degli enzimi è \textbf{abbassare l'energia di attivazione}, stabilizzando lo stato di transizione e permettendo alla reazione di procedere a velocità fisiologicamente rilevanti, senza alterare il \(\Delta G\) complessivo.
\begin{esempio}[Fosforilazione del glucosio]
La prima tappa della glicolisi (glucosio \(\rightarrow\) glucosio-6-fosfato) ha un \(\Delta G\) negativo, ma un'elevata \(E_a\). In assenza di catalisi, la reazione sarebbe trascurabile. L'enzima \textbf{esochinasi} abbassa la \(E_a\) favorendo l'interazione tra glucosio e ATP, consentendo la rapida formazione dei prodotti.
\end{esempio}

\section{Specificità e modelli di riconoscimento del substrato}
Per abbassare l'\(E_a\), gli enzimi legano reversibilmente i substrati formando un \textbf{complesso enzima-substrato (ES)}. Il legame avviene in una regione specifica della proteina detta \textbf{sito attivo}, composta da residui amminoacidici coinvolti sia nel riconoscimento del substrato sia nella catalisi. Durante il processo si formano intermedi transitori (ES ed EP) prima del rilascio del prodotto.
\subsection{Modelli di interazione}
Storicamente, il \textbf{modello chiave-serratura} ipotizzava una complementarità geometrica rigida e preformata tra enzima e substrato. Tuttavia, un legame eccessivamente stabile impedirebbe il successivo rilascio del prodotto.
Il modello attualmente accettato è quello dell'\textbf{adattamento indotto}: il sito attivo non è perfettamente complementare al substrato nello stato libero. L'interazione iniziale (mediata da legami deboli e transitori) induce un cambiamento conformazionale nell'enzima che ottimizza il posizionamento del substrato, favorisce la catalisi e agevola il successivo distacco del prodotto. Un esempio è la chiusura conformazionale dell'esochinasi attorno al glucosio dopo il legame.
\subsection{Stereospecificità}
Oltre alla specificità chimica, gli enzimi sono altamente \textbf{stereospecifici}. Poiché le proteine sono sintetizzate esclusivamente da L-amminoacidi, i siti attivi presentano un'architettura chirale asimmetrica. Di conseguenza, riconoscono e trasformano selettivamente uno specifico stereoisomero, e rilasciano prodotti con configurazione spaziale definita. Esempio: la lattato deidrogenasi muscolare è specifica per l'\textbf{L-lattato}.

\section{Cinetica enzimatica e fattori di modulazione}
L'attività enzimatica è finemente modulata dalla cellula per adattare il flusso metabolico alle richieste energetiche. I principali fattori che influenzano la velocità di reazione sono: concentrazione del substrato, pH, temperatura e concentrazione dell'enzima.
\subsection{Concentrazione del substrato e costante di Michaelis}
In una reazione non catalizzata, la velocità aumenta linearmente con la concentrazione del substrato \([S]\). Nelle reazioni enzimatiche, la velocità aumenta linearmente solo a basse concentrazioni; oltre una certa soglia, la curva diventa \textbf{iperbolica} fino a raggiungere un plateau, la \textbf{velocità massima} (\(V_{\max}\)), quando tutti i siti attivi sono saturi.
La relazione è descritta dalla \textbf{costante di Michaelis} (\(K_M\)), definita come la concentrazione di substrato alla quale la velocità di reazione è pari a \(V_{\max}/2\).
\begin{itemize}
\item \(K_M\) \textbf{bassa} \(\rightarrow\) alta affinità enzima-substrato.
\item \(K_M\) è specifica per ogni coppia enzima-substrato.
\item In condizioni fisiologiche, \([S]\) è generalmente \(< K_M\), collocando l'enzima nella zona lineare della curva: piccole variazioni di substrato si traducono in ampie variazioni della velocità di reazione, garantendo una risposta metabolica sensibile.
\end{itemize}
\subsection{Effetto del pH e della temperatura}
\begin{itemize}
\item \textbf{pH}: La maggior parte degli enzimi presenta un optimum attorno a pH 7. Eccezioni includono enzimi digestivi gastrici (pH 1--2) e lisosomiali (pH 4,5--5). Deviazioni dal pH ottimale alterano lo stato di ionizzazione dei residui del sito attivo, riducendo l'efficienza catalitica. Durante esercizio intenso, il pH muscolare può scendere a 6,5, inibendo enzimi metabolici e contribuendo all'insorgenza della fatica.
\item \textbf{Temperatura}: L'aumento termico incrementa l'energia cinetica e la frequenza delle collisioni molecolari, accelerando la reazione fino a \(\approx 50^\circ\mathrm{C}\). Oltre questa soglia, i legami deboli che stabilizzano la struttura terziaria si rompono, causando \textbf{denaturazione} e perdita irreversibile dell'attività. Il riscaldamento pre-esercizio sfrutta questo principio per ottimizzare l'efficienza enzimatica muscolare.
\end{itemize}
\subsection{Concentrazione dell'enzima}
Quando l'enzima opera in condizioni di saturazione (\([S] \gg K_M\)), la velocità di reazione è direttamente proporzionale alla concentrazione dell'enzima \([E]\), mentre la \(K_M\) rimane invariata. La capacità delle cellule di aumentare l'espressione proteica enzimatica è il fondamento degli \textbf{adattamenti metabolici all'allenamento}, in particolare di tipo aerobico.

\section{Cofattori, coenzimi e vitamine}
Molti enzimi richiedono molecole non proteiche per esercitare la loro funzione catalitica:
\begin{itemize}
\item \textbf{Cofattori}: ioni inorganici essenziali (es. \(\mathrm{Cu^{2+}}, \mathrm{K^+}, \mathrm{Fe^{2+}}, \mathrm{Mg^{2+}}, \mathrm{Mn^{2+}}, \mathrm{Zn^{2+}}\)).
\item \textbf{Coenzimi}: molecole organiche che partecipano direttamente alla reazione, spesso subendo modificazioni temporanee (agendo come secondi substrati). Esempio: \(\mathrm{NAD^+}\) nella lattato deidrogenasi.
\item \textbf{Gruppo prostetico}: cofattore o coenzima legato stabilmente (spesso covalentemente) all'enzima (es. gruppo eme nell'emoglobina).
\end{itemize}
Molte \textbf{vitamine} (specie del gruppo B) sono precursori essenziali di coenzimi. Carenze vitaminiche compromettono l'attività enzimatica e causano patologie metaboliche. I coenzimi \(\mathrm{NAD^+}\) e \(\mathrm{FAD}\) sono fondamentali come accettori/donatori di elettroni nelle reazioni redox mediate dalle deidrogenasi.

\section{Classificazione e nomenclatura enzimatica}
Gli enzimi sono classificati in \textbf{6 classi principali} in base al tipo di reazione catalizzata. La maggior parte presenta il suffisso \textit{-asi}. Ogni enzima è identificato da un codice EC (Enzyme Commission) a 4 numeri (Classe.Sottoclasse.Sottosottoclasse.Seriale).
\begin{table}[htbp]
\centering
\caption{Classificazione principale degli enzimi}
\label{tab:enzimi-classi}
\renewcommand{\arraystretch}{1.2}
\begin{tabular}{@{} p{0.25\textwidth} p{0.70\textwidth} @{}}
\toprule
\textbf{Classe} & \textbf{Tipo di reazione} \\
\midrule
1. Ossidoreduttasi & Reazioni redox, trasferimento di atomi di idrogeno/elettroni \\
2. Transferasi & Trasferimento di gruppi funzionali (es. chinasi: trasferiscono fosfato) \\
3. Idrolasi & Rottura di legami per aggiunta di \(\mathrm{H_2O}\) (idrolisi) \\
4. Liasi & Eliminazione di gruppi per formare doppi legami, o addizione inversa \\
5. Isomerasi & Riarrangiamento intramolecolare (isomerizzazione) \\
6. Ligasi & Formazione di nuovi legami (C-C, C-O, C-N) accoppiata a idrolisi di ATP \\
\bottomrule
\end{tabular}
\end{table}

\section{Regolazione e inibizione dell'attività enzimatica}
Oltre a pH e temperatura, l'attività enzimatica è regolata finemente attraverso due meccanismi principali: modificazioni covalenti e regolazione allosterica.
\subsection{Modificazioni covalenti}
Consistono nell'aggiunta o rimozione reversibile di gruppi funzionali che alterano la conformazione e l'attività dell'enzima. Il meccanismo più comune è la \textbf{fosforilazione}: l'aggiunta di un gruppo fosfato (donato dall'ATP) a residui di Serina, Treonina o Tirosina, catalizzata da \textbf{chinasi}. La rimozione è operata da \textbf{fosfatasi} (defosforilazione). Sebbene la fosforilazione spesso attivi l'enzima, esistono eccezioni (es. inibizione della glicogeno sintasi).
\subsection{Regolazione allosterica}
Molecole chiamate \textbf{effettori allosterici} si legano a siti regolatori distinti dal sito attivo, inducendo cambiamenti conformazionali o elettrostatici che modulano l'affinità per il substrato. Gli effettori possono essere \textbf{positivi} (attivatori) o \textbf{negativi} (inibitori). Questo meccanismo è cruciale per il controllo fine della produzione di energia e biomolecole.
\subsection{Inibizione enzimatica}
L'attività può essere compromessa da molecole inibitrici:
\begin{itemize}
\item \textbf{Inibizione irreversibile}: legame covalente o estremamente stabile che inattiva permanentemente l'enzima (tipica di farmaci o tossine).
\item \textbf{Inibizione reversibile}: legame non covalente transitorio, fondamentale per la regolazione fisiologica.
\begin{itemize}
    \item \textit{Competitiva}: l'inibitore compete con il substrato per il sito attivo (strutturalmente simile), riducendo l'efficienza senza alterare \(V_{\max}\) (superabile aumentando \([S]\)).
    \item \textit{Non competitiva}: l'inibitore si lega a un sito alternativo, alterando la conformazione e riducendo \(V_{\max}\) senza influenzare \(K_M\).
\end{itemize}
\end{itemize}
% =========================================================
% Fine Capitolo 6
% =========================================================

% =========================================================
% CAPITOLO 7: Carboidrati
% =========================================================
\chapter{Carboidrati}
\section{Introduzione e funzioni biologiche}
I \textbf{carboidrati} (o \textit{glucidi}, \textit{zuccheri}, \textit{saccaridi}) sono macromolecole biologiche costituite da carbonio, idrogeno e ossigeno, con formula bruta generale \(\mathrm{C_n(H_2O)_n}\) (\(n \geq 3\)). Il ciclo energetico della biosfera si basa sulla loro interconversione: tramite la \textbf{fotosintesi}, gli organismi autotrofi fissano \(\mathrm{CO_2}\) e \(\mathrm{H_2O}\) utilizzando energia luminosa per sintetizzare glucosio e ossigeno; in assenza di luce, o negli organismi eterotrofi, queste molecole vengono ossidate per produrre energia sotto forma di \textbf{ATP}:
\[
\mathrm{C_n(H_2O)_n + nO_2 \rightarrow nCO_2 + nH_2O + Energia (ATP)}
\]
Oltre alla produzione energetica, i carboidrati svolgono funzioni fondamentali:
\begin{itemize}
\item \textbf{Riserva energetica}: amido (vegetali, cloroplasti/amiloplasti) e glicogeno (animali, fegato e muscolo);
\item \textbf{Strutturale}: cellulosa (pareti vegetali), chitina (esoscheletri), peptidoglicani (batteri);
\item \textbf{Informazionale e di riconoscimento}: glicoproteine e glicolipidi di membrana, coinvolti nella segnalazione cellulare e nel riconoscimento immunitario.
\end{itemize}
In ambito fisiologico, i carboidrati rappresentano il substrato preferenziale per la contrazione muscolare ad alta intensità. Le riserve di glicogeno muscolare si esauriscono in tempi inversamente proporzionali all'intensità dell'esercizio (\(\% \mathrm{VO_2max}\)), e la loro disponibilità è strettamente correlata alla performance aerobica prolungata.

\section{Struttura e classificazione dei monosaccaridi}
I \textbf{monosaccaridi} sono le unità fondamentali dei carboidrati. La loro struttura consiste in una catena carboniosa (3--7 atomi di C) contenente un gruppo \textbf{carbonilico} (\(\mathrm{C=O}\)):
\begin{itemize}
\item Se il carbonile è terminale (gruppo aldeidico, \(-\mathrm{CHO}\)), il monosaccaride è un \textbf{aldoso};
\item Se il carbonile è interno (gruppo chetonico, \(\mathrm{C=O}\)), è un \textbf{chetoso}.
\end{itemize}
I carboni rimanenti legano gruppi idrossilici (\(-\mathrm{OH}\)) e idrogeno. In base al numero di atomi di carbonio si distinguono: triosi (\(\mathrm{C_3H_6O_3}\), es. gliceraldeide), tetrosi, pentosi ed esosi (\(\mathrm{C_6H_{12}O_6}\), es. glucosio).
\subsection{Stereoisomeria e serie D/L}
I monosaccaridi presentano \textbf{centri chirali} (carboni legati a 4 gruppi diversi), eccetto il diidrossiacetone. Ciò genera stereoisomeri non sovrapponibili, detti \textbf{enantiomeri}. La serie di appartenenza (\textbf{D} o \textbf{L}) si determina confrontando la configurazione del carbonio chirale più lontano dal gruppo carbonilico con quella della gliceraldeide. In natura, i monosaccaridi sono quasi esclusivamente della serie \textbf{D}.
Quando due stereoisomeri differiscono per la configurazione di un solo centro chirale, sono detti \textbf{epimeri} (es. glucosio e galattosio sono epimeri al C4).
\subsection{Forma ciclica e anomeria}
In soluzione acquosa, i monosaccaridi con 5 o più atomi di carbonio ciclizzano spontaneamente tramite reazione intramolecolare tra il gruppo carbonilico e un gruppo \(-\mathrm{OH}\), formando un \textbf{emiacetale} o \textbf{emichetale}. Si ottengono:
\begin{itemize}
\item \textbf{Forme piranosiche}: anello a 6 atomi (simile al pirano);
\item \textbf{Forme furanosiche}: anello a 5 atomi (simile al furano).
\end{itemize}
Il carbonio che diventa chirale durante la ciclizzazione è detto \textbf{carbonio anomerico} (C1 negli aldoesosi). Può assumere due configurazioni spaziali:
\begin{itemize}
\item \(\boldsymbol{\alpha}\)-anomero: il gruppo \(-\mathrm{OH}\) anomerico è orientato \textit{sotto} il piano dell'anello;
\item \(\boldsymbol{\beta}\)-anomero: il gruppo \(-\mathrm{OH}\) anomerico è orientato \textit{sopra} il piano.
\end{itemize}
Le due forme sono in equilibrio dinamico con la forma lineare (mutarotazione). Per il D-glucosio, in soluzione prevale la forma \(\beta\) (\(\approx 64\%\)), seguita dalla \(\alpha\) (\(\approx 36\%\)) e da una quota minima della forma aperta (\(<1\%\)).

\section{Derivati dei monosaccaridi}
I monosaccaridi possono subire modificazioni strutturali che ne ampliano le funzioni biologiche:
\begin{itemize}
\item \textbf{Deossizuccheri}: sostituzione di un \(-\mathrm{OH}\) con un \(-\mathrm{H}\). Il \textbf{2-deossiribosio} è il componente dello scheletro del DNA, prodotto dalla riduzione del ribosio catalizzata dalla ribonucleotide reduttasi.
\item \textbf{Derivati ridotti (polioli)}: riduzione del gruppo carbonilico a gruppo \(-\mathrm{OH}\). Esempio: \textbf{glicerolo}, precursore dei trigliceridi. I polioli sono utilizzati come dolcificanti a basso indice glicemico e ridotta cariogenicità.
\item \textbf{Derivati amminici (amminozuccheri)}: sostituzione di un \(-\mathrm{OH}\) con un gruppo amminico (\(-\mathrm{NH_2}\)). Sono costituenti strutturali di glicolipidi, glicoproteine e polisaccaridi batterici.
\end{itemize}

\section{Oligosaccaridi e legame glicosidico}
Due o più monosaccaridi si uniscono tramite un \textbf{legame glicosidico}, formatosi per \textbf{condensazione} tra il carbonio anomerico di uno zucchero e un carbonio idrossilico di un altro, con rilascio di una molecola d'acqua. La natura del legame dipende dalla configurazione anomerica del primo residuo:
\begin{itemize}
\item \(\boldsymbol{\alpha}\)-\textbf{glicosidico}: se il C1 è in conformazione \(\alpha\);
\item \(\boldsymbol{\beta}\)-\textbf{glicosidico}: se il C1 è in conformazione \(\beta\).
\end{itemize}
Il legame più comune è il \((1\to4)\), ma nei polisaccaridi ramificati è fondamentale il \((1\to6)\).
\subsection{Disaccaridi principali}
\begin{itemize}
\item \textbf{Maltosio}: glucosio + glucosio, legame \(\alpha(1\to4)\).
\item \textbf{Cellobiosio}: glucosio + glucosio, legame \(\beta(1\to4)\). Unità ripetitiva della cellulosa.
\item \textbf{Lattosio}: galattosio + glucosio, legame \(\beta(1\to4)\). Presente nel latte.
\item \textbf{Saccarosio}: glucosio + fruttosio, legame \(\alpha(1\to2)\beta\). Zucchero da tavola.
\end{itemize}
La specificità enzimatica digestiva è legata alla configurazione del legame glicosidico. Gli enzimi umani (es. amilasi, maltasi) idrolizzano i legami \(\alpha\), ma non i legami \(\beta\) della cellulosa. Il lattosio viene idrolizzato dalla \textbf{lattasi} (\(\beta\)-galattosidasi). La \textbf{persistenza della lattasi} in età adulta è un adattamento genetico presente in alcune popolazioni; la sua assenza causa \textbf{intolleranza al lattosio}, con fermentazione batterica intestinale e produzione di gas (\(\mathrm{H_2}\), \(\mathrm{CH_4}\), \(\mathrm{CO_2}\)).

\section{Polisaccaridi}
I polisaccaridi sono polimeri di monosaccaridi legati da legami glicosidici. Si classificano in:
\begin{itemize}
\item \textbf{Omopolisaccaridi}: costituiti da un solo tipo di monomero (es. amido, glicogeno, cellulosa);
\item \textbf{Eteropolisaccaridi}: costituiti da due o più tipi di monomeri (es. glicosamminoglicani).
\end{itemize}
\subsection{Omopolisaccaridi del glucosio}
\begin{itemize}
\item \textbf{Amido} (riserva vegetale): composto da \textbf{amilosio} (catena lineare, legami \(\alpha(1\to4)\), \(\approx 20\%\)) e \textbf{amilopectina} (catena ramificata, legami \(\alpha(1\to4)\) e ramificazioni \(\alpha(1\to6)\) ogni 24--30 residui, \(\approx 80\%\)). Si organizza in granuli compatti e insolubili; la cottura ne aumenta la digeribilità.
\item \textbf{Cellulosa} (strutturale vegetale): polimero lineare non ramificato con legami \(\beta(1\to4)\). La configurazione \(\beta\) induce una conformazione estesa e lineare, favorendo la formazione di microfibrille resistenti tramite legami idrogeno intercatena. Indigeribile dall'uomo, ma fermentata dalla flora batterica dei ruminanti.
\item \textbf{Glicogeno} (riserva animale): struttura simile all'amilopectina, ma con ramificazioni \(\alpha(1\to6)\) molto più frequenti (ogni 8--10 residui). Questa elevata ramificazione aumenta la solubilità e fornisce numerosi siti terminali per l'attacco simultaneo degli enzimi glicogenolitici, garantendo un rilascio rapido di glucosio. Accumulato in granuli citoplasmatici, nel fegato mantiene la glicemia, mentre nel muscolo fornisce energia locale. Deficit enzimatici nel suo metabolismo causano \textbf{glicogenosi}.
\end{itemize}
\subsection{Eteropolisaccaridi e matrice extracellulare}
Gli eteropolisaccaridi sono componenti fondamentali della \textbf{matrice extracellulare} (MEC), un reticolo idratato che circonda le cellule, fornendo supporto strutturale, idratazione e vie di diffusione per nutrienti e segnali. La MEC è costituita da fibre proteiche (collagene, elastina, fibronectina, laminina) e da una rete polisaccaridica di \textbf{glicosamminoglicani (GAG)}.
I GAG sono polimeri lineari di disaccaridi ripetuti, spesso solforati o carichi negativamente, che legano grandi quantità di acqua, conferendo turgore e proprietà lubrificanti alla MEC. Sono spesso legati covalentemente a proteine core formando \textbf{proteoglicani}.
\begin{esempio}[Acido ialuronico]
Il più noto GAG, costituito dalla ripetizione di acido glucuronico e N-acetilglucosammina legati da legami \(\beta(1\to3)\) e \(\beta(1\to4)\). Non è solforato, non è legato covalentemente a proteine e può raggiungere lunghezze elevate (\(\approx 50.000\) unità). Abbondante in cartilagini, tendini, liquido sinoviale e umor vitreo, svolge funzioni di ammortizzazione meccanica, lubrificazione articolare e mantenimento dell'idratazione tissutale. La sua degradazione con l'età contribuisce alla comparsa di rughe e artrosi.
\end{esempio}
% =========================================================
% Fine Capitolo 7
% =========================================================

% =========================================================
% CAPITOLO 8: Lipidi
% =========================================================
\chapter{Lipidi}
\section{Introduzione e classificazione}
I \textbf{lipidi} costituiscono una classe eterogenea di composti organici caratterizzati da una marcata \textbf{idrofobicità} e insolubilità in acqua. La loro struttura è prevalentemente costituita da idrocarburi, ovvero catene contenenti solo atomi di carbonio e idrogeno. Nell'organismo umano, i lipidi rappresentano circa il \textbf{17\% del peso corporeo} e svolgono funzioni diversificate, classificabili in tre categorie principali:
\begin{itemize}
\item \textbf{Lipidi di riserva}: trigliceridi (grassi e oli), deputati all'immagazzinamento di energia chimica;
\item \textbf{Lipidi strutturali}: fosfolipidi, sfingolipidi e steroli, componenti fondamentali delle membrane biologiche;
\item \textbf{Lipidi regolatori}: presenti in piccole quantità, fungono da precursori di vitamine liposolubili, ormoni steroidei e molecole di segnalazione (es. eicosanoidi).
\end{itemize}
\subsection{Ruolo metabolico nell'esercizio fisico}
A differenza dei carboidrati, i lipidi non possono essere utilizzati come fonte energetica primaria durante esercizi ad alta intensità, poiché il loro catabolismo ossidativo richiede necessariamente ossigeno e tempi di attivazione più lunghi. Tuttavia, rappresentano la fonte energetica predominante durante:
\begin{itemize}
\item Esercizi aerobici prolungati a bassa-moderata intensità;
\item Fasi di recupero tra sessioni di allenamento;
\item Condizioni di digiuno o deplezione glicogenica.
\end{itemize}
Il concetto di \textbf{crossover point} descrive l'intensità di esercizio alla quale il contributo energetico dei carboidrati eguaglia quello dei lipidi: aumentando l'intensità, la quota lipidica diminuisce progressivamente a favore dei glucidi. L'allenamento aerobico cronico induce adattamenti metabolici che spostano il crossover point verso intensità maggiori, migliorando l'efficienza nell'ossidazione lipidica.

\section{Lipidi semplici: acilgliceroli e cere}
\subsection{Struttura degli acilgliceroli}
Gli \textbf{acilgliceroli} sono esteri formati dalla condensazione tra una molecola di \textbf{glicerolo} (un triolo derivato metabolicamente dai carboidrati) e uno o più \textbf{acidi grassi}. In base al numero di acidi grassi esterificati, si distinguono:
\begin{itemize}
\item \textbf{Monoacilgliceroli}: un solo acido grasso legato al glicerolo;
\item \textbf{Diacilgliceroli}: due acidi grassi legati;
\item \textbf{Triacilgliceroli (TAG) o trigliceridi}: tre acidi grassi legati, forma principale di deposito lipidico nell'organismo.
\end{itemize}
\begin{definizione}[Trigliceridi]
I \textbf{trigliceridi} sono esteri del glicerolo con tre acidi grassi. Rappresentano la forma di immagazzinamento energetico più efficiente, grazie all'elevato contenuto di legami C-H ridotti e all'assenza di acqua di idratazione.
\end{definizione}
\subsection{Immagazzinamento e mobilizzazione}
I trigliceridi sono conservati sotto forma di \textbf{gocce lipidiche} nel citosol cellulare:
\begin{itemize}
\item \textbf{Adipociti}: cellule specializzate del tessuto adiposo, contenenti una singola grande goccia lipidica che occupa la maggior parte del volume cellulare;
\item \textbf{Fibre muscolari scheletriche}: contengono gocce lipidiche più piccole e numerose, posizionate strategicamente vicino ai mitocondri per un rapido accesso energetico durante l'esercizio.
\end{itemize}
Per essere utilizzati a scopo energetico, gli acidi grassi devono essere liberati dal glicerolo tramite il processo di \textbf{lipolisi}, catalizzato da enzimi specifici detti \textbf{lipasi} (es. lipasi ormone-sensibile, HSL). Gli acidi grassi liberi (FFA) vengono quindi trasportati nel sangue legati all'\textbf{albumina} (fino a 7 molecole di FFA per molecola di albumina) e captati dai tessuti target tramite trasportatori di membrana specifici (es. FAT/CD36, FABPpm).
\begin{nota}[Lipotossicità]
Un eccesso cronico di acidi grassi liberi circolanti può indurre \textbf{lipotossicità}, con accumulo di intermedi lipidici (diacilgliceroli, ceramidi) in tessuti non adiposi (fegato, muscolo, pancreas), compromettendo la segnalazione insulinica e la funzione cellulare.
\end{nota}

\section{Acidi grassi: struttura, nomenclatura e proprietà}
\subsection{Struttura chimica}
Gli \textbf{acidi grassi} sono acidi carbossilici a catena alifatica lineare, caratterizzati da:
\begin{itemize}
\item Un \textbf{gruppo carbossilico} terminale (\(-\mathrm{COOH}\)), polare e ionizzabile;
\item Una \textbf{catena idrocarburica} apolare, generalmente composta da 4 a 36 atomi di carbonio (più comunemente 12--24).
\end{itemize}
La maggior parte degli acidi grassi naturali presenta un numero \textbf{pari} di atomi di carbonio, conseguenza del meccanismo biosintetico che aggiunge unità bicarboniose (acetil-CoA).
\subsection{Nomenclatura}
La nomenclatura standard degli acidi grassi segue il formato:
\[
\text{C:X } \Delta^{\text{posizioni}}
\]
dove:
\begin{itemize}
\item \textbf{C} = numero di atomi di carbonio;
\item \textbf{X} = numero di doppi legami;
\item \textbf{\(\Delta^{\text{posizioni}}\)} = posizione dei doppi legami, contando dal carbonio carbossilico (C1).
\end{itemize}
\begin{esempio}[Esempi di nomenclatura]
\begin{itemize}
\item \textbf{Acido stearico}: \(18:0\) (saturo, 18 C, nessun doppio legame);
\item \textbf{Acido oleico}: \(18:1\,\Delta^9\) (monoinsaturo, doppio legame in posizione 9);
\item \textbf{Acido \(\alpha\)-linolenico (ALA)}: \(18:3\,\Delta^{9,12,15}\) (polinsaturo omega-3).
\end{itemize}
\end{esempio}
\subsection{Saturazione e isomeria geometrica}
In base alla presenza di doppi legami nella catena idrocarburica, gli acidi grassi si classificano in:
\begin{itemize}
\item \textbf{Saturi}: assenza di doppi legami (es. acido palmitico \(16:0\), acido stearico \(18:0\));
\item \textbf{Monoinsaturi}: un solo doppio legame (es. acido oleico \(18:1\));
\item \textbf{Polinsaturi}: due o più doppi legami (es. acido linoleico \(18:2\), acido arachidonico \(20:4\)).
\end{itemize}
Il doppio legame introduce rigidità conformazionale e genera \textbf{isomeria geometrica cis-trans}:
\begin{itemize}
\item \textbf{Isomero \textit{cis}}: gli atomi di carbonio adiacenti al doppio legame si trovano dallo \textit{stesso lato}; la catena presenta una \textbf{piega} che impedisce l'impaccamento compatto. La maggior parte degli acidi grassi insaturi naturali è in configurazione \textit{cis}.
\item \textbf{Isomero \textit{trans}}: gli atomi di carbonio si trovano da \textit{lati opposti}; la catena mantiene una conformazione lineare simile ai saturi.
\end{itemize}
\begin{definizione}[Acidi grassi essenziali]
L'organismo umano non possiede gli enzimi (desaturasi) per introdurre doppi legami oltre il carbonio 9 (contando dal carbossile). Pertanto, gli acidi grassi con insaturazioni in posizioni \(\omega\)-3 e \(\omega\)-6 devono essere assunti con la dieta e sono definiti \textbf{essenziali}:
\begin{itemize}
\item \textbf{Omega-3}: primo doppio legame al carbonio 3 dall'estremità metilica (\(\omega\)); es. acido \(\alpha\)-linolenico (ALA, \(18:3\,\omega\)-3);
\item \textbf{Omega-6}: primo doppio legame al carbonio 6 dall'estremità \(\omega\); es. acido linoleico (LA, \(18:2\,\omega\)-6).
\end{itemize}
\end{definizione}
\subsection{Proprietà fisico-chimiche}
Le caratteristiche strutturali degli acidi grassi influenzano direttamente le loro proprietà:
\begin{itemize}
\item \textbf{Punto di fusione}: gli acidi grassi saturi e \textit{trans}, potendo impaccarsi in strutture ordinate, presentano punti di fusione più elevati (solidi a temperatura ambiente, es. burro). Gli insaturi \textit{cis}, con catene piegate, hanno punti di fusione più bassi (liquidi, es. oli vegetali).
\item \textbf{Lunghezza della catena}: catene più lunghe aumentano le interazioni idrofobiche e il punto di fusione (es. grasso di bue vs. burro).
\item \textbf{Solubilità}: il gruppo carbossilico conferisce una modesta idrofilicità, ma la predominanza della catena apolare rende gli acidi grassi a catena lunga (\(>10\) C) praticamente insolubili in acqua. Solo gli acidi a catena molto corta (4--6 C) mostrano una limitata solubilità acquosa.
\end{itemize}
\subsection{Idrogenazione e acidi grassi trans}
Il processo industriale di \textbf{idrogenazione} aggiunge idrogeno ai doppi legami degli acidi grassi insaturi:
\begin{itemize}
\item \textbf{Idrogenazione totale}: saturazione completa, conversione in acidi grassi saturi;
\item \textbf{Idrogenazione parziale}: saturazione incompleta, con isomerizzazione \textit{cis} \(\rightarrow\) \textit{trans} di alcuni doppi legami residui.
\end{itemize}
L'idrogenazione parziale è utilizzata per aumentare la stabilità ossidativa e la consistenza solida degli oli vegetali (es. margarine). Tuttavia, gli \textbf{acidi grassi trans} industriali sono associati a un aumentato rischio cardiovascolare, poiché elevano il colesterolo LDL ("cattivo") e riducono l'HDL ("buono"), con effetti più marcati rispetto ai grassi saturi.

\section{Lipoproteine e trasporto plasmatico dei lipidi}
I lipidi, essendo insolubili in ambiente acquoso, richiedono sistemi di trasporto specializzati nel plasma. Le \textbf{lipoproteine} sono complessi supramolecolari costituiti da:
\begin{itemize}
\item \textbf{Lipidi}: trigliceridi, colesterolo esterificato, fosfolipidi, colesterolo libero;
\item \textbf{Apolipoproteine}: componenti proteiche con funzioni strutturali, di riconoscimento recettoriale e di attivazione enzimatica.
\end{itemize}
Le interazioni tra i componenti sono prevalentemente non covalenti (idrofobiche, elettrostatiche).
\begin{table}[htbp]
\centering
\caption{Classificazione delle principali lipoproteine}
\label{tab:lipoproteine}
\renewcommand{\arraystretch}{1.2}
\begin{tabular}{@{} p{0.25\textwidth} p{0.70\textwidth} @{}}
\toprule
\textbf{Lipoproteina} & \textbf{Origine e funzione principale} \\
\midrule
\textbf{Chilomicroni} & Sintetizzati dall'enterocita intestinale; trasportano trigliceridi alimentari ai tessuti periferici (muscolo, adiposo) \\
\textbf{VLDL} (Very Low Density Lipoprotein) & Prodotte dal fegato; veicolano trigliceridi endogeni ai tessuti extraepatici \\
\textbf{LDL} (Low Density Lipoprotein) & Derivano dal catabolismo delle VLDL; trasportano colesterolo esterificato ai tessuti periferici (elevati livelli associati a rischio cardiovascolare) \\
\textbf{HDL} (High Density Lipoprotein) & Prodotte da fegato e intestino; mediano il \textit{trasporto inverso del colesterolo} dai tessuti periferici al fegato per l'eliminazione (effetto protettivo) \\
\bottomrule
\end{tabular}
\end{table}
La densità delle lipoproteine è inversamente proporzionale al contenuto lipidico: i chilomicroni sono i meno densi (più ricchi di trigliceridi), le HDL le più dense (più ricche di proteine). Le LDL, pur non essendo le più dense, contengono la maggiore percentuale di colesterolo e sono comunemente indicate come "colesterolo cattivo" per il loro ruolo nell'aterogenesi.

\section{Funzioni biologiche dei trigliceridi}
\begin{enumerate}
\item \textbf{Riserva energetica}: l'ossidazione completa degli acidi grassi produce circa \(9\,\mathrm{kcal/g}\), più del doppio rispetto a carboidrati e proteine (\(\approx 4\,\mathrm{kcal/g}\)), grazie allo stato di riduzione più elevato degli atomi di carbonio nella catena idrocarburica.
\item \textbf{Efficienza di stoccaggio}: l'idrofobicità dei trigliceridi permette un accumulo anidro, occupando un volume cellulare minimo rispetto al glicogeno (che lega \(\approx 2\,\mathrm{g}\) di acqua per grammo).
\item \textbf{Isolamento termico}: il tessuto adiposo sottocutaneo funge da barriera termica, riducendo la dispersione di calore corporeo.
\item \textbf{Protezione meccanica}: i cuscinetti adiposi proteggono organi vitali (reni, cuore) da traumi fisici.
\end{enumerate}

\section{Cere}
Le \textbf{cere} sono lipidi semplici non polari costituiti dall'esterificazione di un acido grasso a catena lunga con un alcol a catena lunga (non glicerolo). Grazie alla loro elevata idrofobicità, svolgono funzioni protettive di \textbf{impermeabilizzazione} in:
\begin{itemize}
\item Cuticole vegetali (foglie, frutti);
\item Piume e piumini degli uccelli;
\item Esoscheletri di insetti e artropodi;
\item Cera d'api (produzione dell'alveare).
\end{itemize}
% =========================================================
% Fine Capitolo 8: Lipidi semplici e trasporto
% =========================================================

% =========================================================
% CAPITOLO 9: Nucleotidi
% =========================================================
\chapter{Nucleotidi}
\section{Introduzione e funzioni biologiche}
I \textbf{nucleotidi} costituiscono una classe eterogenea di molecole organiche con struttura simile ma funzioni estremamente diversificate. Nel metabolismo cellulare svolgono ruoli fondamentali:
\begin{itemize}
\item \textbf{Costituenti degli acidi nucleici}: monomeri di DNA e RNA;
\item \textbf{Trasportatori di energia chimica}: ATP, GTP, CTP, UTP;
\item \textbf{Secondi messaggeri}: AMP ciclico (cAMP), GMP ciclico (cGMP);
\item \textbf{Trasferimento di elettroni}: coenzimi redox (NAD$^+$, NADP$^+$, FAD);
\item \textbf{Gruppi prostetici e cofattori}: Coenzima A (CoA).
\end{itemize}

\section{Struttura generale dei nucleotidi}
Un nucleotide è composto da tre componenti fondamentali:
\begin{enumerate}
\item \textbf{Base azotata}: molecola eterociclica aromatica contenente azoto;
\item \textbf{Pentoso}: zucchero a 5 atomi di carbonio (ribosio o desossiribosio);
\item \textbf{Gruppo fosfato}: uno o più residui fosforici legati al C5' del pentoso.
\end{enumerate}
\begin{definizione}[Nucleoside vs Nucleotide]
Un \textbf{nucleoside} è formato esclusivamente dalla base azotata legata al pentoso tramite legame \textbf{$\beta$-N-glicosidico}. L'aggiunta di uno o più gruppi fosfato al C5' del nucleoside genera il \textbf{nucleotide}.
\end{definizione}
\subsection{Le basi azotate}
Le basi azotate si classificano in due famiglie strutturali:
\begin{itemize}
\item \textbf{Purine} (struttura a doppio anello):
\begin{itemize}
    \item \textbf{Adenina (A)}: 6-amminopurina;
    \item \textbf{Guanina (G)}: 2-ammino-6-ossipurina.
\end{itemize}
\item \textbf{Pirimidine} (struttura ad anello singolo):
\begin{itemize}
    \item \textbf{Citosina (C)}: 2-ossi-4-amminopirimidina;
    \item \textbf{Timina (T)}: 2,4-diossi-5-metilpirimidina (presente solo nel DNA);
    \item \textbf{Uracile (U)}: 2,4-diossipirimidina (presente solo nell'RNA).
\end{itemize}
\end{itemize}
\subsection{Il pentoso e il legame N-glicosidico}
Lo zucchero presente nei nucleotidi è un aldopentoso:
\begin{itemize}
\item \textbf{Ribosio}: presente nell'RNA e nei nucleotidi energetici/regolatori;
\item \textbf{2-Desossiribosio}: presente nel DNA, ottenuto per riduzione del ribosio a livello del C2'.
\end{itemize}
La base azotata si lega al C1' del pentoso tramite un legame \textbf{$\beta$-N-glicosidico}:
\begin{itemize}
\item Per le purine: legame con l'azoto in posizione N9;
\item Per le pirimidine: legame con l'azoto in posizione N1.
\end{itemize}
I nucleosidi derivati sono: adenosina, guanosina, citidina, uridina (RNA) e deossiadenosina, deossiguanosina, deossicitidina, deossitimidina (DNA).
\subsection{I gruppi fosfato}
I nucleotidi contengono uno o più gruppi fosfato legati all'ossidrile del C5' del pentoso:
\begin{itemize}
\item \textbf{Nucleoside monofosfato (NMP)}: un gruppo fosfato legato tramite legame \textbf{fosfoestereo};
\item \textbf{Nucleoside difosfato (NDP)}: due gruppi fosfato, il secondo legato tramite legame \textbf{fosfoanidridico};
\item \textbf{Nucleoside trifosfato (NTP)}: tre gruppi fosfato, con due legami fosfoanidridici.
\end{itemize}
\begin{nota}[Energia dei legami fosfoanidridici]
L'idrolisi dei legami fosfoanidridici (tra i gruppi fosfato) rilascia una quantità significativa di energia ($\approx 30\,\mathrm{kJ/mol}$ per legame), rendendo i nucleotidi trifosfato eccellenti donatori di energia per le reazioni endoergoniche cellulari.
\end{nota}

\section{Nucleotidi modificati: desossinucleotidi}
I desossinucleotidi, precursori del DNA, derivano dalla riduzione del ribosio a livello del C2':
\[
\text{Ribonucleoside difosfato (NDP)} \xrightarrow{\text{Ribonucleotide reduttasi}} \text{Desossiribonucleoside difosfato (dNDP)}
\]
Questa reazione, catalizzata dall'enzima \textbf{ribonucleotide reduttasi}, è essenziale per la sintesi del DNA e strettamente regolata per mantenere l'equilibrio tra i pool di nucleotidi ribo- e desossiribo-.

\section{Acidi nucleici: DNA e RNA}
I nucleotidi si polimerizzano per formare gli acidi nucleici tramite legami \textbf{fosfodiesterici} tra il gruppo fosfato in posizione 5' di un nucleotide e l'ossidrile in posizione 3' del nucleotide successivo.
\begin{definizione}[Polarità degli acidi nucleici]
La catena polinucleotidica presenta polarità direzionale: un'estremità libera con gruppo fosfato 5' e un'estremità libera con ossidrile 3'. Per convenzione, le sequenze si scrivono in direzione \textbf{5' $\rightarrow$ 3'}.
\end{definizione}
\subsection{Acido desossiribonucleico (DNA)}
\begin{itemize}
\item Struttura a \textbf{doppia elica} formata da due filamenti antiparalleli;
\item Zucchero: \textbf{2-desossiribosio};
\item Basi: A, G, C, T (desossinucleotidi);
\item I due filamenti sono tenuti insieme da \textbf{legami idrogeno} tra basi complementari:
\[
\mathrm{A \equiv T} \quad (2\ \text{legami H}); \qquad \mathrm{G \equiv C} \quad (3\ \text{legami H})
\]
\item La sequenza nucleotidica codifica l'informazione genetica per la sintesi proteica.
\end{itemize}
\subsection{Acido ribonucleico (RNA)}
\begin{itemize}
\item Generalmente a \textbf{filamento singolo}, ma può formare strutture secondarie locali (stem-loop) stabilizzate da appaiamenti intramolecolari;
\item Zucchero: \textbf{ribosio};
\item Basi: A, G, C, U (ribonucleotidi);
\item Funzioni diversificate: mRNA (trasporto dell'informazione), tRNA (adattatore nella traduzione), rRNA (componente strutturale dei ribosomi), RNA regolatori.
\end{itemize}

\section{Nucleotidi come trasportatori di energia}
\subsection{ATP: la moneta energetica}
L'\textbf{adenosina trifosfato (ATP)} è il principale vettore di energia chimica nella cellula. L'idrolisi dei suoi legami fosfoanidridici:
\[
\mathrm{ATP + H_2O \rightarrow ADP + P_i} \quad \Delta G^\circ \approx -30\,\mathrm{kJ/mol}
\]
fornisce l'energia necessaria per:
\begin{itemize}
\item Lavoro meccanico (contrazione muscolare);
\item Trasporto attivo attraverso le membrane;
\item Sintesi di macromolecole (reazioni endoergoniche accoppiate).
\end{itemize}
\begin{esempio}[Accoppiamento energetico]
La sintesi del glucosio-6-fosfato nella prima tappa della glicolisi è endoergonica:
\[
\text{Glucosio} + \mathrm{P_i} \rightarrow \text{Glucosio-6-fosfato} \quad \Delta G^\circ > 0
\]
L'accoppiamento con l'idrolisi dell'ATP rende la reazione complessiva esoergonica:
\[
\text{Glucosio} + \mathrm{ATP} \rightarrow \text{Glucosio-6-fosfato} + \mathrm{ADP} \quad \Delta G^\circ < 0
\]
\end{esempio}
\subsection{Altri nucleotidi energetici}
\begin{itemize}
\item \textbf{GTP}: energia equivalente all'ATP, ma utilizzato preferenzialmente in vie di segnalazione (proteine G), sintesi proteica e gluconeogenesi;
\item \textbf{UTP}: attivatore del glucosio nella sintesi del glicogeno (formazione di UDP-glucosio);
\item \textbf{CTP}: coinvolto nella sintesi dei fosfolipidi di membrana.
\end{itemize}

\section{Nucleotidi come secondi messaggeri}
\subsection{AMP ciclico (cAMP)}
L'\textbf{adenosina 3',5'-monofosfato ciclico (cAMP)} è un nucleotide modificato che funge da \textbf{secondo messaggero} in numerose vie di segnalazione ormonale (es. adrenalina, glucagone).
\begin{itemize}
\item Sintetizzato dall'\textbf{adenilato ciclasi} a partire dall'ATP, con rimozione dei fosfati $\beta$ e $\gamma$ e ciclizzazione del fosfato $\alpha$ con il C3' del ribosio;
\item Attiva la \textbf{protein chinasi A (PKA)}, che fosforila proteine bersaglio modulandone l'attività;
\item Degradato dalla \textbf{fosfodiesterasi} ad AMP, terminando il segnale.
\end{itemize}
\subsection{GMP ciclico (cGMP)}
Strutturalmente analogo al cAMP, il \textbf{cGMP} (sintetizzato dalla guanilato ciclasi) media segnali come quelli dell'ossido nitrico (NO) nella vasodilatazione e nella fototrasduzione retinica.

\section{Nucleotidi come trasportatori di elettroni}
\subsection{Coenzimi nicotinammidici: NAD$^+$ e NADP$^+$}
La \textbf{nicotinammide adenina dinucleotide (NAD$^+$)} e la sua forma fosforilata \textbf{NADP$^+$} derivano dalla vitamina B3 (niacina) e fungono da accettori/donatori di elettroni nelle reazioni redox.
\begin{definizione}[Reazioni redox mediate da NAD(P)$^+$]
\[
\mathrm{NAD^+ + 2H^+ + 2e^- \rightleftharpoons NADH + H^+}
\]
\begin{itemize}
\item \textbf{NADH}: prodotto nel catabolismo (glicolisi, ciclo di Krebs), dona elettroni alla catena respiratoria mitocondriale per la sintesi di ATP;
\item \textbf{NADPH}: generato principalmente dalla via dei pentoso-fosfati, fornisce potere riducente per le biosintesi (es. acidi grassi) e per i sistemi antiossidanti (glutatione ridotto).
\end{itemize}
\end{definizione}
\subsection{Coenzimi flavinici: FAD e FMN}
La \textbf{flavina adenina dinucleotide (FAD)} e la \textbf{flavina mononucleotide (FMN)} derivano dalla vitamina B2 (riboflavina) e partecipano a reazioni redox come cofattori di flavoproteine.
\[
\mathrm{FAD + 2H^+ + 2e^- \rightleftharpoons FADH_2}
\]
Il \textbf{FADH$_2$} prodotto nel ciclo di Krebs e nella $\beta$-ossidazione degli acidi grassi dona elettroni alla catena respiratoria (complesso II), contribuendo alla sintesi di ATP.

\section{Coenzima A (CoA)}
Il \textbf{Coenzima A} è un derivato della vitamina B5 (acido pantotenico) e svolge un ruolo centrale nel metabolismo intermedio.
\begin{itemize}
\item Struttura: 3'-fosfoadenosina difosfato + acido pantotenico + $\beta$-mercaptoetilamina;
\item Il gruppo tiolico ($-\mathrm{SH}$) terminale forma legami \textbf{tioesterici} ad alta energia con gruppi acilici;
\item \textbf{Acetil-CoA}: intermedio chiave che connette il catabolismo di glucidi, lipidi e amminoacidi con il ciclo di Krebs e le vie biosintetiche.
\end{itemize}
\[
\text{Acetato} + \mathrm{CoA-SH} + \mathrm{ATP} \rightarrow \text{Acetil-CoA} + \mathrm{AMP} + \mathrm{PP_i}
\]

\section{Sintesi funzionale}
\begin{itemize}
\item I nucleotidi sono molecole versatili: costituenti degli acidi nucleici, trasportatori di energia e elettroni, messaggeri intracellulari e cofattori enzimatici;
\item La struttura di base (base + pentoso + fosfato) è modulata per adattarsi a funzioni specifiche;
\item Molti nucleotidi derivano da vitamine del gruppo B: carenze vitaminiche compromettono il metabolismo energetico e la sintesi di DNA/RNA;
\item La regolazione del pool nucleotidico è essenziale per l'omeostasi cellulare, la proliferazione e la risposta agli stimoli.
\end{itemize}
% =========================================================
% Fine Capitolo 9
% =========================================================

% =========================================================
% CAPITOLO 10: Introduzione al Metabolismo Energetico
% =========================================================
\chapter{Introduzione al Metabolismo Energetico}
\section{Introduzione e funzioni del metabolismo}
Il \textbf{metabolismo} rappresenta l'insieme di tutte le reazioni chimiche che avvengono nelle cellule e sono necessarie al loro funzionamento. Queste reazioni sono organizzate in \textbf{vie metaboliche} (es. glicolisi, ciclo di Krebs), che convertono specifici substrati in prodotti attraverso una serie di passaggi catalizzati da enzimi. Durante queste trasformazioni si generano numerose molecole intermedie dette \textbf{metaboliti}. Le vie metaboliche sono strettamente interconnesse e finemente regolate per adattarsi alle richieste cellulari.
Le principali funzioni del metabolismo sono:
\begin{enumerate}
\item \textbf{Produzione di energia chimica}: catturata dalla luce solare (fotosintesi) o ottenuta dall'ossidazione dei nutrienti;
\item \textbf{Conversione dei nutrienti}: trasformazione delle molecole assunte con la dieta nei precursori delle macromolecole cellulari;
\item \textbf{Biosintesi}: polimerizzazione di monomeri (amminoacidi, monosaccaridi, nucleotidi) in macromolecole complesse (proteine, polisaccaridi, acidi nucleici).
\end{enumerate}

\section{Catabolismo e Anabolismo}
Il metabolismo si suddivide in due processi complementari ma distinti:
\begin{itemize}
\item \textbf{Catabolismo}: via degradativa. Le macromolecole complesse vengono ossidate e scomposte in molecole più semplici o completamente mineralizzate a \(\mathrm{CO_2}\) e \(\mathrm{H_2O}\). Durante questo processo, l'energia immagazzinata nei legami chimici viene rilasciata e utilizzata per sintetizzare ATP. I coenzimi nucleotidici (\(\mathrm{NAD^+}\), \(\mathrm{NADP^+}\), \(\mathrm{FAD}\)) acquisiscono equivalenti di riduzione (elettroni e \(\mathrm{H^+}\)), trasformandosi nelle forme ridotte (\(\mathrm{NADH}\), \(\mathrm{NADPH}\), \(\mathrm{FADH_2}\)).
\item \textbf{Anabolismo}: via biosintetica. Molecole semplici (precursori) ed energia (ATP) vengono utilizzate per costruire macromolecole complesse. Le reazioni anaboliche sono riduttive: consumano \(\mathrm{NADH}\)/\(\mathrm{NADPH}\) e \(\mathrm{FADH_2}\), che vengono riconvertiti nelle forme ossidate.
\end{itemize}
I due processi sono interdipendenti e mantenuti in equilibrio dinamico.

\section{Le tre fasi del catabolismo}
La degradazione dei nutrienti avviene attraverso tre stadi progressivi:
\begin{enumerate}
\item \textbf{Fase 1 (Digestione e assorbimento)}: Le macromolecole alimentari (proteine, amido, lipidi) vengono idrolizzate in unità monomeriche (amminoacidi, glucosio, acidi grassi, glicerolo) e assorbite dalle cellule.
\item \textbf{Fase 2 (Conversione in intermedi)}: I monomeri vengono degradati ulteriormente fino a formare \textbf{acetil-CoA} o altri intermedi metabolici chiave (es. \(\alpha\)-chetoglutarato, succinil-CoA, fumarato).
\item \textbf{Fase 3 (Ossidazione completa)}: L'acetil-CoA entra nel \textbf{ciclo di Krebs}, dove viene completamente ossidato a \(\mathrm{CO_2}\). I coenzimi ridotti prodotti (\(\mathrm{NADH}\), \(\mathrm{FADH_2}\)) alimentano la \textbf{fosforilazione ossidativa} mitocondriale, generando la maggior parte dell'ATP cellulare.
\end{enumerate}
Gli intermedi catabolici, oltre all'ATP e agli elettroni conservati nei coenzimi, possono essere dirottati verso vie anaboliche per sintetizzare nuove macromolecole, dimostrando la notevole \textbf{plasticità metabolica} cellulare.

\section{Differenze fondamentali e compartimentalizzazione}
Sebbene catabolismo e anabolismo possano apparire come processi inversi, presentano differenze strutturali e regolatorie fondamentali:
\begin{itemize}
\item \textbf{Non sono percorsi speculari}: utilizzano enzimi diversi nelle tappe chiave, permettendo una regolazione indipendente. Esempio: la degradazione del glicogeno a lattato richiede 12 enzimi, mentre la gluconeogenesi (via inversa) ne utilizza 9.
\item \textbf{Bilancio energetico}: le vie cataboliche sono esoergoniche (\(\Delta G < 0\)), mentre quelle anaboliche sono endoergoniche e consumano generalmente più energia di quanta ne venga rilasciata dalla degradazione diretta.
\item \textbf{Regolazione fine}: i due percorsi sono controllati in modo antagonista per evitare \textbf{cicli futili} (sintesi e degradazione simultanea senza guadagno netto, che sprecherebbe ATP).
\end{itemize}
Inoltre, le vie metaboliche sono \textbf{compartimentalizzate} all'interno della cellula (citoplasma, mitocondri, nucleo, reticolo endoplasmatico, ecc.). Questa organizzazione spaziale permette di isolare reazioni incompatibili, ottimizzare le concentrazioni dei substrati e regolare finemente i flussi metabolici in base alle esigenze tissutali.

\section{Il ruolo dell'ATP e i sistemi di rigenerazione}
L'energia chimica dei nutrienti viene convertita in \textbf{lavoro cellulare} (chimico, meccanico, di trasporto) principalmente attraverso l'\textbf{ATP} (adenosina trifosfato). L'idrolisi dell'ATP:
\[
\mathrm{ATP + H_2O \rightarrow ADP + P_i} \quad \Delta G^\circ \approx -30{,}5\,\mathrm{kJ/mol}
\]
rilascia energia utilizzabile direttamente per la contrazione muscolare o accoppiata a reazioni endoergoniche. La quantità di ATP libero nelle cellule è estremamente ridotta (\(\approx 40\text{--}50\,\mathrm{g}\) nel muscolo), sufficiente solo per \(2\text{--}4\,\mathrm{s}\) di attività intensa. Pertanto, l'ATP deve essere continuamente rigenerato tramite tre sistemi principali:
\begin{enumerate}
\item \textbf{Fosfageni (fosfocreatina, PCr)}: via anaerobica alattacida. La \textbf{creatina chinasi} catalizza il trasferimento di un gruppo fosfato dalla PCr all'ADP:
\[
\mathrm{PCr + ADP \rightleftharpoons ATP + Creatina}
\]
Previene l'accumulo di ADP e sostiene sforzi massimali fino a \(\approx 10\,\mathrm{s}\).
\item \textbf{Glicolisi anaerobica}: via anaerobica lattacida. Degradazione del glucosio/glicogeno a piruvato e lattato. Fornisce ATP rapidamente ma si esaurisce in \(10\text{--}30\,\mathrm{s}\) a causa dell'accumulo di \(\mathrm{H^+}\).
\item \textbf{Metabolismo aerobico}: ossidazione completa di carboidrati e lipidi nel mitocondrio. Richiede ossigeno e tempi di attivazione maggiori (\(>60\,\mathrm{s}\)), ma garantisce produzione sostenuta di ATP.
\end{enumerate}
\begin{esempio}[Accoppiamento energetico]
La fosforilazione del glucosio è endoergonica (\(\Delta G^\circ = +13{,}8\,\mathrm{kJ/mol}\)). Accoppiandola all'idrolisi dell'ATP (\(\Delta G^\circ = -30{,}5\,\mathrm{kJ/mol}\)), la reazione complessiva diventa esoergonica (\(\Delta G^\circ = -16{,}7\,\mathrm{kJ/mol}\)), procedendo spontaneamente.
\end{esempio}
\subsection{Recupero della fosfocreatina e ruolo dell'AMP}
Dopo uno sforzo, la PCr viene rigenerata nei mitocondri tramite la \textbf{creatina chinasi mitocondriale}, utilizzando ATP prodotto dal metabolismo aerobico. Il recupero richiede ossigeno: circa il 75\% della PCr si ripristina nel primo minuto di riposo, con recupero quasi completo in \(3\text{--}5\,\mathrm{min}\). L'interruzione del flusso ematico compromette tale recupero, evidenziando l'importanza del riposo attivo tra ripetizioni intense.
Quando l'ADP si accumula rapidamente, l'enzima \textbf{adenilato chinasi} (o miochinasi) catalizza:
\[
\mathrm{2\,ADP \rightleftharpoons ATP + AMP}
\]
Sebbene produca poco ATP, genera \textbf{AMP}, un potente regolatore allosterico che attiva enzimi catabolici (es. glicogeno fosforilasi, PFK-1) per ripristinare l'equilibrio energetico.

\section{Cinetiche delle fonti energetiche}
Le diverse vie metaboliche differiscono radicalmente per \textbf{velocità massima di produzione di ATP} e \textbf{capacità di riserva}:
\begin{itemize}
\item \textbf{PCr}: velocità massima più elevata, ma capacità limitata.
\item \textbf{Glicolisi anaerobica}: velocità \(\approx \frac{1}{2}\) della PCr, capacità moderata.
\item \textbf{Ossidazione aerobica dei carboidrati}: velocità \(\approx \frac{1}{2}\) della glicolisi, capacità elevata.
\item \textbf{Ossidazione aerobica dei lipidi}: velocità \(\approx \frac{1}{2}\) dei carboidrati (e \(\approx \frac{1}{9}\) della PCr), ma capacità praticamente illimitata grazie ai depositi di trigliceridi.
\end{itemize}
Questa differenza cinetica spiega il \textbf{crossover metabolico}: a bassa intensità predominano i lipidi; all'aumentare dell'intensità (\(\% \mathrm{VO_2max}\)), il contributo lipidico diminuisce a favore dei carboidrati, poiché le vie aerobiche non riescono a sostenere la richiesta di ATP ad alte potenze. L'allenamento aerobico sposta il crossover verso intensità maggiori, migliorando l'efficienza ossidativa lipidica.

\section{Metabolismo delle proteine e smaltimento dell'azoto}
Quando carboidrati e lipidi sono esauriti, o in presenza di eccesso proteico, gli \textbf{amminoacidi} vengono catabolizzati a scopo energetico. A differenza di glucidi e lipidi, gli amminoacidi contengono \textbf{azoto} (gruppo amminico). La sua rimozione genera \textbf{ammoniaca (\(\mathrm{NH_3}\))}, molecola altamente tossica per il sistema nervoso. L'organismo la converte rapidamente in \textbf{urea} tramite il \textbf{ciclo dell'urea} (principalmente epatico) per l'eliminazione renale.
In base al destino metabolico dello scheletro carbonioso, gli amminoacidi si classificano in:
\begin{itemize}
\item \textbf{Glucogenici}: convertiti in piruvato o intermedi del ciclo di Krebs, utilizzabili per la gluconeogenesi.
\item \textbf{Chetogenici}: convertiti in acetil-CoA o acetoacetato, precursori dei corpi chetonici.
\end{itemize}

\section{Specificità metaboliche dei tessuti}
I diversi organi specializzano il proprio metabolismo per svolgere funzioni sistemiche:
\begin{itemize}
\item \textbf{Fegato}: centro dell'omeostasi metabolica. Accumula e rilascia glucosio (glicogenosintesi/glicogenolisi), sintetizza glucosio \textit{de novo} (gluconeogenesi), ossida acidi grassi producendo corpi chetonici e VLDL, catabolizza amminoacidi (tranne BCAA) e smaltisce l'ammoniaca (ciclo dell'urea).
\item \textbf{Muscolo scheletrico}: accumula glicogeno e trigliceridi intramuscolari per uso locale. Catabolizza attivamente i \textbf{BCAA} durante l'esercizio per produrre ATP e intermedi energetici. Non rilascia glucosio nel sangue (manca di glucosio-6-fosfatasi).
\item \textbf{Tessuto adiposo}: riserva principale di trigliceridi, rilascio di acidi grassi liberi e glicerolo durante il digiuno o l'esercizio.
\item \textbf{Cervello}: dipendente quasi esclusivamente da glucosio e, in condizioni di digiuno prolungato, da corpi chetonici.
\item \textbf{Cuore}: ossidatore preferenziale di acidi grassi e lattato, altamente efficiente e resistente alla fatica.
\end{itemize}
Questa divisione del lavoro metabolico garantisce un rifornimento energetico coordinato e l'adattamento dinamico alle condizioni fisiologiche e ambientali.
% =========================================================
% Fine Capitolo 10
% =========================================================

% =========================================================
% CAPITOLO 11: Metabolismo degli Zuccheri
% =========================================================
\chapter{Metabolismo degli Zuccheri}
\section{Digestione e assorbimento dei carboidrati}
I carboidrati alimentari (amido, saccarosio) vengono idrolizzati nei loro monomeri (glucosio, fruttosio, galattosio) tramite digestione enzimatica iniziata in bocca e completata nell'intestino tenue. L'assorbimento avviene attraverso la membrana apicale degli enterociti, cellule polarizzate con trasportatori specifici:
\begin{itemize}
\item \textbf{Glucosio e galattosio}: entrano contro gradiente mediante \textbf{trasporto attivo secondario} accoppiato al sodio, mediato da \textbf{SGLT1}.
\item \textbf{Fruttosio}: entra a favore di gradiente per \textbf{diffusione facilitata} tramite \textbf{GLUT5}.
\end{itemize}
Dalla membrana basolaterale, tutti e tre i monosaccaridi escono nel circolo sanguigno tramite \textbf{GLUT2}.
\begin{definizione}[Famiglia GLUT]
I trasportatori GLUT (GLUcose Transporters) differiscono per affinità, distribuzione tissutale e regolazione:
\begin{itemize}
\item \textbf{GLUT1}: alta affinità, ubiquitario, garantisce il trasporto basale.
\item \textbf{GLUT2}: bassa affinità, mai saturo, funge da sensore glicemico nel pancreas.
\item \textbf{GLUT3}: altissima affinità, espresso in cervello e placenta per garantire apporto anche in ipoglicemia.
\item \textbf{GLUT4}: espresso in muscolo e tessuto adiposo, la sua traslocazione in membrana è stimolata dall'insulina.
\item \textbf{GLUT5}: specifico per il fruttosio.
\end{itemize}
\end{definizione}

\section{Fosforilazione del glucosio}
Una volta internalizzato, il glucosio viene immediatamente fosforilato a glucosio-6-fosfato (G6P) consumando una molecola di ATP. La reazione è catalizzata da \textbf{esochinasi} (ubiquitaria) o \textbf{glucochinasi} (epatica).
Questa tappa svolge tre funzioni:
\begin{enumerate}
\item \textbf{Intrappolamento}: il G6P carico non è riconosciuto dai GLUT e rimane nella cellula.
\item \textbf{Attivazione}: aumenta l'energia libera della molecola, rendendola reattiva.
\item \textbf{Regolazione}: è un passo irreversibile e il primo punto di controllo metabolico.
\end{enumerate}

\section{La glicolisi}
La \textbf{glicolisi} è una via citosolica che scinde una molecola di glucosio in due molecole di \textbf{piruvato}, producendo una resa netta di 2 ATP e 2 NADH. Non richiede ossigeno ed è l'unica via in grado di generare ATP in condizioni anaerobiche. Si articola in due fasi:
\subsection{Fase di investimento energetico (reazioni 1-5)}
Consuma 2 ATP per preparare il substrato:
\begin{enumerate}
\item \textbf{Esochinasi/Glucochinasi}: glucosio a G6P.
\item \textbf{Fosfoglucoisomerasi}: G6P (aldoso) a fruttosio-6-fosfato (chetoso).
\item \textbf{Fosfofruttochinasi-1 (PFK-1)}: fruttosio-6-fosfato a fruttosio-1,6-bisfosfato. Consuma il secondo ATP. È il \textbf{principale punto di regolazione} della via.
\item \textbf{Aldolasi}: scinde il fruttosio-1,6-bisfosfato in gliceraldeide-3-fosfato (G3P) e diidrossiacetone fosfato (DHAP).
\item \textbf{Trioso fosfato isomerasi}: converte il DHAP in G3P. A questo punto il glucosio è trasformato in due molecole di G3P.
\end{enumerate}
\subsection{Fase di recupero energetico (reazioni 6-10)}
Produce 4 ATP e 2 NADH dall'ossidazione dei due triosi:
\begin{enumerate}[start=6]
\item \textbf{GAPDH}: ossida il G3P a 1,3-bisfosfoglicerato (1,3-BPG), riducendo il NAD+ a NADH. Il fosfato aggiunto deriva dal fosfato inorganico, non dall'ATP.
\item \textbf{Fosfoglicerato chinasi}: trasferisce il fosfato ad alta energia dall'1,3-BPG all'ADP, producendo ATP e 3-fosfoglicerato. Prima fosforilazione a livello del substrato.
\item \textbf{Fosfoglicerato mutasi}: sposta il fosfato dal C3 al C2, formando 2-fosfoglicerato.
\item \textbf{Enolasi}: disidrata il 2-fosfoglicerato formando fosfoenolpiruvato (PEP), composto ad altissima energia.
\item \textbf{Piruvato chinasi}: trasferisce il fosfato dal PEP all'ADP, producendo ATP e \textbf{piruvato}. Terzo punto di regolazione.
\end{enumerate}
\textbf{Bilancio netto}: 2 ATP + 2 NADH + 2 piruvato per molecola di glucosio.

\section{Destino del piruvato}
Il piruvato segue percorsi diversi in base alla disponibilità di ossigeno:
\begin{itemize}
\item \textbf{Condizioni aerobiche}: entra nel mitocondrio e viene convertito in \textbf{acetil-CoA} dalla piruvato deidrogenasi, producendo NADH e CO2. L'acetil-CoA entra nel ciclo di Krebs.
\item \textbf{Condizioni anaerobiche (muscolo)}: viene ridotto a \textbf{lattato} dalla lattato deidrogenasi (LDH), rigenerando il NAD+ necessario per continuare la glicolisi.
\item \textbf{Fermentazione alcolica (lieviti)}: il piruvato viene convertito in etanolo e CO2.
\end{itemize}
\begin{definizione}[Isoforme LDH]
L'enzima LDH presenta isoforme tissuto-specifiche. Nel muscolo:
\begin{itemize}
\item \textbf{LDH-1 (H-tipo)}: predominante nelle fibre ossidative lente, favorisce la conversione di lattato in piruvato.
\item \textbf{LDH-5 (M-tipo)}: predominante nelle fibre glicolitiche veloci, favorisce la produzione di lattato.
\end{itemize}
\end{definizione}

\section{Metabolismo di altri monosaccaridi}
Altri esosi possono entrare nella glicolisi dopo conversioni specifiche:
\begin{itemize}
\item \textbf{Galattosio}: metabolizzato nel fegato tramite la via di Leloir, che lo converte in glucosio-1-fosfato e poi in G6P.
\item \textbf{Mannosio}: fosforilato a mannosio-6-fosfato e isomerizzato a fruttosio-6-fosfato.
\item \textbf{Fruttosio}: nel muscolo viene fosforilato a fruttosio-6-fosfato dall'esochinasi. Nel \textbf{fegato}, la glucochinasi non lo riconosce; interviene la fruttochinasi che produce fruttosio-1-fosfato, scisso in DHAP e gliceraldeide. Entrambi entrano nella glicolisi \textbf{a valle della PFK-1}, eludendo il principale controllo. Un eccesso cronico di fruttosio epatico può quindi stimolare la lipogenesi e l'accumulo di trigliceridi.
\end{itemize}

\section{Glicogenolisi: degradazione del glicogeno}
Il glicogeno è una riserva altamente ramificata di glucosio. La sua degradazione richiede tre enzimi:
\begin{enumerate}
\item \textbf{Glicogeno fosforilasi}: scinde i legami alfa(1-4) tramite fosforolisi, producendo glucosio-1-fosfato (G1P) senza consumare ATP. Punto di regolazione principale.
\item \textbf{Fosfoglucomutasi}: converte il G1P in glucosio-6-fosfato (G6P).
\item \textbf{Enzima deramificante}: sposta i residui vicino alle ramificazioni e idrolizza il legame alfa(1-6), liberando glucosio libero.
\end{enumerate}
\textbf{Destino del G6P}: nel muscolo entra nella glicolisi per produrre ATP locale. Nel fegato viene defosforilato a glucosio dalla \textbf{glucosio-6-fosfatasi} e rilasciato in circolo per mantenere la glicemia. Il muscolo manca di questo enzima.

\section{Via del pentoso fosfato}
Via citosolica che devia dal G6P e svolge due funzioni principali:
\begin{enumerate}
\item \textbf{Produzione di NADPH}: nella fase ossidativa, il G6P viene ossidato producendo 2 NADPH. Il NADPH fornisce potere riducente per biosintesi (es. acidi grassi) e per i sistemi antiossidanti (rigenerazione del glutatione).
\item \textbf{Produzione di ribosio-5-fosfato}: precursore per la sintesi di nucleotidi e acidi nucleici.
\end{enumerate}
Se il ribosio non è richiesto, la fase non ossidativa riconverte i pentosi in fruttosio-6-fosfato e gliceraldeide-3-fosfato, riconnettendosi alla glicolisi.
\begin{definizione}[Regolazione e deficit di G6PD]
L'enzima chiave \textbf{glucosio-6-fosfato deidrogenasi (G6PD)} è inibito da alti livelli di NADPH. Un suo deficit genetico (favismo) compromette la difesa antiossidante degli eritrociti, scatenando crisi emolitiche in risposta a stress ossidativi (es. ingestione di fave o certi farmaci).
\end{definizione}

\section{Gluconeogenesi: sintesi di glucosio}
Via anabolica che sintetizza glucosio da precursori non glucidici (piruvato, lattato, glicerolo, amminoacidi glucogenici). Avviene principalmente nel fegato e nella corteccia renale, attivata durante il digiuno o l'esercizio prolungato.
Non è il semplice inverso della glicolisi: le tre tappe irreversibili della glicolisi sono bypassate da enzimi specifici:
\begin{enumerate}
\item \textbf{Piruvato $\rightarrow$ Fosfoenolpiruvato}: in due step. Nel mitocondrio, la piruvato carbossilasi converte il piruvato in ossalacetato. Questo viene traslocato nel citosol e convertito in fosfoenolpiruvato dalla PEP carbossichinasi (consuma GTP).
\item \textbf{Fruttosio-1,6-bisfosfato $\rightarrow$ Fruttosio-6-fosfato}: catalizzato dalla fruttosio-1,6-bisfosfatasi, che rimuove il fosfato senza produrre ATP.
\item \textbf{Glucosio-6-fosfato $\rightarrow$ Glucosio}: catalizzato dalla glucosio-6-fosfatasi (reticolo endoplasmatico), assente nel muscolo.
\end{enumerate}
\textbf{Costo energetico}: sintesi di 1 glucosio = consumo di 4 ATP + 2 GTP + 2 NADH. La via è rigidamente regolata per evitare cicli futili con la glicolisi.

\section{Glicogenosintesi: ripristino delle riserve}
Attivata in condizioni di abbondanza di glucosio e alti livelli di insulina. La via utilizza UDP-glucosio come donatore attivo:
\begin{enumerate}
\item \textbf{G6P $\rightarrow$ G1P}: tramite fosfoglucomutasi.
\item \textbf{Attivazione}: G1P + UTP $\rightarrow$ UDP-glucosio + pirofosfato.
\item \textbf{Allungamento}: la \textbf{glicogeno sintasi} trasferisce il glucosio all'estremità non riducente del glicogeno preesistente, formando legami alfa(1-4). Principale punto di regolazione.
\item \textbf{Ramificazione}: l'\textbf{enzima ramificante} crea legami alfa(1-6) spostando oligosaccaridi terminali.
\end{enumerate}
\begin{nota}[Glicogenina]
La glicogeno sintasi richiede un "innesco" di almeno 4 residui. La \textbf{glicogenina} è una proteina autocatalitica che si glicosila formando un primer di 8 residui, permettendo l'intervento della sintasi.
\end{nota}

\section{Integrazione metabolica e regolazione}
Le vie del metabolismo glucidico sono regolate reciprocamente per adattarsi allo stato nutrizionale e alla richiesta energetica:
\begin{itemize}
\item \textbf{Insulina} (post-prandiale): attiva glicogenosintesi e glicolisi; inibisce gluconeogenesi e glicogenolisi.
\item \textbf{Glucagone e adrenalina} (digiuno/esercizio): attivano glicogenolisi e gluconeogenesi; inibiscono glicogenosintesi.
\item \textbf{Regolazione allosterica}: il fruttosio-2,6-bisfosfato è un potente attivatore della PFK-1 e inibitore della fruttosio-1,6-bisfosfatasi, coordinando le due vie in risposta ai segnali ormonali.
\end{itemize}
\begin{esempio}[Ciclo di Cori]
Durante l'esercizio intenso, il muscolo produce lattato. Questo viene rilasciato nel sangue, captato dal fegato e riconvertito in glucosio tramite gluconeogenesi. Il glucosio torna al muscolo, completando un ciclo che ricicla il carbonio e mantiene la glicemia a spese di ATP epatico.
\end{esempio}
% =========================================================
% Fine Capitolo 11
% =========================================================

% =========================================================
% CAPITOLO 12: Ciclo di Krebs e Fosforilazione Ossidativa
% =========================================================
\chapter{Ciclo di Krebs e Fosforilazione Ossidativa}
\section{Introduzione al Ciclo di Krebs}
La \textbf{Fase 3} del catabolismo aerobico di carboidrati, lipidi e amminoacidi converge in un intermedio comune: l'\textbf{acetil-CoA} (2 atomi di carbonio). Nella Fase 3a, l'acetil-CoA entra nel \textbf{ciclo di Krebs} (noto anche come ciclo dell'acido citrico o ciclo degli acidi tricarbossilici), una via metabolica localizzata nella \textbf{matrice mitocondriale}.
Sebbene il ciclo non utilizzi direttamente ossigeno, esso si arresta in condizioni di ipossia per esaurimento dei cofattori ossidati (\(\mathrm{NAD^+}\) e \(\mathrm{FAD}\)). Il ciclo non produce ATP direttamente, ma genera coenzimi ridotti (\(\mathrm{NADH}\), \(\mathrm{FADH_2}\)) e GTP, che verranno convertiti in ATP tramite la fosforilazione ossidativa.
\subsection{Struttura e bilancio generale}
Il ciclo è composto da \textbf{8 tappe enzimatiche} (di cui 3 irreversibili) e si articola in due fasi logiche:
\begin{enumerate}
\item \textbf{Fase I (Ossidazione dell'acetile)}: i due atomi di carbonio dell'acetil-CoA si uniscono all'ossalacetato e vengono completamente ossidati a \(\mathrm{CO_2}\).
\item \textbf{Fase II (Rigenerazione dell'ossalacetato)}: lo scheletro carbonioso residuo viene riorganizzato e ossidato per riformare l'ossalacetato, pronto per un nuovo ciclo.
\end{enumerate}
Per ogni molecola di acetil-CoA processata, il bilancio netto è:
\[
\text{Acetil-CoA} + 3\,\mathrm{NAD^+} + \mathrm{FAD} + \mathrm{GDP} + \mathrm{P_i} + 2\,\mathrm{H_2O} \rightarrow 2\,\mathrm{CO_2} + 3\,\mathrm{NADH} + \mathrm{FADH_2} + \mathrm{GTP} + \text{CoA-SH} + 2\,\mathrm{H^+}
\]
\subsection{Le otto tappe del ciclo}
\begin{enumerate}
\item \textbf{Formazione del citrato}: l'acetil-CoA (2C) si condensa con l'ossalacetato (4C) formando citrato (6C). Catalizzata dalla \textbf{citrato sintasi}. Reazione irreversibile.
\item \textbf{Isomerizzazione del citrato}: il citrato viene convertito in isocitrato dall'\textbf{aconitasi}. L'isocitrato viene immediatamente rimosso dalla tappa successiva, spostando l'equilibrio.
\item \textbf{Decarbossilazione ossidativa dell'isocitrato}: l'isocitrato (6C) è ossidato ad \(\alpha\)-chetoglutarato (5C) dall'\textbf{isocitrato deidrogenasi}. Rilascia la prima molecola di \(\mathrm{CO_2}\) e produce il primo \(\mathrm{NADH}\). Reazione irreversibile.
\item \textbf{Decarbossilazione ossidativa dell'\(\alpha\)-chetoglutarato}: l'\(\alpha\)-chetoglutarato (5C) è convertito in succinil-CoA (4C) dal complesso della \(\alpha\)-\textbf{chetoglutarato deidrogenasi}. Rilascia la seconda molecola di \(\mathrm{CO_2}\) e produce il secondo \(\mathrm{NADH}\). Reazione irreversibile. Termine della Fase I.
\item \textbf{Conversione del succinil-CoA in succinato}: la scissione del legame tioestere del succinil-CoA rilascia energia sufficiente per sintetizzare \textbf{GTP} dal GDP (\textbf{fosforilazione a livello del substrato}). Il GTP è energeticamente equivalente all'ATP.
\item \textbf{Ossidazione del succinato}: il succinato è ossidato a fumarato dalla \textbf{succinato deidrogenasi}, un enzima legato alla membrana mitocondriale interna. Gli elettroni rimossi riducono il \(\mathrm{FAD}\) a \(\mathrm{FADH_2}\).
\item \textbf{Idratazione del fumarato}: la \textbf{fumarasi} aggiunge una molecola d'acqua al doppio legame del fumarato, formando malato. L'enzima è specifico per l'isomero trans.
\item \textbf{Rigenerazione dell'ossalacetato}: il malato è ossidato ad ossalacetato dalla \textbf{malato deidrogenasi}, producendo il terzo \(\mathrm{NADH}\) e chiudendo il ciclo.
\end{enumerate}
\begin{definizione}[Natura anfibolica]
Il ciclo di Krebs è una via \textbf{anfibolica}: oltre a degradare l'acetil-CoA (catabolismo), fornisce intermedi per vie biosintetiche (anabolismo). Le reazioni che sottraggono intermedi sono dette \textit{cataplerotiche}, mentre quelle che li riforniscono sono dette \textit{anaplerotiche}.
\end{definizione}
\subsection{Regolazione del ciclo}
Le tre reazioni irreversibili costituiscono i principali punti di controllo. Il ciclo è finemente modulato dallo stato energetico cellulare:
\begin{itemize}
\item \textbf{Inibizione}: alte concentrazioni di \(\mathrm{ATP}\), \(\mathrm{NADH}\) e succinil-CoA segnalano un sufficiente stato energetico, inibendo citrato sintasi, isocitrato deidrogenasi e \(\alpha\)-chetoglutarato deidrogenasi.
\item \textbf{Attivazione muscolare}: il \(\mathrm{Ca^{2+}}\), rilasciato per innescare la contrazione, stimola simultaneamente gli enzimi del ciclo per rigenerare rapidamente l'ATP consumato.
\end{itemize}

\section{Fosforilazione Ossidativa e Catena di Trasporto degli Elettroni}
I coenzimi ridotti (\(\mathrm{NADH}\), \(\mathrm{FADH_2}\)) prodotti nel ciclo di Krebs e nelle fasi precedenti del catabolismo convogliano i loro elettroni nella \textbf{catena di trasporto degli elettroni (ETC)}, localizzata nella \textbf{membrana mitocondriale interna}. Il flusso di elettroni attraverso quattro complessi proteici genera un gradiente elettrochimico di protoni, la cui energia viene sfruttata dall'\textbf{ATP sintetasi} per fosforilare l'ADP.
\subsection{Navette per il NADH citosolico}
La glicolisi produce \(\mathrm{NADH}\) nel citosol, ma la membrana mitocondriale è impermeabile al \(\mathrm{NADH}\). Il trasferimento degli equivalenti riducenti avviene tramite sistemi navetta (\textit{shuttle}):
\begin{itemize}
\item \textbf{Navetta malato-aspartato} (fegato, cuore): il \(\mathrm{NADH}\) citosolico riduce l'ossalacetato a malato, che entra nel mitocondrio e viene riossidato, \textbf{rigenerando NADH nella matrice}. Resa massima.
\item \textbf{Navetta del glicerolo-3-fosfato} (muscolo scheletrico, cervello): gli elettroni sono trasferiti al diidrossiacetone fosfato (DHAP), formando glicerolo-3-fosfato. Un enzima mitocondriale lo riossida, riducendo il \(\mathrm{FAD}\) a \(\mathrm{FADH_2}\). Resa leggermente inferiore, ma rende il muscolo meno dipendente dalla disponibilità di \(\mathrm{NAD^+}\).
\end{itemize}
\subsection{I complessi della catena di trasporto}
\begin{enumerate}
\item \textbf{Complesso I} (NADH:ubichinone ossidoreduttasi): accetta elettroni dal \(\mathrm{NADH}\), li trasferisce all'ubichinone (\(\mathrm{Q}\)) e pompa \textbf{4 protoni} nello spazio intermembrana.
\item \textbf{Complesso II} (succinato deidrogenasi): identico all'enzima del ciclo di Krebs. Trasferisce elettroni dal \(\mathrm{FADH_2}\) all'ubichinone. \textbf{Non pompa protoni}.
\item \textbf{Complesso III} (ubichinone-citocromo c ossidoreduttasi): trasferisce elettroni dall'ubichinone ridotto al citocromo c (proteina mobile nello spazio intermembrana). Pompa \textbf{4 protoni}.
\item \textbf{Complesso IV} (citocromo c ossidasi): accetta elettroni dal citocromo c e li trasferisce all'ossigeno molecolare (\(\mathrm{O_2}\)), riducendolo ad \(\mathrm{H_2O}\). Pompa \textbf{2 protoni}. Il trasferimento simultaneo di 4 elettroni previene la formazione di specie reattive dell'ossigeno (ROS).
\end{enumerate}
\begin{definizione}[Forza motrice protonica]
La membrana interna è impermeabile ai protoni (ricca di cardiolipina). Il pompaggio attivo crea un gradiente di pH (\(\Delta \mathrm{pH}\)) e un potenziale elettrico (\(\Delta \Psi\)) tra matrice e spazio intermembrana. Questa energia potenziale combinata è la \textbf{forza motrice protonica} (\(\Delta G\)), che guida il rientro dei protoni nella matrice.
\end{definizione}
\subsection{ATP Sintasi e rapporto P/O}
L'\textbf{ATP sintetasi} (Complesso V) è l'unico canale che permette il rientro dei protoni. È un motore molecolare reversibile composto da:
\begin{itemize}
\item \textbf{Porzione \(F_o\)}: canale transmembrana attraverso cui i protoni fluiscono secondo gradiente.
\item \textbf{Porzione \(F_1\)}: macchina catalitica citosolica/matrix. Il flusso protonico induce una \textbf{rotazione meccanica} che modifica la conformazione dei siti attivi, catalizzando la sintesi di ATP da ADP e \(\mathrm{P_i}\).
\end{itemize}
L'efficienza di accoppiamento è misurata dal \textbf{rapporto P/O} (ATP prodotti per atomo di ossigeno ridotto). Nei mitocondri di mammifero i valori accettati sono:
\begin{itemize}
\item \(\approx \mathbf{2{,}5}\) ATP per molecola di \(\mathrm{NADH}\)
\item \(\approx \mathbf{1{,}5}\) ATP per molecola di \(\mathrm{FADH_2}\)
\end{itemize}

\section{Resa energetica totale del glucosio}
Sommando i prodotti della glicolisi, dell'ossidazione del piruvato e del ciclo di Krebs, l'ossidazione completa di una molecola di glucosio produce:
\begin{table}[htbp]
\centering
\caption{Bilancio energetico dell'ossidazione completa del glucosio}
\label{tab:resa-glucosio}
\renewcommand{\arraystretch}{1.2}
\begin{tabular}{@{} p{0.4\textwidth} p{0.25\textwidth} p{0.25\textwidth} @{}}
\toprule
\textbf{Processo} & \textbf{Prodotto} & \textbf{ATP equivalenti} \\
\midrule
Glicolisi & 2 NADH + 2 ATP & \(5 + 2 = 7\) \\
Ossidazione del Piruvato & 2 NADH & \(5\) \\
Ciclo di Krebs & 6 NADH + 2 FADH\(_2\) + 2 GTP & \(15 + 3 + 2 = 20\) \\
\midrule
\textbf{Totale per molecola di glucosio} & & \(\approx \mathbf{32}\) ATP \\
\bottomrule
\end{tabular}
\end{table}
\begin{nota}[Efficienza comparata]
La resa aerobica (\(\approx 32\) ATP) è circa 16 volte superiore a quella della glicolisi anaerobica (2 ATP). La differenza risiede nell'ossidazione completa del carbonio a \(\mathrm{CO_2}\) e nell'efficiente recupero energetico tramite la catena respiratoria e il gradiente protonico.
\end{nota}
% =========================================================
% Fine Capitolo 12
% =========================================================

% =========================================================
% CAPITOLO 13: Metabolismo dei Lipidi
% =========================================================
\chapter{Metabolismo dei Lipidi}
\section{Digestione, Assorbimento e Mobilizzazione}
Il catabolismo lipidico a fini energetici utilizza principalmente acidi grassi derivanti dalla dieta, dai trigliceridi (TAG) di riserva o dalla \textit{de novo} sintesi. La digestione dei lipidi alimentari avviene nell'intestino tenue, dove la bile (prodotta dal fegato e immagazzinata nella cistifellea) emulsiona le gocce lipidiche, facilitando l'azione delle lipasi. I prodotti finali (acidi grassi, diacilgliceroli, lisofosfolipidi) formano micelle e vengono assorbiti per diffusione semplice negli enterociti.
\begin{itemize}
\item \textbf{Acidi grassi a catena lunga}: all'interno dell'enterocita vengono risintetizzati in TAG, assemblati con colesterolo e apolipoproteine nei \textbf{chilomicroni}, e rilasciati nel sistema linfatico/sanguigno. A livello dei capillari tissutali, la lipoproteina lipasi idrolizza i TAG rilasciando acidi grassi e glicerolo per l'uso o lo stoccaggio locale.
\item \textbf{Acidi grassi a catena corta e media}: passano direttamente nel circolo sanguigno legati all'albumina e raggiungono il fegato.
\end{itemize}
\subsection{Mobilizzazione dei TAG di riserva}
In risposta a bassi livelli di glucosio, ormoni come glucagone e adrenalina stimolano la lipolisi negli adipociti tramite la via del cAMP. Tre enzimi agiscono in sequenza:
\begin{enumerate}
\item \textbf{ATGL} (Adipose Triglyceride Lipase): stacca il primo acido grasso, formando diacilglicerolo (DAG).
\item \textbf{HSL} (Hormone-Sensitive Lipase): stacca il secondo acido grasso, formando monoacilglicerolo (MAG).
\item \textbf{MGL} (Monoacylglycerol Lipase): stacca l'ultimo acido grasso, liberando glicerolo.
\end{enumerate}
Gli acidi grassi rilasciati viaggiano nel sangue legati all'albumina e vengono captati dalle cellule target (es. muscolo) tramite trasportatori specifici (FABP, FATP, CD36), la cui espressione è aumentata dall'allenamento di endurance.

\section{Ossidazione degli Acidi Grassi (\(\beta\)-ossidazione)}
L'ossidazione degli acidi grassi avviene nella matrice mitocondriale e procede attraverso un ciclo ripetuto di quattro reazioni che rimuove unità bicarboniose sotto forma di acetil-CoA.
\subsection{Attivazione e trasporto mitocondriale}
Gli acidi grassi a catena lunga devono essere prima attivati nel citosol dalla \textbf{acil-CoA sintetasi}, consumando ATP:
\[
\mathrm{Acido\ grasso + ATP + CoA\text{-}SH \rightarrow Acil\text{-}CoA + AMP + PP_i}
\]
L'acil-CoA non può attraversare direttamente la membrana mitocondriale interna e richiede lo \textbf{shuttle della carnitina}:
\begin{enumerate}
\item \textbf{CPT1} (Carnitina Aciltransferasi I): sulla membrana esterna, trasferisce l'acile alla carnitina formando acil-carnitina. \textit{Tappa limitante e principale punto di regolazione}.
\item \textbf{Traslocasi}: scambia acil-carnitina con carnitina libera attraverso la membrana interna.
\item \textbf{CPT2} (nella matrice): trasferisce l'acile dalla carnitina al CoA mitocondriale, rigenerando acil-CoA.
\end{enumerate}
Il glicerolo rilasciato dalla lipolisi viene invece ossidato a diidrossiacetone fosfato (DHAP) e può entrare nella glicolisi o nella gluconeogenesi (principalmente in fegato e rene).
\subsection{Il ciclo della \(\beta\)-ossidazione}
Ogni ciclo comprende quattro passaggi enzimatici:
\begin{enumerate}
\item \textbf{Deidrogenazione}: l'acil-CoA deidrogenasi introduce un doppio legame \(\mathrm{trans}\) tra C2 e C3, producendo \(\mathrm{FADH_2}\).
\item \textbf{Idratazione}: l'enoi-CoA idratasi aggiunge \(\mathrm{H_2O}\) al doppio legame, formando \(\beta\)-idrossiacil-CoA.
\item \textbf{Deidrogenazione}: la \(\beta\)-idrossiacil-CoA deidrogenasi ossida il gruppo \(-\mathrm{OH}\) a chetone, producendo \(\mathrm{NADH}\).
\item \textbf{Tiolisi}: la tiolasi scinde la molecola con un CoA, rilasciando \textbf{acetil-CoA} e un acil-CoA accorciato di due atomi di carbonio.
\end{enumerate}
Il ciclo si ripete fino alla completa conversione dell'acido grasso in acetil-CoA.
\begin{definizione}[Resa energetica]
Il palmitato (\(16:0\)) richiede 7 cicli di \(\beta\)-ossidazione, producendo 8 acetil-CoA, 7 \(\mathrm{NADH}\) e 7 \(\mathrm{FADH_2}\). Considerando i rapporti P/O (\(\mathrm{NADH} \approx 2{,}5\) ATP, \(\mathrm{FADH_2} \approx 1{,}5\) ATP) e sottraendo i 2 equivalenti ATP per l'attivazione, la resa netta è di \(\mathbf{106}\) ATP per molecola di palmitato, significativamente superiore ai 32 ATP del glucosio.
\end{definizione}
\subsection{Ossidazione di acidi grassi insaturi e a catena dispari}
\begin{itemize}
\item \textbf{Insaturi}: richiedono enzimi ausiliari (\textit{enoi-CoA isomerasi} per convertire legami \textit{cis} in \textit{trans}, e \textit{2,4-dienoi-CoA reduttasi} per acidi poliinsaturi) per permettere la continuazione della \(\beta\)-ossidazione.
\item \textbf{A catena dispari}: l'ultimo ciclo produce \textbf{propionil-CoA} (3C), convertito in succinil-CoA (intermedio del ciclo di Krebs) tramite biotina e vitamina B12.
\end{itemize}

\section{Regolazione del Catabolismo Lipidico}
Il flusso degli acidi grassi verso il mitocondrio è finemente regolato per evitare cicli futili con la sintesi:
\begin{itemize}
\item \textbf{Malonil-CoA}: primo intermedio della sintesi degli acidi grassi, inibisce potentemente la \textbf{CPT1}, bloccando l'ingresso degli acidi grassi nel mitocondrio e dirottandoli verso la sintesi di TAG/fosfolipidi nel citosol.
\item \textbf{Feedback energetico}: alti rapporti \(\mathrm{NADH}/\mathrm{NAD^+}\) e alte concentrazioni di acetil-CoA inibiscono rispettivamente la \(\beta\)-idrossiacil-CoA deidrogenasi e la tiolasi, segnalando abbondanza energetica.
\end{itemize}

\section{Biosintesi degli Acidi Grassi e dei Trigliceridi}
La lipogenesi avviene nel citosol (fegato e tessuto adiposo) e utilizza acetil-CoA (esportato dal mitocondrio tramite la navetta del citrato) e NADPH (dalla via del pentoso fosfato).
\subsection{Tappe della sintesi}
\begin{enumerate}
\item \textbf{Formazione del malonil-CoA}: l'\textbf{acetil-CoA carbossilasi} (ACC) carbossila l'acetil-CoA consumando ATP e biotina. È la \textbf{tappa limitante e regolata}.
\item \textbf{Complesso acido grasso sintasi}: catalizza cicli ripetuti di condensazione, riduzione, disidratazione e riduzione, aggiungendo 2 atomi di carbonio (dal malonil-CoA) per ciclo fino al \textbf{palmitato} (16C).
\end{enumerate}
Il palmitato può essere successivamente allungato (nel RE o mitocondri) e desaturato (da desaturasi che inseriscono doppi legami \textit{cis}). Poiché i mammiferi non possono inserire doppi legami oltre il C9, gli acidi grassi \(\omega\)-3 e \(\omega\)-6 sono \textbf{essenziali}.
\subsection{Sintesi di TAG e fosfolipidi}
Entrambi derivano dall'\textbf{acido fosfatidico}, formato dall'esterificazione di due acidi grassi al glicerolo-3-fosfato. L'aggiunta di una terza testa polare o di un terzo acido grasso determina la formazione di fosfolipidi o trigliceridi, rispettivamente.

\section{Sintesi del Colesterolo}
Il colesterolo è sintetizzato nel citosol/RE a partire da acetil-CoA. Le tappe chiave sono:
\begin{enumerate}
\item Condensazione di 2 acetil-CoA in acetoacetil-CoA, poi aggiunta di un terzo acetil-CoA per formare \textbf{HMG-CoA}.
\item Riduzione dell'HMG-CoA a \textbf{mevalonato} da parte della \textbf{HMG-CoA reduttasi} (consuma 2 NADPH). Questa è la \textbf{tappa limitante}.
\item Fosforilazione e decarbossilazione del mevalonato per formare unità isopreniche attivate, che si condensano in \textbf{squalene} e ciclizzano a colesterolo.
\end{enumerate}
\begin{definizione}[Regolazione della sintesi]
La HMG-CoA reduttasi è regolata a:
\begin{itemize}
\item \textit{Breve termine}: fosforilazione (inattivazione) da AMPK quando l'ATP è basso.
\item \textit{Lungo termine}: degradazione proteasomica in presenza di alto colesterolo; regolazione trascrizionale mediata da SREBP in condizioni di carenza.
\end{itemize}
\end{definizione}

\section{Produzione e Utilizzo dei Corpi Chetonici}
In condizioni di digiuno prolungato o diabete, il fegato devia l'acetil-CoA (che non può essere completamente ossidato per carenza di ossalacetato, dirottato nella gluconeogenesi) verso la \textbf{chetogenesi}.
\subsection{Sintesi epatica}
Avviene nei mitocondri epatici. Due acetil-CoA formano acetoacetil-CoA, che condensa con un terzo acetil-CoA a formare HMG-CoA. La \textbf{HMG-CoA liasi} scinde quest'ultimo in \textbf{acetoacetato} e acetil-CoA. L'acetoacetato può essere ridotto a \(\beta\)-\textbf{idrossibutirrato} o decarbossilato ad \textbf{acetone}.
\subsection{Utilizzo extraepatico}
I corpi chetonici sono esportati nel sangue e captati da cervello, muscolo e cuore. La \(\beta\)-\textbf{chetoacil-CoA trasferasi} (tioforasi) attiva l'acetoacetato usando succinil-CoA, formando acetoacetil-CoA, poi scisso in 2 acetil-CoA per il ciclo di Krebs.
\begin{nota}[Specificità tissutale]
Il fegato \textbf{non esprime la tioforasi}, quindi produce corpi chetonici ma non può utilizzarli, garantendo un rifornimento energetico costante per i tessuti periferici in condizioni di carenza glucidica.
\end{nota}
% =========================================================
% Fine Capitolo 13
% =========================================================

% =========================================================
% CAPITOLO 14: Metabolismo degli Amminoacidi
% =========================================================
\chapter{Metabolismo degli Amminoacidi}
\section{Digestione, assorbimento e pool degli amminoacidi}
Le proteine alimentari vengono idrolizzate nel tratto gastrointestinale ad opera di proteasi (pepsina nello stomaco, tripsina, chimotripsina ed elastasi pancreatiche nell'intestino) e peptidasi di membrana enterocitaria, fino a liberare amminoacidi singoli.
L'assorbimento avviene sulla membrana apicale degli enterociti tramite \textbf{trasporto attivo secondario} accoppiato al sodio (\(\mathrm{Na^+}\)), che mantiene un gradiente elettrochimico grazie all'azione basolaterale della \(\mathrm{Na^+/K^+}\)-ATPasi.
Gli amminoacidi liberati confluiscono nel \textbf{pool degli amminoacidi liberi}, una riserva intracellulare molto esigua (circa l'1\% del totale proteico), che non esiste in forme di deposito specializzate come glicogeno o trigliceridi. Questo pool viene utilizzato principalmente per la sintesi proteica. Solo in condizioni di eccesso dietetico, digiuno prolungato o sforzo massimale protratto, gli amminoacidi vengono catabolizzati a scopo energetico, fornendo circa il 10--15\% del fabbisogno giornaliero.

\section{Catabolismo e rimozione del gruppo amminico}
A differenza di carboidrati e lipidi, gli amminoacidi contengono azoto sotto forma di gruppo amminico (\(-\mathrm{NH_2}\)). Prima che lo scheletro carbonioso possa essere ossidato per produrre energia, il gruppo amminico deve essere rimosso, generando \textbf{ammoniaca} (\(\mathrm{NH_3}\)), molecola altamente tossica che viene protonata a ione ammonio (\(\mathrm{NH_4^+}\)) a pH fisiologico.
\subsection{Transaminazione}
Il gruppo amminico viene trasferito reversibilmente a un \(\alpha\)-chetoglutarato, formando \textbf{glutammato} e l'\(\alpha\)-chetoacido corrispondente all'amminoacido di partenza. La reazione è catalizzata da \textbf{aminotrasferasi} (o transaminasi) e richiede il cofattore \textbf{piridossal fosfato (PLP)}, derivato della vitamina B6.
\[
\text{Amminoacido}_1 + \alpha\text{-chetoglutarato} \rightleftharpoons \alpha\text{-chetoacido}_1 + \text{Glutammato}
\]
Le transaminasi più rilevanti clinicamente sono l'\textbf{ALT} (alanina aminotransferasi) e l'\textbf{AST} (aspartato aminotransferasi). Elevati livelli ematici di questi enzimi sono marcatori di danno epatico o muscolare.
\subsection{Deaminazione ossidativa}
Il glutammato è l'unico amminoacido che può subire deaminazione diretta nel mitocondrio. La \textbf{glutammato deidrogenasi} catalizza la rimozione del gruppo amminico, rigenerando \(\alpha\)-chetoglutarato e liberando \(\mathrm{NH_4^+}\), con produzione simultanea di \(\mathrm{NAD(P)H}\).
\[
\text{Glutammato} + \mathrm{NAD^+} + \mathrm{H_2O} \rightleftharpoons \alpha\text{-chetoglutarato} + \mathrm{NH_4^+} + \mathrm{NADH} + \mathrm{H^+}
\]

\section{Il ciclo dell'urea}
Per smaltire l'azoto tossico, i mammiferi convertono l'\(\mathrm{NH_4^+}\) in \textbf{urea} tramite il \textbf{ciclo dell'urea}, una via metabolica presente nella sua forma completa esclusivamente nel \textbf{fegato}.
\subsection{Trasporto dell'azoto al fegato}
I tessuti extraepatici non possono eseguire il ciclo dell'urea e devono trasportare l'\(\mathrm{NH_4^+}\) al fegato:
\begin{itemize}
\item \textbf{Glutammina}: formata in molti tessuti dalla glutammina sintetasi (\(\text{Glutammato} + \mathrm{NH_4^+} + \mathrm{ATP} \rightarrow \text{Glutammina}\)).
\item \textbf{Ciclo glucosio-alanina} (muscolo): il gruppo amminico del glutammato viene trasferito al piruvato dall'ALT, formando \textbf{alanina}. L'alanina raggiunge il fegato, dove viene riconvertita in piruvato (utilizzato per gluconeogenesi) e glutammato, dal quale viene liberato \(\mathrm{NH_4^+}\) per il ciclo.
\end{itemize}
\subsection{Tappe del ciclo}
Il ciclo incorpora due atomi di azoto (uno dall'\(\mathrm{NH_4^+}\), uno dall'aspartato) e consuma l'equivalente di \textbf{4 legami ad alta energia} (3 ATP, di cui uno idrolizzato ad AMP + \(\mathrm{PP_i}\)):
\begin{enumerate}
\item \textbf{Sintesi del carbamil fosfato} (matrice mitocondriale): \(\mathrm{NH_4^+} + \mathrm{HCO_3^-} + 2\mathrm{ATP} \xrightarrow{\text{CPS I}} \text{Carbamil fosfato}\). Tappa limitante, attivata allostericamente dall'\textbf{N-acetilglutammato} (segno di eccesso di glutammato/amminoacidi).
\item \textbf{Formazione della citrullina} (matrice): \(\text{Carbamil fosfato} + \text{Ornitina} \xrightarrow{\text{OTC}} \text{Citrullina}\). La citrullina esce nel citosol.
\item \textbf{Formazione dell'argininosuccinato} (citosol): \(\text{Citrullina} + \text{Aspartato} + \mathrm{ATP} \xrightarrow{\text{Argininosuccinato sintetasi}} \text{Argininosuccinato}\).
\item \textbf{Scissione in arginina e fumarato}: \(\text{Argininosuccinato} \xrightarrow{\text{Argininosuccinasi}} \text{Arginina} + \text{Fumarato}\). Il fumarato entra nel ciclo di Krebs, collegando le due vie.
\item \textbf{Idrolisi dell'arginina}: \(\text{Arginina} + \mathrm{H_2O} \xrightarrow{\text{Arginasi}} \text{Urea} + \text{Ornitina}\). L'ornitina rientra nel mitocondrio per ricominciare il ciclo.
\end{enumerate}
\begin{nota}[Complicanze metaboliche]
Deficit ereditari di qualsiasi enzima del ciclo causano \textbf{iperammonemia}, con accumulo di \(\mathrm{NH_4^+}\) nel sangue e gravi danni neurologici (coma, ritardo cognitivo) per la sua capacità di attraversare la barriera ematoencefalica.
\end{nota}

\section{Destino metabolico degli scheletri carboniosi}
Una volta rimosso l'azoto, lo scheletro carbonioso degli amminoacidi viene convertito in intermedi metabolici chiave, classificabili in base al loro destino finale:
\begin{itemize}
\item \textbf{Glucogenici}: convertiti in piruvato o intermedi del ciclo di Krebs (ossalacetato, \(\alpha\)-chetoglutarato, succinil-CoA, fumarato), utilizzabili per la gluconeogenesi epatica.
\item \textbf{Chetogenici}: convertiti in acetil-CoA o acetoacetato, precursori dei corpi chetonici e degli acidi grassi.
\item \textbf{Entrambi}: alcuni amminoacidi (es. fenilalanina, tirosina, triptofano, isoleucina, treonina) danno origine a intermedi sia glucogenici che chetogenici. Solo leucina e lisina sono esclusivamente chetogenici.
\end{itemize}
\subsection{Metabolismo dei BCAA}
Gli amminoacidi a catena ramificata (\textbf{BCAA}: valina, leucina, isoleucina) sono essenziali e rappresentano il 35\% degli amminoacidi presenti nelle proteine muscolari. A differenza della maggior parte degli amminoacidi, \textbf{non sono catabolizzati nel fegato} (che manca dell'enzima BCAT), ma vengono transaminati e ossidati \textbf{direttamente nel muscolo scheletrico}.
Durante l'esercizio prolungato o il digiuno, il muscolo degrada le proprie proteine, ossida i BCAA per produrre ATP e utilizza i gruppi amminici per sintetizzare alanina e glutammina, che vengono rilasciate in circolo per il fegato.

\section{Biosintesi e ruolo degli amminoacidi come precursori}
\subsection{Essenziali e non essenziali}
L'organismo umano può sintetizzare \textbf{amminoacidi non essenziali} a partire da intermedi della glicolisi, del ciclo di Krebs o della via del pentoso fosfato. Gli \textbf{amminoacidi essenziali} (istidina, isoleucina, leucina, lisina, metionina, fenilalanina, treonina, triptofano, valina) devono essere introdotti con la dieta. In condizioni di stress o sviluppo, alcuni non essenziali diventano \textbf{condizionalmente essenziali}.
\subsection{Precursori di composti azotati}
Oltre alla sintesi proteica, gli amminoacidi sono precursori di molecole biologiche fondamentali:
\begin{itemize}
\item \textbf{Glicina}: precursore dell'\textbf{eme} (gruppo prostetico di emoglobina e citocromi) e della \textbf{glutatione} (tripeptide antiossidante GSH, insieme a cisteina e glutammato).
\item \textbf{Arginina}: precursore dell'\textbf{ossido nitrico (NO)}, potente vasodilatatore sintetizzato dalla NO sintasi (NOS).
\item \textbf{Triptofano}: precursore della \textbf{serotonina} (neurotrasmettitore) e della \textbf{melatonina} (regolatore del ciclo sonno-veglia).
\item \textbf{Tirosina}: precursore delle \textbf{catecolamine} (dopamina, noradrenalina, adrenalina) e degli ormoni tiroidei.
\end{itemize}
% =========================================================
% Fine Capitolo 14
% =========================================================

% =========================================================
% CAPITOLO 15: Principi di regolazione del metabolismo
% =========================================================
\chapter{Principi di regolazione del metabolismo}
\section{Introduzione e omeostasi metabolica}
I processi metabolici cellulari devono essere costantemente regolati per mantenere l'equilibrio dinamico (\textbf{omeostasi}) tra la produzione di energia (catabolismo) e la sintesi di macromolecole (anabolismo). Tale regolazione è essenziale per permettere le normali funzioni fisiologiche e gli adattamenti a stimoli esterni come variazioni termiche, disponibilità di nutrienti o attività fisica. Sbilanciamenti protratti in questi meccanismi possono portare all'insorgenza di disordini metabolici sistemici.

\section{Ormoni: classificazione e meccanismi d'azione}
La regolazione metabolica avviene a più livelli: dalla modulazione fine dell'attività enzimatica tramite effettori allosterici, alla regolazione sistemica mediata da \textbf{ormoni}. Gli ormoni sono molecole segnale secrete da ghiandole endocrine, trasportate tramite il circolo sanguigno fino a raggiungere cellule bersaglio dotate di recettori specifici.
In base alla natura chimica, si distinguono due classi principali:
\begin{itemize}
\item \textbf{Ormoni peptidici o derivati da amminoacidi} (es. insulina, adrenalina, glucagone, GH): molecole idrofiliche che \textbf{non} attraversano la membrana plasmatica. Agiscono legandosi a recettori di superficie, attivando cascate di segnalazione intracellulare che modificano rapidamente l'attività di enzimi esistenti.
\item \textbf{Ormoni steroidei} (es. cortisolo, testosterone, estrogeni): derivati dal colesterolo, sono lipofilici e \textbf{attraversano} facilmente la membrana. Si legano a recettori citoplasmatici o nucleari, modulando la \textbf{trascrizione genica} e quindi la sintesi \textit{de novo} di enzimi metabolici (effetto più lento ma duraturo).
\end{itemize}

\section{Regolazione metabolica durante l'esercizio fisico}
Durante un'attività fisica sostenuta (es. 70\% $\mathrm{VO_2max}$), si osserva una netta modificazione del profilo ormonale:
\begin{itemize}
\item \textbf{Diminuisce l'insulina} (ormone anabolico).
\item \textbf{Aumentano adrenalina e glucagone} (ormoni catabolici).
\end{itemize}
Questa modificazione stimola la mobilizzazione dei substrati energetici: aumento della glicogenolisi e glicolisi (con produzione di lattato) e aumento della lipolisi (rilascio di glicerolo e acidi grassi liberi). L'obiettivo finale è il \textbf{mantenimento della glicemia} e il rifornimento continuo di ATP ai muscoli attivi. Al termine dell'esercizio, il profilo ormonale si inverte per favorire il recupero (risintesi di glicogeno e incorporazione dei nutrienti).

\section{Cascate di segnalazione ormonale}
\subsection{Adrenalina e glucagone: attivazione del catabolismo}
L'adrenalina (secretata dalla midollare del surrene) e il glucagone (dal pancreas) agiscono tramite recettori accoppiati a \textbf{proteine G}. Il legame ormone-recettore attiva la proteina G, che stimola l'\textbf{adenilato ciclasi} a convertire ATP in \textbf{AMP ciclico (cAMP)}. Il cAMP attiva a sua volta la \textbf{Proteina Chinasi A (PKA)}, che fosforila enzimi bersaglio:
\begin{itemize}
\item \textbf{Glicogenolisi}: la PKA fosforila e attiva la glicogeno fosforilasi chinasi, che a sua volta attiva la \textbf{glicogeno fosforilasi}, enzima chiave della degradazione del glicogeno.
\item \textbf{Lipolisi}: la PKA fosforila la \textbf{lipasi ormone-sensibile (HSL)}, promuovendo il distacco di acidi grassi dal diacilglicerolo.
\end{itemize}
Al termine dello stimolo, l'azione delle fosfatasi defosforila questi enzimi, spegnendo la risposta.
\subsection{Insulina e IGF-1: attivazione dell'anabolismo}
L'insulina si lega a un recettore con attività \textbf{tirosina-chinasica} intrinseca. L'attivazione innesca una cascata che coinvolge IRS e la chinasi \textbf{AKT}:
\begin{itemize}
\item \textbf{Glicogenosintesi}: AKT fosforila e inattiva la \textbf{GSK3} (glicogeno sintasi chinasi 3). Poiché GSK3 inattiva normalmente la glicogeno sintasi, la sua inibizione da parte di AKT \textbf{promuove la sintesi di glicogeno}.
\item \textbf{Inibizione della lipolisi}: AKT attiva la \textbf{fosfodiesterasi 3 (PDE3)}, che degrada il cAMP in AMP. La riduzione di cAMP inattiva la PKA e, di conseguenza, la HSL, bloccando la degradazione dei trigliceridi.
\item \textbf{Sintesi proteica}: AKT attiva \textbf{mTOR}, un regolatore centrale che promuove la traduzione proteica e la trascrizione genica (via S6K1 e 4EBP), fondamentale per l'ipertrofia muscolare. Lo stesso pathway è attivato dall'\textbf{IGF-1}.
\end{itemize}
\begin{nota}[mTOR e insulino-resistenza]
mTOR è attivato non solo da insulina/IGF-1, ma anche da amminoacidi (leucina) e lipidi. Tuttavia, tramite un feedback negativo (via S6K1 che inibisce IRS1), un'attivazione cronica di mTOR per eccesso nutrizionale riduce la sensibilità cellulare all'insulina, contribuendo all'insulino-resistenza.
\end{nota}
\subsection{Glucocorticoidi (cortisolo)}
Il cortisolo, secreto dalla corteccia surrenale, ha effetti catabolici più lenti (via genomica). Nel fegato stimola gli enzimi della gluconeogenesi. Nel muscolo inibisce la sintesi di glicogeno e proteica (attivando la miostatina, inibitore di mTOR), favorendo il catabolismo per fornire substrati al fegato e inducendo insulino-resistenza periferica.

\section{Regolazione allosterica e sensori energetici}
Mentre gli ormoni agiscono in minuti o ore, la cellula risponde a cambiamenti energetici immediati tramite \textbf{effettori allosterici}: molecole (substrati, prodotti, ATP, AMP, $\mathrm{Ca^{2+}}$) che si legano a siti regolatori degli enzimi, modificandone l'attività in pochi secondi.
\subsection{Regolazione di enzimi chiave}
\begin{itemize}
\item \textbf{Glicogeno fosforilasi}: attivata direttamente dal $\mathrm{Ca^{2+}}$ rilasciato dal reticolo sarcoplasmatico durante la contrazione e dall'AMP (indicatore di deficit energetico). Questo meccanismo anticipa l'arrivo dell'adrenalina.
\item \textbf{Piruvato deidrogenasi (PDH)}: controlla l'ingresso del piruvato nel ciclo di Krebs. È inattivata da acetil-CoA e NADH (segnali di abbondanza) e attivata da piruvato, $\mathrm{Ca^{2+}}$ e $\mathrm{Mg^{2+}}$ (segnali di contrazione/richiesta).
\item \textbf{Fosfofruttochinasi-1 (PFK-1)}: principale punto di controllo della glicolisi. Inibita da ATP e citrato; attivata da ADP e AMP.
\end{itemize}
\subsection{Il ruolo dell'AMP e dell'AMPK}
Quando l'ATP viene consumato rapidamente, l'enzima \textbf{adenilato chinasi} (o miochinasi) converte 2 ADP in 1 ATP + 1 \textbf{AMP}. L'aumento di AMP è il segnale più sensibile di carenza energetica.
L'AMP attiva la \textbf{Proteina Chinasi attivata dall'AMP (AMPK)}, il principale sensore energetico cellulare:
\begin{itemize}
\item \textbf{Promuove il catabolismo}: aumenta la captazione di glucosio (GLUT4) e acidi grassi (CD36), stimola il trasporto mitocondriale (CPT1) e il metabolismo ossidativo.
\item \textbf{Inibisce l'anabolismo}: blocca le vie biosintetiche e agisce come \textbf{antagonista fisiologico di mTOR}.
\end{itemize}
Questa regolazione garantisce che, in condizioni di stress energetico, la cellula priorizzi la produzione di ATP a scapito della crescita e della sintesi.
% =========================================================
% Fine Capitolo 15
% =========================================================

% =========================================================
% CAPITOLO 16: Il muscolo scheletrico e la contrazione muscolare
% =========================================================
\chapter{Il muscolo scheletrico e la contrazione muscolare}
\section{Introduzione al tessuto muscolare}
Il tessuto muscolare rappresenta circa il \textbf{40\% della massa corporea totale} ed è deputato alla conversione di energia chimica (ATP) in energia meccanica (movimento). Nell'organismo umano si distinguono tre tipi di muscolo:
\begin{itemize}
\item \textbf{Muscolo scheletrico}: striato e volontario, responsabile della locomozione e del movimento consapevole;
\item \textbf{Muscolo cardiaco}: striato e involontario, specializzato nella contrazione ritmica del cuore;
\item \textbf{Muscolo liscio}: non striato e involontario, presente nelle pareti di organi cavi e vasi sanguigni.
\end{itemize}
Oltre alla funzione motoria, il muscolo scheletrico funge da \textbf{riserva proteica} mobilizzabile in condizioni di digiuno prolungato o stress metabolico.

\section{Struttura della fibra muscolare}
Le cellule del muscolo scheletrico, dette \textbf{miociti} o \textbf{fibre muscolari}, presentano caratteristiche peculiari:
\begin{itemize}
\item Forma allungata e cilindrica, con lunghezza variabile da pochi millimetri a diversi centimetri;
\item \textbf{Multinucleazione}: derivano dalla fusione di numerosi mioblasti durante lo sviluppo;
\item \textbf{Postmitoticità}: non si dividono dopo la differenziazione; il numero di fibre è determinato geneticamente prima della nascita.
\end{itemize}
\begin{definizione}[Unità motoria]
Un \textbf{motoneurone} innerva un gruppo variabile di fibre muscolari (da poche unità nei muscoli oculari a migliaia nei muscoli degli arti), formando un'\textbf{unità motoria}. Tutte le fibre di un'unità motoria si contraggono in modo coordinato.
\end{definizione}
\subsection{Organizzazione subcellulare}
La fibra muscolare è delimitata dalla membrana plasmatica, detta \textbf{sarcolemma}, che presenta invaginazioni perpendicolari chiamate \textbf{tubuli T}. Questi permettono al potenziale d'azione di propagarsi rapidamente all'interno della fibra.
Il citoplasma (\textbf{sarcoplasma}) è occupato quasi interamente dalle \textbf{miofibrille}, strutture cilindriche parallele che rappresentano l'apparato contrattile. Nel sarcoplasma sono inoltre presenti:
\begin{itemize}
\item \textbf{Mioglobina}: proteina di deposito dell'ossigeno;
\item \textbf{Mitocondri}: suddivisi in subsarcolemmatici (SS, sotto la membrana) e intramiofibrillari (IMF, a contatto con le miofibrille);
\item \textbf{Gocce lipidiche e granuli di glicogeno}: riserve energetiche locali;
\item \textbf{Reticolo sarcoplasmatico (RS)}: rete di canali membranosi derivata dal reticolo endoplasmatico, specializzata nell'immagazzinamento e nel rilascio di ioni \(\mathrm{Ca^{2+}}\).
\end{itemize}
\begin{nota}[Triade]
In corrispondenza dei tubuli T, il RS si espande a formare le \textbf{cisterne terminali}. L'insieme di un tubulo T flanched da due cisterne terminali costituisce la \textbf{triade}, struttura fondamentale per l'accoppiamento eccitazione-contrazione.
\end{nota}

\section{Le miofibrille e il sarcomero}
Le miofibrille sono composte principalmente da due proteine contrattili:
\begin{itemize}
\item \textbf{Miosina}: forma i \textbf{filamenti spessi}. Ogni molecola presenta una coda ad elica e una testa globulare con attività ATPasica, capace di legare l'actina (ponte trasversale).
\item \textbf{Actina}: polimerizza a formare l'\textbf{actina F} (filamentosa), struttura base dei \textbf{filamenti sottili}.
\end{itemize}
I filamenti sottili contengono anche proteine regolatrici:
\begin{itemize}
\item \textbf{Tropomiosina}: proteina allungata che avvolge il filamento di actina, mascherando i siti di legame per la miosina a riposo;
\item \textbf{Troponina}: complesso ternario (TnI, TnT, TnC) che ancora la tropomiosina e media la regolazione calcio-dipendente.
\end{itemize}
\subsection{Architettura del sarcomero}
Il \textbf{sarcomero} è l'unità funzionale ripetuta delle miofibrille, delimitata dalle \textbf{linee Z} e con una lunghezza a riposo di circa \(2{,}0\text{--}2{,}5\,\mu\mathrm{m}\). La sua struttura presenta bande caratteristiche:
\begin{itemize}
\item \textbf{Banda A} (scura): regione di sovrapposizione tra filamenti spessi e sottili;
\item \textbf{Banda I} (chiara): contiene solo filamenti sottili;
\item \textbf{Linea M}: centro del sarcomero, dove le code di miosina sono ancorate dalla miomesina;
\item \textbf{Linea Z}: punto di ancoraggio dei filamenti sottili, composta da \(\alpha\)-actinina, nebulina e desmina.
\end{itemize}
La proteina \textbf{titina} (la più grande del proteoma umano) connette la linea Z alla linea M, stabilizzando i filamenti spessi e conferendo elasticità passiva al sarcomero.

\section{Accoppiamento eccitazione-contrazione}
La contrazione muscolare è innescata da un segnale nervoso trasmesso attraverso la \textbf{giunzione neuromuscolare}.
\subsection{Trasmissione sinaptica}
Il motoneurone rilascia il neurotrasmettitore \textbf{acetilcolina (ACh)} nella fessura sinaptica. Il legame dell'ACh ai recettori nicotinici del sarcolemma causa un flusso ionico che depolarizza la membrana, generando un \textbf{potenziale d'azione}. Questo si propaga lungo il sarcolemma e nei tubuli T.
\subsection{Rilascio del calcio}
La depolarizzazione dei tubuli T attiva il \textbf{recettore della rianodina} (canale del \(\mathrm{Ca^{2+}}\) del RS), provocando il rilascio massivo di ioni \(\mathrm{Ca^{2+}}\) nel sarcoplasma. La concentrazione citosolica di \(\mathrm{Ca^{2+}}\) aumenta da \(\approx 10^{-7}\,\mathrm{M}\) a \(\approx 10^{-5}\,\mathrm{M}\).
\subsection{Meccanismo molecolare della contrazione}
\begin{enumerate}
\item Il \(\mathrm{Ca^{2+}}\) si lega alla \textbf{troponina C}, inducendo un cambiamento conformazionale nel complesso troponina-tropomiosina.
\item La tropomiosina si sposta, esponendo i siti di legame per la miosina sull'actina.
\item Inizia il \textbf{ciclo dei ponti trasversali}:
\begin{itemize}
    \item \textbf{Distacco}: il legame dell'ATP alla testa di miosina ne causa il distacco dall'actina;
    \item \textbf{Idrolisi}: l'ATP è idrolizzato ad ADP + \(\mathrm{P_i}\); l'energia rilasciata "carica" la testa di miosina in una conformazione ad alta energia;
    \item \textbf{Legame}: la testa di miosina si lega a un nuovo sito sull'actina (se esposto);
    \item \textbf{Colpo di forza}: il rilascio di \(\mathrm{P_i}\) innesca il movimento della testa di miosina, che trascina il filamento sottile verso il centro del sarcomero;
    \item \textbf{Rigor}: dopo il rilascio dell'ADP, la miosina rimane saldamente legata all'actina fino al legame di un nuovo ATP.
\end{itemize}
\end{enumerate}
\begin{nota}[Rigor mortis]
In assenza di ATP (es. dopo la morte), actina e miosina rimangono permanentemente legate, causando la rigidità cadaverica.
\end{nota}
\subsection{Rilassamento}
La contrazione termina quando:
\begin{itemize}
\item L'acetilcolina è degradata dall'\textbf{acetilcolinesterasi}, interrompendo lo stimolo nervoso;
\item La pompa \textbf{SERCA} (\(\mathrm{Ca^{2+}}\)-ATPasi del reticolo sarcoplasmatico) riassorbe attivamente il \(\mathrm{Ca^{2+}}\) nel RS, consumando ATP;
\item La diminuzione del \(\mathrm{Ca^{2+}}\) citosolico permette alla tropomiosina di mascherare nuovamente i siti di legame sull'actina.
\end{itemize}

\section{Ruoli dell'ATP nella contrazione}
L'ATP svolge tre funzioni essenziali:
\begin{enumerate}
\item \textbf{Distacco della miosina}: necessario per interrompere il legame actina-miosina;
\item \textbf{Energia meccanica}: l'idrolisi dell'ATP fornisce l'energia per il colpo di forza;
\item \textbf{Trasporto attivo}: alimenta la pompa SERCA per il riassorbimento del \(\mathrm{Ca^{2+}}\).
\end{enumerate}
Nonostante l'elevato consumo durante l'esercizio, i livelli di ATP nel muscolo sono mantenuti pressoché costanti (\(\approx 5\,\mathrm{mmol/kg}\)) grazie a sistemi di rigenerazione rapidi (fosfocreatina, glicolisi, fosforilazione ossidativa).

\section{Proprietà meccaniche della contrazione}
\subsection{Contrazione rapida e tetanica}
Un singolo impulso nervoso (1--2 ms) genera una \textbf{contrazione rapida} (twitch) della durata di 20--200 ms, preceduta da un breve periodo di latenza necessario per la segnalazione del \(\mathrm{Ca^{2+}}\).
Durante un'attività sostenuta, impulsi ripetuti a frequenza elevata causano la \textbf{sommazione} delle contrazioni. Se la frequenza è sufficiente a impedire il rilassamento completo, si ottiene un \textbf{tetano fuso}, con tensione massima continua.
\subsection{Relazione forza-lunghezza}
La forza generata da una fibra dipende dalla lunghezza iniziale del sarcomero:
\begin{itemize}
\item A \(2{,}2\,\mu\mathrm{m}\) (lunghezza ottimale) la sovrapposizione actina-miosina è massima, così come la forza;
\item A lunghezze minori o maggiori, la ridotta sovrapposizione o l'interferenza sterica diminuiscono il numero di ponti trasversali possibili.
\end{itemize}
\subsection{Relazione forza-velocità}
Esiste un compromesso inverso tra forza e velocità di contrazione:
\begin{itemize}
\item Carichi elevati \(\rightarrow\) contrazione lenta, forza massima;
\item Carichi nulli \(\rightarrow\) contrazione veloce, forza nulla.
\end{itemize}
La \textbf{potenza} (lavoro per unità di tempo, misurata in Watt) è data dal prodotto forza \(\times\) velocità e raggiunge il massimo a circa \(1/3\) della velocità massima (\(V_{\max}\)).

\section{Classificazione delle fibre muscolari}
Le fibre scheletriche si classificano in base alle proprietà contrattili e metaboliche:
\begin{table}[htbp]
\centering
\caption{Caratteristiche delle fibre muscolari scheletriche}
\label{tab:fibre-muscolari}
\renewcommand{\arraystretch}{1.2}
\begin{tabular}{@{} p{0.2\textwidth} p{0.38\textwidth} p{0.38\textwidth} @{}}
\toprule
\textbf{Tipo} & \textbf{Proprietà contrattili} & \textbf{Metabolismo energetico} \\
\midrule
\begin{tabular}[t]{@{}l@{}}\textbf{Tipo I}\\(lente, rosse, ossidative)\end{tabular} &
\begin{tabular}[t]{@{}l@{}}
\textbullet{} Miosina-ATPasi lenta\\
\textbullet{} Contrazione lenta, forza ridotta\\
\textbullet{} Alta resistenza alla fatica
\end{tabular} &
\begin{tabular}[t]{@{}l@{}}
\textbullet{} Molti mitocondri, alta mioglobina\\
\textbullet{} Fosforilazione ossidativa predominante\\
\textbullet{} Utilizzo preferenziale di lipidi e carboidrati\\
\textbullet{} Ideali per endurance
\end{tabular} \\
\midrule
\begin{tabular}[t]{@{}l@{}}\textbf{Tipo IIx}\\(rapide, bianche, glicolitiche)\end{tabular} &
\begin{tabular}[t]{@{}l@{}}
\textbullet{} Miosina-ATPasi veloce\\
\textbullet{} Contrazione rapida, forza elevata\\
\textbullet{} Bassa resistenza alla fatica
\end{tabular} &
\begin{tabular}[t]{@{}l@{}}
\textbullet{} Pochi mitocondri, poca mioglobina\\
\textbullet{} Glicolisi anaerobica e fosfocreatina\\
\textbullet{} Ideali per sprint e sollevamento pesi\\
\textbullet{} Tendenza all'ipertrofia
\end{tabular} \\
\midrule
\begin{tabular}[t]{@{}l@{}}\textbf{Tipo IIa}\\(rapide, ossidative)\end{tabular} &
\begin{tabular}[t]{@{}l@{}}
\textbullet{} Proprietà intermedie\\
\textbullet{} Resistenza moderata alla fatica
\end{tabular} &
\begin{tabular}[t]{@{}l@{}}
\textbullet{} Metabolismo misto (aerobico + anaerobico)\\
\textbullet{} Reclutate a intensità sub-massimali
\end{tabular} \\
\bottomrule
\end{tabular}
\end{table}
\begin{definizione}[Distribuzione delle fibre]
La percentuale di ciascun tipo di fibra varia geneticamente tra individui e muscoli. Atleti di endurance presentano una predominanza di fibre di tipo I (65--75\%), mentre velocisti mostrano una prevalenza di tipo II (70--80\%). L'allenamento può indurre adattamenti fenotipici (es. transizione IIx \(\rightarrow\) IIa) ma non cambia il numero totale di fibre.
\end{definizione}

\section{Reclutamento delle unità motorie}
Il sistema nervoso regola la forza muscolare tramite due meccanismi:
\begin{enumerate}
\item \textbf{Codifica di frequenza}: aumento della frequenza di scarica dei motoneuroni per potenziare la contrazione delle fibre già attive;
\item \textbf{Reclutamento ordinato}: attivazione sequenziale delle unità motorie in base alla soglia di eccitazione.
\end{enumerate}
Secondo il \textbf{principio di Henneman}, le unità motorie di tipo I (piccole, bassa soglia) sono reclutate per prime durante attività a bassa intensità. All'aumentare della richiesta di forza, vengono progressivamente reclutate le unità di tipo IIa e infine IIx (grandi, alta soglia).
\begin{esempio}[Applicazione pratica]
Durante una corsa a bassa intensità, sono attivate principalmente le fibre di tipo I. In uno sprint finale o in un sollevamento massimale, il reclutamento si estende alle fibre rapide di tipo II, garantendo la massima potenza in tempi brevi.
\end{esempio}
% =========================================================
% Fine Capitolo 16
% =========================================================

% =========================================================
% CAPITOLO 17: Biochimica dell'esercizio fisico ad alta intensità
% =========================================================
\chapter{Biochimica dell'esercizio fisico ad alta intensità}
\section{Introduzione: VO\textsubscript{2}max e deficit di ossigeno}
Un parametro fondamentale per quantificare l'intensità di un esercizio è il \textbf{VO\textsubscript{2}}, ovvero il volume di ossigeno consumato per minuto. All'aumentare dell'intensità, aumenta il consumo di ossigeno fino al raggiungimento del \textbf{massimo consumo di ossigeno (VO\textsubscript{2}max)}, valore individuale dipendente da fattori genetici e dal grado di allenamento. Da un punto di vista metabolico, il VO\textsubscript{2}max rappresenta la massima quantità di ATP ottenibile dal muscolo tramite metabolismo ossidativo.
\begin{definizione}[Deficit di ossigeno]
Nelle fasi iniziali di un esercizio a intensità costante ($\leq$ VO\textsubscript{2}max), il consumo effettivo di ossigeno è inferiore al fabbisogno teorico per la produzione di ATP richiesta. Questa differenza è nota come \textbf{deficit di ossigeno}. In questa fase, l'energia proviene da meccanismi anaerobici (fosfocreatina, glicolisi) che compensano l'incapacità del sistema cardiovascolare di fornire immediatamente l'ossigeno necessario e di attivare completamente la piruvato deidrogenasi (PDH).
\end{definizione}
In esercizi ad alta intensità, lo stato stazionario non viene mai raggiunto: parte del fabbisogno energetico è \textbf{sovramassimale} (>100\% VO\textsubscript{2}max) e il deficit si accumula. La differenza totale tra fabbisogno teorico e consumo effettivo è definita \textbf{MAOD} (Maximal Accumulated Oxygen Deficit) ed è un indicatore della capacità anaerobica dell'atleta.

\section{EPOC: il debito di ossigeno post-esercizio}
Al termine dell'esercizio, il consumo di ossigeno rimane elevato rispetto al valore basale per un periodo variabile. Questo consumo extra è definito \textbf{EPOC} (Excess Post-exercise Oxygen Consumption) e serve a:
\begin{itemize}
\item Rigenerare le scorte di \textbf{fosfocreatina (PCr)};
\item Smaltire il \textbf{lattato} accumulato (riconversione in glucosio o ossidazione);
\item Ripristinare le riserve di ossigeno legate a emoglobina e mioglobina;
\item Sostenere l'aumento di temperatura corporea e la ventilazione.
\end{itemize}
La durata e l'entità dell'EPOC dipendono dall'intensità e dalla durata dello sforzo.

\section{Caratteristiche dell'esercizio ad alta intensità (HIE)}
Un esercizio ad alta intensità (\textbf{HIE}, High-Intensity Exercise) è caratterizzato da:
\begin{itemize}
\item \textbf{Breve durata}: da pochi secondi (salto, lancio) fino a circa 1 minuto (sprint 60--200 m);
\item \textbf{Produzione energetica anaerobica}: l'ATP è fornito principalmente da fosforilazione diretta (PCr) e glicolisi anaerobica;
\item \textbf{Reclutamento massivo delle fibre di tipo IIx}: fibre rapide, glicolitiche, ad alta potenza ma bassa resistenza.
\end{itemize}
I processi aerobici contribuiscono in misura significativa solo per durate superiori ai 30 s.
\subsection{Variazioni dei metaboliti durante HIE}
Studi su sprint di 30 s hanno mostrato variazioni rapide dei metaboliti muscolari:
\begin{itemize}
\item \textbf{PCr}: si riduce a $\approx 34\%$ del valore basale già dopo 6 s;
\item \textbf{Glicogeno}: diminuisce significativamente ($\approx 70\%$ del valore iniziale);
\item \textbf{ATP}: rimane stabile o scende lievemente ($>50\text{--}60\%$), poiché la sua deplezione totale è incompatibile con la vita cellulare (rigor mortis).
\end{itemize}
Il contributo relativo della glicolisi anaerobica aumenta con la durata dell'esercizio, mentre quello della PCr diminuisce progressivamente.

\section{Produzione e ruolo del lattato}
Durante un HIE, la glicolisi produce grandi quantità di piruvato che, in carenza di ossigeno sufficiente, non può essere convertito in acetil-CoA. Viene quindi ridotto a \textbf{lattato} dalla lattato deidrogenasi (LDH), rigenerando il NAD\textsuperscript{+} necessario per continuare la glicolisi.
\begin{definizione}[Soglia del lattato]
L'intensità alla quale le concentrazioni di lattato iniziano ad aumentare esponenzialmente è detta \textbf{soglia anaerobica} o soglia del lattato. Oltre questa soglia, l'accumulo di lattato e ioni H\textsuperscript{+} abbassa il pH intracellulare, inibendo enzimi chiave (ATPasi miosinica, SERCA) e contribuendo all'insorgenza della fatica.
\end{definizione}
Il lattato non è un semplice prodotto di scarto:
\begin{enumerate}
\item \textbf{Navetta del lattato}: può essere rilasciato dalle fibre di tipo IIx, captato dalle fibre ossidative (tipo I) o dal fegato e riconvertito in piruvato per l'ossidazione o in glucosio tramite gluconeogenesi (\textbf{Ciclo di Cori}).
\item \textbf{Effetto Bohr}: il co-trasporto di lattato e H\textsuperscript{+} tramite i trasportatori MCT-1/MCT-4 acidifica lo spazio extracellulare, riducendo l'affinità dell'emoglobina per l'O\textsubscript{2} e favorendone il rilascio ai tessuti attivi.
\end{enumerate}

\section{Contributo del metabolismo aerobico}
Anche in esercizi anaerobici, il metabolismo ossidativo contribuisce alla produzione di ATP:
\begin{itemize}
\item Nei primi 10 s di uno sprint: $\approx 3\%$ del totale;
\item Media nei 30 s (test di Wingate): $15\text{--}30\%$.
\end{itemize}
Le riserve di ossigeno legate a emoglobina e mioglobina possono sostenere temporaneamente il metabolismo aerobico anche in assenza di apporto esterno. Con l'aumentare della durata dell'HIE, il contributo aerobico cresce, ma la \textbf{velocità di produzione di ATP} (turnover) diminuisce, causando un calo progressivo della potenza meccanica.

\section{Regolazione metabolica durante HIE}
La necessità di generare ATP in pochi secondi è incompatibile con la regolazione ormonale (troppo lenta). La modulazione avviene tramite \textbf{effettori allosterici} che rispondono istantaneamente alle variazioni cellulari.
\subsection{Ruolo dell'AMP e dell'AMPK}
L'idrolisi massiva di ATP durante uno sforzo intenso aumenta i livelli di ADP. L'enzima \textbf{adenilato chinasi} (o miochinasi) catalizza:
\[
\mathrm{2\,ADP \rightleftharpoons ATP + AMP}
\]
L'\textbf{AMP}, normalmente presente a concentrazioni bassissime, aumenta rapidamente e funge da potente attivatore allosterico di:
\begin{itemize}
\item \textbf{Glicogeno fosforilasi} (glicogenolisi);
\item \textbf{Fosfofruttochinasi-1 (PFK-1)} (glicolisi);
\item \textbf{AMPK} (Proteina Chinasi attivata dall'AMP).
\end{itemize}
\begin{definizione}[AMPK]
L'\textbf{AMPK} è il principale sensore di carenza energetica cellulare. La sua attivazione:
\begin{itemize}
\item \textbf{Promuove il catabolismo}: aumenta la captazione di glucosio (GLUT4) e acidi grassi (CD36), stimola il trasporto mitocondriale degli acidi grassi (CPT1) e il metabolismo ossidativo;
\item \textbf{Inibisce l'anabolismo}: blocca le vie biosintetiche e agisce come antagonista fisiologico di mTOR.
\end{itemize}
Questa regolazione garantisce che, in condizioni di stress energetico, la cellula priorizzi la produzione immediata di ATP.
\end{definizione}
\subsection{Altri effettori allosterici}
\begin{itemize}
\item \textbf{Ca\textsuperscript{2+}}: rilasciato dal reticolo sarcoplasmatico durante la contrazione, attiva direttamente la glicogeno fosforilasi chinasi e la piruvato deidrogenasi fosfatasi, anticipando la risposta ormonale.
\item \textbf{ATP/ADP}: alti livelli di ATP inibiscono la PFK-1 (feedback negativo), mentre l'ADP la attiva.
\end{itemize}

\section{Metabolismo delle purine e soglia dell'ammoniaca}
L'AMP accumulato viene smaltito tramite il \textbf{ciclo dei nucleotidi purinici (PNC)}:
\[
\mathrm{AMP \xrightarrow{\text{AMP deaminasi}} IMP + NH_4^+}
\]
L'ammoniaca (\(\mathrm{NH_3}\)/\(\mathrm{NH_4^+}\)) prodotta diffonde nel plasma. Ad intensità superiori al 100\% del VO\textsubscript{2}max, si osserva un aumento esponenziale dell'ammoniaca plasmatica, definibile come \textbf{soglia dell'ammoniaca}. L'iperammoniemia contribuisce all'insorgenza della fatica centrale e periferica, interferendo con la trasmissione neuromuscolare e il ciclo di Krebs.

\section{Effetti dello stato nutrizionale}
In un HIE di breve durata (30 s), le scorte di glicogeno si riducono del 25--35\%, rendendo l'impatto dello stato nutrizionale meno evidente rispetto a esercizi prolungati. Tuttavia:
\begin{itemize}
\item Il glicogeno diminuisce più rapidamente nelle fibre di tipo IIx rispetto alle tipo I;
\item Una dieta povera di carboidrati anticipa l'insorgenza della fatica;
\item I benefici di una dieta iper-glucidica rispetto a una normo-glucidica sono meno chiari per sforzi molto brevi, ma diventano significativi per ripetizioni multiple o recuperi incompleti.
\end{itemize}

\section{Supporti nutrizionali ergogenici nell'HIE}
Un supporto ergogenico è una sostanza in grado di migliorare la performance. Per l'HIE, le evidenze più solide supportano:
\begin{itemize}
\item \textbf{Creatina}: aumenta le scorte muscolari di fosfocreatina, migliorando la performance in sforzi ripetuti ad alta intensità e potenzialmente la risposta ipertrofica all'allenamento.
\item \textbf{Alcalinizzatori} (bicarbonato, citrato di sodio): tamponano l'acidosi metabolica indotta dal lattato, ritardando la fatica in sforzi di 1--7 minuti.
\item \textbf{Caffeina}: aumenta la disponibilità di Ca\textsuperscript{2+} per la contrazione e modula la percezione dello sforzo; gli effetti sono significativi ma variabili in base al genotype metabolico individuale.
\item \textbf{\(\beta\)-alanina}: precursore della carnosina, principale tampone intracellulare contro l'acidità; l'integrazione cronica aumenta la capacità tampone muscolare e migliora la performance in HIE.
\end{itemize}
\begin{nota}[Individualità della risposta]
L'efficacia degli integratori ergogenici dipende da fattori individuali (genetica, stato di allenamento, dieta). La caffeina, ad esempio, mostra risposte variabili in base alla velocità di metabolizzazione (genotipo CYP1A2).
\end{nota}

\section{Sintesi funzionale}
\begin{itemize}
\item L'HIE (1 s -- 2 min) crea un deficit di ossigeno non compensato dal metabolismo aerobico;
\item Recluta preferenzialmente fibre di tipo IIx, ad alta potenza ma bassa resistenza;
\item L'energia proviene da ATP, PCr, glicogenolisi e glicolisi anaerobica; il metabolismo aerobico contribuisce significativamente solo oltre i 30 s;
\item La regolazione metabolica è mediata da effettori allosterici (AMP, Ca\textsuperscript{2+}, ATP/ADP) e dall'AMPK, non da ormoni;
\item L'accumulo di lattato e ammoniaca contribuisce alla fatica;
\item Integratori come creatina, alcalinizzanti, caffeina e \(\beta\)-alanina possono migliorare la performance in HIE se utilizzati correttamente.
\end{itemize}
% =========================================================
% Fine Capitolo 17
% =========================================================

% =========================================================
% CAPITOLO 18: Biochimica dell'esercizio fisico di endurance
% =========================================================
\chapter{Biochimica dell'esercizio fisico di endurance}
\section{Definizione e caratteristiche dell'endurance}
Un esercizio di \textbf{endurance} (o aerobico) è un'attività fisica costante di durata variabile da pochi minuti a diverse ore. Sono inclusi sia esercizi sulla media distanza (mezzofondo, nuoto) che sulla lunga distanza (maratona, ciclismo). Esercizi superiori alle 4--6 ore sono definiti di \textbf{ultra-endurance}.
L'obiettivo è mantenere la massima potenza possibile per periodi prolungati, con intensità considerata stazionaria (al netto di cambi di ritmo strategici).
Molti studi utilizzano il ciclismo come modello sperimentale per:
\begin{itemize}
\item Facilità di misurazione della spesa energetica e dell'intensità;
\item Posizione stabile delle braccia, utile per prelievi ematici durante l'esercizio;
\item Reclutamento specifico del muscolo vasto laterale, facilmente campionabile tramite biopsia.
\end{itemize}
\begin{nota}[Reclutamento delle fibre]
Negli esercizi di endurance vengono reclutate principalmente le fibre muscolari di \textbf{tipo I} (lente, ossidative) e \textbf{tipo IIa} (rapide, ossidative), poiché non è richiesta la produzione di ATP alla velocità massima tipica delle fibre IIx glicolitiche.
\end{nota}

\section{Fonti energetiche nell'endurance}
La produzione di ATP per la contrazione muscolare è sostenuta principalmente dal \textbf{metabolismo ossidativo} di lipidi e carboidrati. Le fonti possono essere extramuscolari o intramuscolari:
\begin{table}[htbp]
\centering
\caption{Fonti energetiche nell'esercizio di endurance}
\label{tab:fonti-endurance}
\renewcommand{\arraystretch}{1.2}
\begin{tabular}{@{} p{0.45\textwidth} p{0.45\textwidth} @{}}
\toprule
\textbf{Fonti extramuscolari} & \textbf{Fonti intramuscolari} \\
\midrule
\textbullet{} Acidi grassi liberi (FFA) plasmatici, mobilizzati dai TAG del tessuto adiposo & \textbullet{} Glucosio prodotto dal glicogeno muscolare \\
\textbullet{} Glucosio plasmatico, fornito dal glicogeno epatico & \textbullet{} Acidi grassi ricavati dai TAG intramuscolari \\
\bottomrule
\end{tabular}
\end{table}
Le proteine svolgono un ruolo marginale nella sintesi di ATP durante l'endurance, diventando significative solo in condizioni di digiuno prolungato o deplezione glicogenica estrema.

\section{Regolazione metabolica}
Il metabolismo ossidativo è più complesso di quello anaerobico e presenta molteplici punti di controllo per bilanciare l'utilizzo di carboidrati e lipidi. La regolazione avviene a più livelli:
\begin{itemize}
\item \textbf{Regolazione ormonale}: catecolammine (adrenalina, noradrenalina) e insulina modulano l'attività enzimatica tramite cascate di segnalazione (es. via cAMP-PKA).
\item \textbf{Effettori allosterici}: ADP, AMP, Ca\textsuperscript{2+}, Pi rispondono istantaneamente alle variazioni del bilancio energetico cellulare.
\item \textbf{Regolazione a lungo termine}: l'espressione genica degli enzimi metabolici viene modulata dall'allenamento cronico (adattamenti trascrizionali).
\end{itemize}
I principali fattori che influenzano il bilanciamento dei substrati sono:
\begin{enumerate}
\item Intensità dell'esercizio;
\item Durata dell'esercizio;
\item Stato nutrizionale (disponibilità di glicogeno, assunzione di carboidrati);
\item Livello di allenamento (adattamenti metabolici).
\end{enumerate}

\section{Effetto dell'intensità sull'utilizzo dei substrati}
Studi tramite calorimetria indiretta, traccianti isotopici e biopsie muscolari hanno quantificato il contributo delle varie fonti energetiche in funzione dell'intensità.
\begin{definizione}[Relazione intensità-s substrati]
All'aumentare dell'intensità dell'esercizio:
\begin{itemize}
\item Aumenta la spesa energetica totale;
\item Aumenta il contributo dei \textbf{carboidrati} alla produzione di ATP;
\item Diminuisce il contributo dei \textbf{lipidi}.
\end{itemize}
\end{definizione}
Con l'aumentare dell'intensità, inizia ad aumentare la produzione di lattato (\textbf{soglia aerobica}). La produzione di lattato cresce lentamente in funzione dell'utilizzo di glucosio, fino al raggiungimento della \textbf{soglia anaerobica} (o soglia del lattato), oltre la quale si osserva un aumento esponenziale dei livelli di lattato dovuto al reclutamento progressivo delle fibre di tipo IIx.
Per questo motivo, negli sport di endurance è fondamentale mantenersi \textbf{sotto la soglia anaerobica}, dove l'ATP resta a carico del catabolismo ossidativo e l'accumulo di lattato è minimo.
\subsection{Metabolismo dei carboidrati: glicogenolisi}
Con l'aumentare dell'intensità aumenta la degradazione del glicogeno muscolare. Curiosamente, questo non sembra dovuto all'aumento della forma attiva fosforilata (a) della glicogeno fosforilasi, che anzi può diminuire ad alta intensità, probabilmente a causa della riduzione del pH associata all'esercizio intenso.
Un ruolo preponderante sembra averlo l'accumulo di \textbf{fosfato inorganico (Pi)} che si libera in seguito all'idrolisi dell'ATP in ADP. Il Pi è un substrato della glicogeno fosforilasi, che quindi può lavorare più velocemente. Anche l'ADP e l'AMP liberi, che aumentano durante l'esercizio, sono regolatori allosterici dell'enzima attivo e contribuiscono alla glicogenolisi.
\subsection{Metabolismo dei carboidrati: assorbimento e utilizzo del glucosio}
Aumentando l'intensità, aumenta anche l'assorbimento del glucosio ematico a causa di:
\begin{itemize}
\item Aumento del flusso sanguigno muscolare;
\item Reclutamento di un numero maggiore di fibre.
\end{itemize}
All'interno della cellula, l'accumulo di ADP, AMP e Pi attiva la \textbf{PFK-1} (fosfofruttochinasi-1), accelerando la glicolisi e l'utilizzo del glucosio.
L'esercizio aumenta anche la traslocazione del trasportatore \textbf{GLUT4} sulla membrana plasmatica. Tuttavia, questa traslocazione non sembra proporzionale all'intensità, suggerendo che GLUT4 non sia il principale responsabile dell'aumento dell'uptake glucidico in funzione dell'intensità.
\subsection{Metabolismo dei carboidrati: ossidazione del piruvato}
Il piruvato prodotto dalla glicolisi, per essere ossidato nei mitocondri, deve essere prima convertito in \textbf{acetil-CoA} dalla \textbf{piruvato deidrogenasi (PDH)}.
All'aumentare dell'intensità, l'attività della PDH aumenta, grazie:
\begin{itemize}
\item All'aumento del suo substrato, il piruvato;
\item All'aumento del Ca\textsuperscript{2+} citosolico, che attiva la PDH fosfatasi (defosforilazione attivante).
\end{itemize}
\subsection{Metabolismo dei lipidi: disponibilità di acidi grassi}
Il contributo dei lipidi alla produzione di ATP diminuisce con l'aumentare dell'intensità, in particolare sopra il 65\% del VO\textsubscript{2}max.
Un esercizio al 65\% fa aumentare la concentrazione ematica sia degli acidi grassi liberi (FFA) che del glicerolo, prodotti della lipolisi. Aumentando ulteriormente l'intensità (es. 85\% VO\textsubscript{2}max), il glicerolo continua ad aumentare (indice che la lipolisi non è rallentata), ma i FFA diminuiscono.
La riduzione di FFA non è quindi dovuta a una riduzione della lipolisi, ma più probabilmente a un \textbf{ridotto afflusso di sangue al tessuto adiposo}, che limita la disponibilità di albumina per il trasporto degli acidi grassi. In queste condizioni, gli acidi grassi possono rimanere intrappolati negli adipociti.
\subsection{Metabolismo dei lipidi: ingresso mitocondriale}
La ridotta disponibilità di FFA nel sangue contribuisce solo in parte alla ridotta ossidazione lipidica, perché la concentrazione plasmatica rimane comunque superiore alla velocità di ossidazione. Questo suggerisce alterazioni anche nell'ingresso degli FFA nelle cellule e nella loro ossidazione.
Aumentando l'intensità, diminuisce l'assorbimento degli acidi grassi a catena lunga (LCFA, es. oleato) ma non di quelli a catena media (MCFA, es. octanoato).
\begin{itemize}
\item I LCFA necessitano del legame con la \textbf{carnitina} e del trasportatore \textbf{CPT1} per entrare nel mitocondrio.
\item I MCFA possono entrare direttamente nel mitocondrio senza carnitina.
\end{itemize}
L'elemento limitante non è quindi l'ingresso nella cellula (mediato da CD36, FATP, FABP), ma l'\textbf{ingresso nel mitocondrio} tramite CPT1.
\subsection{Disponibilità di carnitina e accumulo di acetil-CoA}
Il piruvato prodotto dalla glicolisi attiva la PDH, che trasforma il piruvato in acetil-CoA. Se l'acetil-CoA viene prodotto più velocemente di quanto possa essere smaltito dal ciclo di Krebs, si accumula e inibisce la PDH (feedback negativo), bloccando il metabolismo ossidativo e aumentando la produzione di lattato.
Per evitare questo, la \textbf{carnitina} può legare l'acetil-CoA e formare \textbf{acetil-carnitina}:
\[
\text{Acetil-CoA} + \text{Carnitina} \rightleftharpoons \text{Acetil-carnitina} + \text{CoA}
\]
Questo meccanismo tampona l'acetil-CoA in eccesso e permette di proseguire l'esercizio ad intensità elevata.
\begin{nota}[Effetto sulla beta-ossidazione]
Il lato negativo di questo processo è la diminuzione della disponibilità di \textbf{carnitina libera}, necessaria per il trasporto degli acidi grassi a catena lunga nel mitocondrio tramite CPT1. Il risultato finale è un minore ingresso degli acidi grassi e una ridotta ossidazione lipidica ad alta intensità.
\end{nota}
L'assunzione cronica di carnitina (es. 2 g/giorno per 12--24 settimane) è in grado di aumentarne i livelli muscolari e di ridurre la produzione di lattato negli esercizi ad intensità elevata, con un conseguente miglioramento delle performance.
\subsection{Effetto del pH su CPT1}
Un esercizio ad intensità elevata causa un abbassamento del pH muscolare. L'enzima \textbf{CPT1} è sensibile a variazioni di pH e un suo abbassamento ne può ridurre l'attività, contribuendo ulteriormente a limitare l'ingresso degli acidi grassi nei mitocondri e quindi la loro ossidazione.

\section{Effetti della durata dell'esercizio}
Al contrario rispetto all'intensità, l'aumento della \textbf{durata} di un esercizio ad intensità costante e moderata (60\% VO\textsubscript{2}max) provoca:
\begin{itemize}
\item Un \textbf{aumento} dell'ossidazione dei lipidi;
\item Una \textbf{diminuzione} dell'ossidazione dei carboidrati.
\end{itemize}
Nello specifico, aumenta l'utilizzo degli acidi grassi plasmatici a scapito del glicogeno muscolare (e anche dei trigliceridi intramuscolari).
La ridotta disponibilità di glicogeno riduce la quantità di piruvato prodotta dalla glicolisi. La riduzione del piruvato inattiva la PDH (promuovendo l'attivazione della piruvato deidrogenasi chinasi, PDK), portando a una ridotta ossidazione dei carboidrati e favorendo indirettamente l'ossidazione lipidica.

\section{Effetti dello stato nutrizionale}
\subsection{Carico di carboidrati}
Bassi livelli di glicogeno promuovono la produzione di adrenalina, che stimola la lipolisi e l'ossidazione lipidica.
Aumentare le scorte di glicogeno muscolare con una dieta ricca di carboidrati:
\begin{itemize}
\item Aumenta la glicogenolisi e l'utilizzo del glicogeno, indipendentemente dall'intensità;
\item Non interferisce con la captazione del glucosio ematico;
\item Stimola sia l'attività della glicogeno fosforilasi (substrato) sia la piruvato deidrogenasi, aumentando l'ossidazione dei carboidrati a scapito dei lipidi.
\end{itemize}
Alte concentrazioni di glicogeno prima di un esercizio aumentano la produzione di acetil-CoA dal piruvato, che viene legato dalla carnitina, riducendo ulteriormente la disponibilità di carnitina libera per il trasporto degli acidi grassi.
\subsection{Assunzione di carboidrati pre-esercizio}
L'assunzione di carboidrati poco prima di un allenamento induce un aumento delle prestazioni rispetto a una condizione di digiuno, grazie alla maggiore disponibilità di glucosio.
Tuttavia, in un'ottica di perdita di massa grassa, l'assunzione di carboidrati pre-esercizio limita fortemente la lipolisi e l'ossidazione dei lipidi, in parte dovuta all'azione anti-lipolitica dell'\textbf{insulina}. Gli effetti negativi dell'ingestione di carboidrati sulla lipolisi possono persistere per molte ore dopo l'alimentazione.
\begin{definizione}[Indice glicemico]
Visti gli effetti dell'insulina, l'assunzione di carboidrati a \textbf{basso indice glicemico (LGI)} prima di un allenamento comporta una maggiore disponibilità di FFA e una maggiore ossidazione lipidica rispetto all'assunzione di carboidrati ad alto indice glicemico (HGI), che provocano un picco insulinico più marcato.
\end{definizione}
\subsection{Adattamento lipidico (train low, compete high)}
La diminuzione delle scorte di carboidrati endogeni è un fattore limitante le performance negli esercizi di endurance prolungati (>90 min). Forzare l'organismo all'utilizzo dei grassi (\textbf{adattamento lipidico}) può ridurre l'uso del glicogeno e migliorare le performance.
Questo ha portato allo studio di strategie di allenamento con bassi livelli di carboidrati per poi aumentarli in gara (\textit{train low, compete high}).
Dopo 5 giorni di dieta per l'adattamento lipidico (alto contenuto di grassi e bassi CHO):
\begin{itemize}
\item L'ossidazione dei lipidi risulta aumentata sia al giorno 6 (giorno di ripristino dei carboidrati) sia il giorno successivo;
\item Tra le cause: aumento dei TAG intramuscolari e della lipasi HSL;
\item L'adattamento lipidico riduce l'uso del glicogeno muscolare, ma non altera l'utilizzo del glucosio ematico.
\end{itemize}

\section{Sintesi funzionale}
\begin{itemize}
\item Un esercizio di endurance è un'attività costante della durata da 5 minuti a 4 ore (o più per ultra-endurance).
\item Le fonti di energia sono il metabolismo ossidativo di lipidi e carboidrati, provenienti da riserve intra- ed extramuscolari.
\item Il bilanciamento tra lipidi e carboidrati dipende da intensità, durata, stato nutrizionale e livello di allenamento.
\item All'aumentare dell'intensità c'è una diminuzione dell'utilizzo dei lipidi a favore dei carboidrati, a causa di limitazioni nel trasporto mitocondriale (CPT1, carnitina) e nella disponibilità di FFA.
\item All'aumentare della durata, il metabolismo lipidico è favorito dalla deplezione del glicogeno e dalla conseguente inattivazione della PDH.
\item Protocolli di adattamento lipidico (\textit{train low}) riducono l'ossidazione dei carboidrati e aumentano l'ossidazione lipidica anche in seguito a reintegrazione dei carboidrati.
\item La regolazione metabolica avviene tramite ormoni (adrenalina, insulina) ed effettori allosterici (AMP, Ca\textsuperscript{2+}, Pi) che rispondono alle esigenze energetiche in tempo reale.
\end{itemize}
% =========================================================
% Fine Capitolo 18
% =========================================================

% =========================================================
% CAPITOLO 19: Biochimica del gioco del calcio
% =========================================================
\chapter{Biochimica del gioco del calcio}
\section{Il calcio come esercizio intermittente ad alta intensità}
Il calcio, analogamente ad altri sport di squadra come rugby e hockey, è caratterizzato da un profilo di attività \textbf{intermittente ad alta intensità} (HIIE, High-Intensity Intermittent Exercise). Durante una partita, fasi di sforzo massimale o sovramassimale (sprint, salti, contrasti) si alternano a periodi di attività a bassa intensità (camminata, jogging) o di riposo completo.
Questa alternanza rende la regolazione metabolica particolarmente complessa, poiché coinvolge simultaneamente:
\begin{itemize}
\item \textbf{Sistemi anaerobici}: fosfocreatina (PCr) e glicolisi anaerobica, per gli sforzi esplosivi;
\item \textbf{Sistemi aerobici}: ossidazione di carboidrati e lipidi, per il recupero tra gli sforzi e il mantenimento dell'attività a bassa intensità.
\end{itemize}
I modelli sperimentali per lo studio del calcio utilizzano protocolli lavoro/riposo che cercano di replicare il profilo metabolico reale, sebbene con alcune limitazioni in termini di specificità.
\begin{definizione}[HIIE]
Un esercizio \textbf{HIIE} è caratterizzato da ripetuti brevi sforzi ad alta intensità (da pochi secondi a 1-2 minuti) intervallati da periodi di recupero attivo o passivo. Il calcio rappresenta un esempio fisiologico di HIIE prolungato (90 minuti).
\end{definizione}

\section{Produzione di energia negli sprint ripetuti}
Studi condotti su modelli di sprint ripetuti (es. 10 sprint massimali di 6 s su cicloergometro, con 30 s di recupero) hanno fornito informazioni fondamentali sul metabolismo muscolare durante attività intermittenti.
\subsection{Declino della potenza e contributo energetico}
Durante il primo sprint, la potenza media prodotta è massima e il contributo della PCr alla produzione di ATP è prossimo al 50\%. Al decimo sprint si osserva:
\begin{itemize}
\item Una riduzione del 65\% della potenza media prodotta;
\item Una diminuzione significativa del contributo della glicolisi anaerobica;
\item Un aumento relativo del contributo del metabolismo aerobico.
\end{itemize}
\begin{table}[htbp]
\centering
\caption{Produzione di ATP durante sprint ripetuti (dati medi)}
\label{tab:sprint-atp}
\renewcommand{\arraystretch}{1.2}
\begin{tabular}{@{} p{0.3\textwidth} p{0.3\textwidth} p{0.3\textwidth} @{}}
\toprule
\textbf{Parametro} & \textbf{1° Sprint} & \textbf{10° Sprint} \\
\midrule
Produzione ATP & 14,9 mmol/kg/s & 5,3 mmol/kg/s \\
Contributo PCr & $\approx$ 50\% & Ridotto \\
Contributo glicolisi & Elevato & Ridotto \\
Contributo aerobico & Minore & Aumentato \\
\bottomrule
\end{tabular}
\end{table}
\subsection{Analisi dei metaboliti muscolari}
Biopsie del muscolo vasto laterale prelevate prima e dopo gli sprint hanno rivelato:
\begin{itemize}
\item \textbf{PCr}: il recupero di 30 s non è sufficiente per rigenerare completamente le scorte di fosfocreatina;
\item \textbf{Glicogenolisi}: la degradazione del glicogeno è ridotta nel decimo sprint rispetto al primo, nonostante alti livelli di adrenalina che dovrebbero stimolarla;
\item \textbf{Lattato}: non si accumula ulteriormente negli sprint finali, indicando un minore flusso glicolitico.
\end{itemize}
Questi dati suggeriscono che, in esercizi intermittenti prolungati, il \textbf{metabolismo aerobico} diventa progressivamente più importante per sostenere la produzione di ATP.

\section{Regolazione enzimatica durante l'HIIE}
L'attività degli enzimi chiave del metabolismo glucidico è modulata dinamicamente durante gli sprint ripetuti.
\subsection{Glicogeno fosforilasi}
La percentuale di glicogeno fosforilasi nella forma attiva (fosforilata, forma \textit{a}) aumenta rapidamente nei primi 15 s del primo sprint, probabilmente in risposta a:
\begin{itemize}
\item Aumento del \(\mathrm{Ca^{2+}}\) citosolico (rilasciato dal reticolo sarcoplasmatico);
\item Accumulo di \(\mathrm{P_i}\) (fosfato inorganico) dall'idrolisi dell'ATP;
\item Aumento dell'AMP (indicatore di deficit energetico).
\end{itemize}
Nei successivi sprint, questa attivazione è attenuata, probabilmente a causa dell'acidificazione intracellulare (\(\mathrm{H^+}\)) che inibisce l'enzima.
\subsection{Piruvato deidrogenasi (PDH)}
La PDH, enzima che converte il piruvato in acetil-CoA per l'ingresso nel ciclo di Krebs, mostra un profilo di attivazione diverso:
\begin{itemize}
\item Nel primo sprint, l'attività della PDH aumenta progressivamente, raggiungendo il massimo negli ultimi 15 s;
\item Nel terzo sprint, l'attivazione è più rapida e sostenuta.
\end{itemize}
L'attivazione è favorita dal \(\mathrm{Ca^{2+}}\) (che attiva la PDH fosfatasi) e, paradossalmente, dagli ioni \(\mathrm{H^+}\) che stimolano la defosforilazione attivante.
\begin{nota}[Bilancio piruvato-lattato]
Nel primo sprint, la velocità di produzione del piruvato supera la capacità della PDH di ossidarlo: il piruvato in eccesso viene convertito in lattato, che si accumula. Negli sprint successivi, la maggiore attività della PDH e la ridotta glicogenolisi prevengono l'accumulo di piruvato e, di conseguenza, di lattato.
\end{nota}

\section{Effetti del rapporto lavoro/riposo}
La durata della fase di sforzo e del recupero influenza significativamente la performance e lo stress metabolico.
\begin{table}[htbp]
\centering
\caption{Effetto della distanza dello sprint sulla performance (recupero fisso: 30 s)}
\label{tab:sprint-distanza}
\renewcommand{\arraystretch}{1.2}
\begin{tabular}{@{} p{0.25\textwidth} p{0.35\textwidth} p{0.35\textwidth} @{}}
\toprule
\textbf{Distanza} & \textbf{Numero sprint} & \textbf{Effetto sulla performance} \\
\midrule
15 m & 40 & Nessun calo significativo \\
30 m & 20 & Calo del 6\% \\
40 m & 15 & Calo del 12\% \\
\bottomrule
\end{tabular}
\end{table}
Sprint più lunghi (30-40 m) inducono un maggiore stress metabolico, evidenziato da:
\begin{itemize}
\item Aumento di \textbf{acido urico} e \textbf{ipoxantina} plasmatici (metaboliti della degradazione delle purine, indicativi di consumo di ATP e accumulo di AMP);
\item Maggiore accumulo di lattato ematico.
\end{itemize}
Un recupero più lungo riduce lo stress metabolico: dopo uno sprint di 6 s, 30 s di recupero non rigenerano completamente la PCr; dopo uno sprint di 30 s, 4 minuti di recupero permettono un ripristino quasi completo.

\section{Recupero della fosfocreatina e capacità aerobica}
La rigenerazione della PCr avviene nei mitocondri durante il recupero, utilizzando ATP prodotto dalla fosforilazione ossidativa. Questo processo dipende dall'afflusso di sangue e ossigeno al muscolo.
\begin{definizione}[Implicazioni per l'allenamento]
Atleti con una \textbf{elevata capacità aerobica} (\(\mathrm{VO_2max}\)) mostrano una più rapida resintesi della PCr tra gli sforzi intermittenti. Questo permette di mantenere prestazioni più elevate per periodi prolungati, ritardando l'affaticamento.
\end{definizione}

\section{Utilizzo dei substrati energetici nell'HIIE}
Confrontando esercizi intermittenti ad alta intensità con esercizi continui a intensità media equivalente (stessa spesa energetica totale), si osserva:
\begin{itemize}
\item \textbf{Maggiore ossidazione dei carboidrati} nell'HIIE;
\item \textbf{Minore ossidazione dei lipidi} nell'HIIE.
\end{itemize}
Questo è probabilmente dovuto all'elevato flusso glicolitico che riduce la disponibilità di carnitina libera (sequestrata come acetil-carnitina), limitando l'ingresso degli acidi grassi a catena lunga nei mitocondri tramite CPT1.
\subsection{Ruolo dei lipidi nel calcio}
Nonostante la predominanza dei carboidrati, i lipidi svolgono un ruolo importante, specialmente nel \textbf{secondo tempo}:
\begin{itemize}
\item Si osserva un aumento degli acidi grassi liberi (FFA) plasmatici, indicativo di aumentata lipolisi nel tessuto adiposo;
\item Questo è favorito dalla combinazione di: minore intensità media, aumento delle catecolamine (adrenalina) e diminuzione dell'insulina (per deplezione dei carboidrati).
\end{itemize}

\section{Effetti dello stato nutrizionale}
\subsection{Livelli di glicogeno muscolare}
Bassi livelli iniziali di glicogeno muscolare sono associati a:
\begin{itemize}
\item Aumento della lipolisi (FFA e glicerolo plasmatici elevati);
\item Minori livelli di glucosio ematico;
\item \textbf{Ridotta percentuale di sforzi ad alta intensità} durante la partita.
\end{itemize}
Questo evidenzia l'importanza di adeguate riserve di glicogeno per sostenere la performance in sport intermittenti.
\subsection{Assunzione di carboidrati pre-esercizio}
In condizioni di digiuno, l'ossidazione lipidica è maggiore rispetto all'assunzione di carboidrati prima dell'esercizio. Tuttavia, a differenza degli esercizi continui sottomassimali, \textbf{l'indice glicemico} del pasto pre-esercizio (alto vs. basso) non influenza significativamente:
\begin{itemize}
\item L'ossidazione lipidica;
\item La performance fisica nell'HIIE.
\end{itemize}
Eventuali piccole differenze sono probabilmente annullate dall'assunzione di carboidrati \textit{durante} l'esercizio, pratica comune nel calcio.
\subsection{Assunzione di carboidrati durante l'esercizio}
L'ingestione di carboidrati (es. maltodestrine al 7,5\%, 1 ml/kg ogni 12 min) durante un esercizio HIIE prolungato:
\begin{itemize}
\item Mantiene elevati i livelli di insulina e glucosio ematico;
\item Riduce i livelli di glicerolo e FFA (minore lipolisi);
\item \textbf{Aumenta la capacità di sostenere l'esercizio} (es. 158 min vs. 131 min fino all'esaurimento);
\item Non altera significativamente l'utilizzo complessivo del glicogeno muscolare.
\end{itemize}
Il beneficio principale sembra derivare dalla maggiore disponibilità di glucosio ematico, che supporta l'elevata velocità di ossidazione glucidica richiesta negli sforzi intermittenti.
\begin{esempio}[Miglioramento delle abilità tecniche]
L'integrazione con carboidrati durante l'esercizio migliora le performance in test specifici per il calcio:
\begin{itemize}
\item Agilità: riduzione del tempo di esecuzione;
\item Dribbling: maggiore precisione e velocità;
\item Precisione del tiro: migliore accuratezza.
\end{itemize}
Questi benefici sono attribuibili al mantenimento della funzione neuromuscolare e cognitiva, supportata da adeguati livelli di glucosio.
\end{esempio}

\section{Sintesi funzionale}
\begin{itemize}
\item Il calcio è un esercizio \textbf{intermittente ad alta intensità} (HIIE) che richiede l'integrazione di sistemi energetici anaerobici e aerobici.
\item La produzione di ATP deriva inizialmente da PCr e glicolisi anaerobica; con il protrarsi dell'esercizio, il contributo aerobico aumenta progressivamente.
\item Il recupero della PCr dipende dalla durata del riposo e dalla capacità aerobica dell'atleta.
\item L'ossidazione dei carboidrati è predominante nell'HIIE; i lipidi contribuiscono maggiormente nelle fasi tardive e a bassa intensità.
\item Adeguati livelli di glicogeno muscolare sono essenziali per sostenere gli sforzi ad alta intensità.
\item L'assunzione di carboidrati \textit{durante} l'esercizio migliora la capacità di performance prolungata e le abilità tecniche specifiche, probabilmente tramite il mantenimento della glicemia e della funzione neuromuscolare.
\end{itemize}
% =========================================================
% Fine Capitolo 19
% =========================================================

% =========================================================
% CAPITOLO 20: Adattamenti metabolici all'allenamento
% =========================================================

\chapter{Adattamenti metabolici all'allenamento}

\section{Introduzione e segnali cellulari}
La ripetizione cronica dell'esercizio fisico induce specifici adattamenti metabolici a livello delle fibre muscolari e dell'intero organismo. Tali adattamenti sono reversibili e dipendono strettamente dal tipo di stimolo (alta intensità vs. endurance), dalla durata e dalla frequenza dell'allenamento.
Il muscolo scheletrico è un tessuto altamente plastico: lo stress meccanico, metabolico, ossidativo e ipossico derivante dalla contrazione attiva vie di trasduzione del segnale che portano all'aumento dell'espressione e dell'attività di enzimi e proteine specifiche. I principali segnali coinvolti sono:
\begin{itemize}
\item \textbf{Stress meccanico}: accorciamenti e stiramenti ripetuti attivano meccanosensori come canali ionici attivati dallo stiramento (ingresso di \(\mathrm{Na^+}\) e \(\mathrm{Ca^{2+}}\)) e integrine di membrana, che trasducono segnali di crescita e ipertrofia.
\item \textbf{Molecole segnale}: ormoni (adrenalina, cortisolo, GH) e fattori di crescita (IGF-I, MGF) promuovono risposte anaboliche, mentre inibitori della crescita come la miostatina vengono soppressi.
\item \textbf{Stress metabolico}: l'accumulo di AMP attiva la chinasi AMPK, sensore energetico che, tramite l'attivazione di PGC-1\(\alpha\), promuove la biogenesi mitocondriale e il catabolismo lipidico.
\item \textbf{Stress ossidativo}: l'esercizio aumenta moderatamente le specie reattive dell'ossigeno (ROS), che a livelli controllati fungono da molecole di segnalazione attivando vie adattative (NRF2, MAPK).
\item \textbf{Stress ipossico}: il rapido consumo di \(\mathrm{O_2}\) nei mitocondri attiva fattori indotti dall'ipossia (HIFs), stimolando l'espressione di geni per GLUT4, mioglobina ed enzimi glicolitici.
\end{itemize}

\section{Adattamenti all'allenamento ad alta intensità (HIT)}
L'allenamento ad alta intensità migliora la capacità di svolgere esercizi massimali o sovramassimali, con adattamenti circoscritti principalmente al metabolismo anaerobico e alla struttura contrattile:
\begin{itemize}
\item \textbf{Metabolismo anaerobico}: aumento dell'attività della glicolisi anaerobica e della creatina chinasi (CK); lieve incremento delle riserve di creatina/fosfocreatina (PCr).
\item \textbf{Metabolismo glucidico}: modesto aumento delle riserve di glicogeno intramuscolare e miglioramento della capacità di tamponare l'acidosi intracellulare (aumento della carnosina).
\item \textbf{Metabolismo proteico e ionico}: aumento del contenuto proteico muscolare e up-regulation della pompa \(\mathrm{Na^+/K^+}\)-ATPasi.
\end{itemize}
La pompa \(\mathrm{Na^+/K^+}\)-ATPasi contrasta l'accumulo extracellulare di \(\mathrm{K^+}\) (che depolarizza il sarcolemma riducendo l'eccitabilità), preservando la generazione della forza durante sforzi ripetuti.

\subsection{Ipertrofia muscolare e cellule satellite}
Uno degli effetti principali del HIT è l'aumento del diametro delle fibre muscolari (\textbf{ipertrofia}), mediato dall'attivazione della via di segnalazione AKT-mTOR. Questa via, stimolata da IGF-I/MGF, amminoacidi (leucina) e stress meccanico, attiva la sintesi proteica necessaria alla costruzione di nuove miofibrille, reticolo sarcoplasmatico e citoscheletro.
Le \textbf{cellule satellite} (staminali muscolari) giocano un ruolo cruciale: in risposta al danno da esercizio, si attivano, proliferano e possono fondersi con fibre preesistenti, aumentando il numero di nuclei per fibra (\textbf{iperplasia funzionale}). Questo aumento del rapporto sarcoplasma/nuclei precede e supporta l'ipertrofia. Inoltre, si ipotizza l'esistenza di una \textbf{memoria muscolare}: anche dopo periodi di inattività e perdita di massa, i mionuclei aggiunti persistono, accelerando il recupero ipertrofico alla ripresa dell'allenamento.

\section{Adattamenti all'allenamento di endurance (ET)}
L'ET induce profondi adattamenti fisiologici, tra cui l'aumento del \(\mathrm{VO_{2max}}\) e lo spostamento della soglia del lattato verso intensità maggiori, indice di un migliorato metabolismo ossidativo.

\subsection{Biogenesi mitocondriale e angiogenesi}
L'adattamento più evidente è l'aumento della densità mitocondriale (per numero e dimensione), guidato dal co-attivatore trascrizionale PGC-1\(\alpha\), attivato da AMPK e vie MAPK. Parallelamente, lo stress ipossico e metabolico stimola l'\textbf{angiogenesi}, aumentando la capillarizzazione e migliorando l'apporto di \(\mathrm{O_2}\) e nutrienti.

\subsection{Metabolismo dei carboidrati}
Nei soggetti allenati si osserva un \textbf{risparmio di glicogeno} muscolare durante esercizi sottoposti, dovuto a ridotti livelli di ADP, AMP e \(\mathrm{P_i}\) (attivatori della glicogeno fosforilasi). Questo posticipa la fatica. Conseguentemente, si riducono la produzione di piruvato e lattato e l'attività della PDH.
Si osserva inoltre un aumento dell'espressione di GLUT4, che non tanto durante l'esercizio, ma soprattutto a riposo/post-esercizio, migliora la captazione di glucosio (sensibilità insulinica) e l'attività della glicogeno sintasi, favorendo un maggiore accumulo di glicogeno.

\subsection{Metabolismo dei lipidi}
L'ET sposta il metabolismo ossidativo dai carboidrati verso i lipidi. Questo non deriva da un aumento della lipolisi sistemica, ma da:
\begin{itemize}
\item Maggiore assorbimento degli acidi grassi plasmatici, mediato dall'up-regulation dei trasportatori di membrana FABPm e FAT/CD36.
\item Maggiore utilizzo dei trigliceridi intramuscolari (IMTG), favorito dall'aumento della lipasi ATGL.
\item Aumentata attività di CPT1 e degli enzimi della \(\beta\)-ossidazione e del ciclo di Krebs, che potenziano la capacità ossidativa lipidica.
\end{itemize}

\subsection{Metabolismo delle proteine}
L'ET aumenta l'espressione della BCKAD (deidrogenasi degli \(\alpha\)-chetoacidi a catena ramificata), migliorando il potenziale ossidativo dei BCAA. Tuttavia, durante l'esercizio, l'allenamento riduce l'attivazione di BCKAD (probabilmente via aumento della BCKAD chinasi), preservando i BCAA e limitando il catabolismo proteico durante lo sforzo.

\section{Adattamenti all'allenamento intermittente ad alta intensità (HIIT)}
L'HIIT produce adattamenti sovrapponibili a quelli dell'ET ad alto volume, nonostante la minore durata totale dello stimolo:
\begin{itemize}
\item Riduzione dell'utilizzo di glicogeno e della produzione di lattato.
\item Aumento della vascolarizzazione e della biogenesi mitocondriale (via PGC-1\(\alpha\)).
\item Potenziamento del metabolismo ossidativo (aumento di GLUT4, FABPm, CD36).
\end{itemize}
Inoltre, l'HIIT incrementa specificamente l'espressione dei trasportatori del lattato (MCT) e della pompa \(\mathrm{Na^+/K^+}\)-ATPasi, riducendo l'accumulo plasmatico di \(\mathrm{K^+}\) e migliorando la clearance del lattato.

\section{Disallenamento e atrofia muscolare}
L'inattività prolungata o l'immobilizzazione causano una rapida perdita degli adattamenti metabolici e una riduzione della massa muscolare (\textbf{atrofia}). Questo processo è mediato dall'attivazione di ligasi ubiquitina-specifiche come \textbf{Atrogin-1} e \textbf{MuRF-1}, che promuovono la degradazione proteica via proteasoma.
L'atrofia da disuso è fisiologica e reversibile con la ripresa dell'attività. Al contrario, l'atrofia patologica (da denervazione, ischemia o distrofie) porta a morte cellulare e perdita permanente di fibre. In questi casi, l'esercizio fisico adattato rimane utile per rallentare il declino e preservare la funzionalità residua.

\section{Sintesi funzionale}
\begin{itemize}
\item Gli adattamenti metabolici sono specifici per il tipo di allenamento e reversibili in caso di inattività.
\item L'HIT migliora la capacità anaerobica, il tamponamento del pH e induce ipertrofia tramite la via mTOR e il reclutamento delle cellule satellite.
\item L'ET potenzia il metabolismo ossidativo tramite biogenesi mitocondriale (PGC-1\(\alpha\)), angiogenesi, risparmio di glicogeno e aumento dell'ossidazione lipidica (tramite trasportatori specifici e CPT1).
\item L'HIIT unisce adattamenti ossidativi simili all'ET a miglioramenti nella gestione del lattato e dell'omeostasi ionica (\(\mathrm{Na^+/K^+}\)-ATPasi).
\item Il disallenamento attiva vie proteolitiche (Atrogin-1/MuRF-1) causando atrofia, reversibile con la ripresa dell'esercizio se non di natura patologica.
\end{itemize}

% =========================================================
% Fine Capitolo 20
% =========================================================

% =========================================================
% CAPITOLO 21: Esercizio fisico e stress ossidativo
% =========================================================

\chapter{Esercizio fisico e stress ossidativo}

\section{Definizione di stress ossidativo}
Lo \textbf{stress ossidativo} è una condizione patofisiologica caratterizzata dalla perdita dell'equilibrio (omeostasi) tra la produzione di specie ossidanti, in particolare le \textbf{specie reattive dell'ossigeno} (ROS, Reactive Oxygen Species), e la capacità di difesa antiossidante della cellula. Quando la produzione di ROS supera le capacità di neutralizzazione, si instaura uno squilibrio che può danneggiare le macromolecole biologiche.

\begin{definizione}[Reazioni di ossidoriduzione]
In una reazione redox:
\begin{itemize}
\item L'\textbf{agente ossidante} acquista elettroni, ossidando la molecola vicina;
\item L'\textbf{agente riducente} cede elettroni, riducendo la molecola vicina.
\end{itemize}
Le molecole antiossidanti (es. glutatione) agiscono come agenti riducenti, venendo esse stesse ossidate per interrompere le reazioni di ossidazione a catena.
\end{definizione}

\section{Ossigeno molecolare e ROS}
L'ossigeno molecolare (\(\mathrm{O_2}\)) costituisce circa il \SI{20}{\percent} dell'atmosfera ed è essenziale per la respirazione mitocondriale. Tuttavia, è anche una molecola reattiva con potenziali effetti dannosi (\textit{paradosso dell'ossigeno}).
I ROS sono molecole altamente reattive derivanti dalla \textbf{riduzione parziale} dell'ossigeno. Non tutti i ROS sono radicali liberi (specie con elettrone spaiato), ma tutti possiedono forte potere ossidante.

\begin{table}[htbp]
\centering
\caption{Principali specie reattive dell'ossigeno}
\label{tab:ros}
\renewcommand{\arraystretch}{1.2}
\begin{tabular}{@{} p{0.28\textwidth} p{0.3\textwidth} p{0.3\textwidth} @{}}
\toprule
\textbf{Specie} & \textbf{Formula} & \textbf{Tipo} \\
\midrule
Anione superossido & \(\mathrm{O_2^{\bullet-}}\) & Radicale libero \\
Perossido di idrogeno & \(\mathrm{H_2O_2}\) & Non radicale \\
Radicale idrossilico & \(\mathrm{HO^{\bullet}}\) & Radicale libero \\
\bottomrule
\end{tabular}
\end{table}

Nel complesso IV (citocromo c ossidasi) della catena respiratoria, 4 elettroni vengono trasferiti \textit{contemporaneamente} all'\(\mathrm{O_2}\) per formare \(\mathrm{H_2O}\), evitando la produzione di ROS intermedi. Quando questo trasferimento non è coordinato, si generano le specie reattive sopra elencate.

\subsection{Fonti di ROS}
\begin{itemize}
\item \textbf{Endogene} (prodotte dalle cellule):
\begin{itemize}
\item Mitocondri (fuoriuscita di elettroni dalla catena respiratoria);
\item Perossisomi (ossidazione di acidi grassi a catena molto lunga);
\item NADPH ossidasi (NOX) del reticolo sarcoplasmatico e dei tubuli T;
\item Xantina ossidasi (degradazione delle purine);
\item Ossido nitrico sintasi (produzione di NO, una specie reattiva dell'azoto, RNS).
\end{itemize}
\item \textbf{Esogene} (fattori ambientali):
\begin{itemize}
\item Radiazioni ionizzanti e UV;
\item Farmaci e xenobiotici;
\item Inquinamento atmosferico;
\item Fumo di tabacco;
\item Alimentazione (eccesso di grassi, metalli pesanti).
\end{itemize}
\end{itemize}

\subsection{Bersagli molecolari dei ROS}
I ROS possono danneggiare tutte le macromolecole biologiche:
\begin{itemize}
\item \textbf{DNA}: ossidazione delle basi, rotture del filamento, mutazioni, accelerazione dell'invecchiamento e rischio oncogenico;
\item \textbf{Proteine}: ossidazione di residui amminoacidici (Cys, Met), formazione di ponti disolfuro aberranti, aggregazione, perdita di funzione;
\item \textbf{Lipidi}: perossidazione degli acidi grassi polinsaturi di membrana, alterazione della fluidità e della permeabilità, produzione di aldeidi tossiche (es. malondialdeide).
\end{itemize}

\section{Sistemi di difesa antiossidante}
Per prevenire o limitare i danni ossidativi, le cellule possiedono sistemi di difesa articolati.

\begin{definizione}[Classificazione degli antiossidanti]
\begin{itemize}
\item \textbf{Esogeni}: assunti con la dieta (vitamine, polifenoli, carotenoidi);
\item \textbf{Endogeni}:
\begin{itemize}
\item \textit{Enzimatici}: proteine che catalizzano la neutralizzazione dei ROS;
\item \textit{Non enzimatici}: molecole a basso peso molecolare (tioli, bilirubina, acido urico).
\end{itemize}
\end{itemize}
\end{definizione}

\subsection{Antiossidanti enzimatici}
\begin{itemize}
\item \textbf{Superossido dismutasi (SOD)}: catalizza la dismutazione dell'anione superossido in ossigeno e perossido di idrogeno:
\[
\mathrm{2\,O_2^{\bullet-} + 2\,H^+ \xrightarrow{SOD} H_2O_2 + O_2}
\]
Esistono isoforme citosoliche (Cu/Zn-SOD) e mitocondriali (Mn-SOD).
\item \textbf{Catalasi (CAT)}: localizzata nei perossisomi, degrada il perossido di idrogeno:
\[
\mathrm{2\,H_2O_2 \xrightarrow{CAT} 2\,H_2O + O_2}
\]
\item \textbf{Glutatione perossidasi (GPx)}: utilizza glutatione ridotto (GSH) per detossificare \(\mathrm{H_2O_2}\) e perossidi lipidici:
\[
\mathrm{H_2O_2 + 2\,GSH \xrightarrow{GPx} 2\,H_2O + GSSG}
\]
Il glutatione ossidato (GSSG) viene rigenerato dalla glutatione reduttasi a spese di NADPH.
\end{itemize}

\subsection{Antiossidanti non enzimatici endogeni}
\begin{itemize}
\item \textbf{Glutatione (GSH)}: tripeptide (Glu-Cys-Gly) con gruppo tiolico reattivo; principale antiossidante intracellulare;
\item \textbf{Bilirubina}: prodotto della degradazione dell'eme; a basse concentrazioni protegge i lipidi di membrana dalla perossidazione;
\item \textbf{Acido urico}: prodotto finale del metabolismo delle purine; principale antiossidante del plasma (ad alte concentrazioni può precipitare causando gotta).
\end{itemize}

\subsection{Antiossidanti esogeni (dietari)}
\begin{itemize}
\item \textbf{Vitamina C} (acido ascorbico): antiossidante idrosolubile, rigenera la vitamina E ossidata;
\item \textbf{Vitamina E} (tocoferoli): antiossidante liposolubile, protegge le membrane dalla perossidazione lipidica;
\item \textbf{Carotenoidi} (es. \(\beta\)-carotene, licopene): scavenger di radicali liberi, proteggono da danni foto-ossidativi;
\item \textbf{Polifenoli} (es. resveratrolo, quercetina): attività antiossidante e modulazione di vie di segnalazione.
\end{itemize}

\section{Esercizio fisico e produzione di ROS}
Numerosi studi dimostrano che esercizi di endurance o di breve durata ma alta intensità causano un aumento dei marcatori di stress ossidativo nel sangue e nel muscolo. Al contrario, esercizi di breve durata (\(<1\) min) a bassa intensità (\(\approx \SI{30}{\percent}\) \(\mathrm{VO_{2max}}\)) non inducono un aumento significativo.

\begin{nota}[La contrazione muscolare come fonte di ROS]
La causa principale dell'aumento di ROS durante l'esercizio è la \textbf{contrazione muscolare}, non il metabolismo mitocondriale. I mitocondri contribuiscono in misura minore rispetto ad altre fonti cellulari.
\end{nota}

\subsection{Principali fonti di ROS durante l'esercizio}
\begin{itemize}
\item \textbf{Reticolo sarcoplasmatico / Tubuli T}: contengono NADPH ossidasi (NOX) che producono anione superossido in risposta alla depolarizzazione e al \(\mathrm{Ca^{2+}}\);
\item \textbf{Perossisomi}: la contrazione attiva la xantina ossidasi (XO) durante la degradazione dell'AMP (via delle purine), producendo ROS;
\item \textbf{Ossido nitrico sintasi (NOS)}: i muscoli esprimono le isoforme neuronale (nNOS; fibre veloci) ed endoteliale (eNOS). La contrazione aumenta la produzione di ossido nitrico (\(\mathrm{NO}\)), una specie reattiva dell'azoto (RNS) che può reagire con \(\mathrm{O_2^{\bullet-}}\) formando perossinitrito (\(\mathrm{ONOO^-}\)), un potente ossidante.
\end{itemize}

\section{Effetti dei ROS sulla funzione muscolare}
\subsection{ROS e forza muscolare}
Livelli basali di ROS sono necessari al muscolo non affaticato per generare forza ottimale. Esperimenti su fibre isolate mostrano una relazione a campana:
\begin{itemize}
\item \textbf{Bassi ROS} (trattamento con antiossidanti): forza ridotta;
\item \textbf{ROS moderati}: aumento della forza;
\item \textbf{ROS elevati}: diminuzione della forza e danno cellulare.
\end{itemize}

\begin{definizione}[Meccanismi di modulazione della forza]
\begin{itemize}
\item \textbf{Effetti positivi}: regolazione dei flussi di \(\mathrm{Ca^{2+}}\) nel reticolo sarcoplasmatico e aumento della sensibilità delle miofibrille alla stimolazione;
\item \textbf{Effetti negativi}: inibizione della creatina chinasi e di enzimi glicolitici/mitocondriali chiave.
\end{itemize}
\end{definizione}

\subsection{ROS e affaticamento muscolare}
Studi su muscoli murini stimolati tetanicamente mostrano che la diminuzione della forza è accompagnata da un aumento dei ROS, in particolare dell'anione superossido. L'accumulo eccessivo di ROS contribuisce all'insorgenza della fatica attraverso:
\begin{itemize}
\item Alterazione del rilascio/riassorbimento del \(\mathrm{Ca^{2+}}\);
\item Danno ossidativo alle proteine contrattili;
\item Inibizione degli enzimi energetici.
\end{itemize}

\subsection{Danno strutturale e atrofia}
Se presenti a concentrazioni eccessive, i ROS possono:
\begin{itemize}
\item Aumentare l'espressione delle ubiquitin-ligasi \textbf{Atrogin-1} e \textbf{MuRF-1}, promuovendo la degradazione proteasomale delle miofibrille;
\item Inibire le vie di sintesi proteica (es. via AKT-mTOR), sbilanciando la proteostasi verso il catabolismo;
\item Danneggiare le \textbf{cellule satellite}, riducendo la capacità rigenerativa del muscolo o spingendole verso differenziamento adiposo.
\end{itemize}

\section{Adattamenti all'esercizio cronico: l'ormesi}
L'aumento moderato dello stress ossidativo è uno dei segnali che inducono adattamenti benefici all'esercizio fisico cronico.

\begin{definizione}[Ormesi]
L'\textbf{ormesi} è il fenomeno per cui l'esposizione cronica a dosi moderate di uno stressore (es. ROS) induce una risposta adattativa benefica, mentre dosi troppo basse o troppo alte hanno effetto nullo o deleterio.
\end{definizione}

\subsection{Meccanismi di adattamento}
Un aumento moderato di ROS attiva vie di segnalazione che promuovono la difesa cellulare:
\begin{itemize}
\item \textbf{NRF2} (Nuclear factor erythroid 2-related factor 2): si dissocia da Keap1 in risposta allo stress ossidativo, trasloca nel nucleo e induce l'espressione di geni antiossidanti (SOD, CAT, GPx, enzimi della sintesi del GSH);
\item \textbf{MAPK} (Mitogen-Activated Protein Kinases): attivano cascate di fosforilazione che regolano la trascrizione di geni adattativi;
\item \textbf{PGC-1\(\alpha\)}: co-attivatore trascrizionale indotto dai ROS, promuove la biogenesi mitocondriale e l'espressione di enzimi antiossidanti, migliorando la capacità ossidativa e la resilienza allo stress.
\end{itemize}
Se la risposta adattativa riesce a bilanciare la produzione di ROS, si mantiene l'omeostasi redox e la cellula diventa più resistente a stress futuri.

\section{Integrazione antiossidante: benefici e limiti}
L'utilizzo di integratori antiossidanti (vitamine C/E, N-acetilcisteina) durante l'esercizio ha effetti complessi e dose-dipendenti:
\begin{itemize}
\item \textbf{A livelli moderati di ROS}: gli antiossidanti possono \textbf{bloccare l'attivazione} delle vie di segnalazione (NRF2, PGC-1\(\alpha\)) che mediano gli adattamenti all'allenamento, annullando i benefici dell'esercizio cronico;
\item \textbf{A livelli elevati di ROS}: gli antiossidanti possono \textbf{proteggere} dai danni ossidativi eccessivi, ritardare la fatica e migliorare la funzionalità muscolare.
\end{itemize}

\begin{esempio}[N-acetilcisteina (NAC)]
Studi su muscoli isolati stimolati continuamente mostrano che la somministrazione di NAC (precursore del glutatione) diminuisce il declino della forza (fatica), suggerendo un ruolo protettivo in condizioni di stress ossidativo acuto e intenso.
\end{esempio}

\section{Sintesi funzionale}
\begin{itemize}
\item Lo stress ossidativo è la perdita di equilibrio tra produzione di ROS e difese antiossidanti.
\item I ROS sono specie parzialmente ridotte dell'ossigeno con forte potere ossidante, prodotte da fonti endogene (NOX, XO, mitocondri) ed esogene.
\item Le cellule possiedono sistemi antiossidanti enzimatici (SOD, CAT, GPx) e non enzimatici (GSH, vitamine).
\item L'esercizio fisico ad intensità moderata/alta aumenta la produzione di ROS, principalmente da reticolo sarcoplasmatico e perossisomi.
\item Livelli moderati di ROS sono necessari per la forza muscolare e inducono adattamenti benefici (ormesi) tramite NRF2, MAPK e PGC-1\(\alpha\).
\item Livelli eccessivi di ROS contribuiscono alla fatica, al danno strutturale e all'atrofia muscolare.
\item L'integrazione antiossidante può essere utile in condizioni di stress acuto intenso, ma può interferire con gli adattamenti cronici se usata indiscriminatamente.
\end{itemize}

% =========================================================
% Fine Capitolo 21
% =========================================================

% =========================================================
% CAPITOLO 22: Basi biochimiche della fatica
% =========================================================

\chapter{Basi biochimiche della fatica}

\section{Introduzione e classificazione}
La \textbf{fatica} può essere definita come una diminuzione delle performance fisiche, associata a un aumento (reale o percepito) della difficoltà nel sostenere un'attività, o come l'incapacità del muscolo di mantenere un determinato livello di forza durante l'esercizio. Lo studio dei suoi meccanismi è complesso a causa della natura multifattoriale del fenomeno.
L'origine della fatica può essere classificata in due categorie principali:
\begin{itemize}
\item \textbf{Fatica centrale}: origina dal sistema nervoso centrale (cervello e midollo spinale) e si manifesta come una diminuzione dell'attivazione volontaria dei muscoli (ridotta frequenza o sincronizzazione dei motoneuroni).
\item \textbf{Fatica periferica}: origina a livello o al di sotto della giunzione neuromuscolare e nel muscolo stesso, caratterizzata da una diminuzione della forza contrattile delle fibre.
\end{itemize}

\section{Fatica centrale}
La causa principale della fatica centrale è l'alterazione delle concentrazioni di neurotrasmettitori nel sistema nervoso centrale, in particolare:
\begin{itemize}
\item \textbf{Serotonina}: il suo aumento è associato a sensazioni di letargia e perdita di motivazione.
\item \textbf{Dopamina, glutammato e GABA}: le loro fluttuazioni modulano l'attivazione motoria, la percezione dello sforzo e la coordinazione neuromuscolare.
\end{itemize}
La fatica centrale è spesso influenzata da fattori metabolici periferici; ad esempio, durante l'esercizio prolungato, l'ossidazione degli amminoacidi a catena ramificata (BCAA) può alterare il rapporto tra triptofano e BCAA nel sangue, favorendo l'ingresso di triptofano nel cervello e la successiva sintesi di serotonina.

\section{Fatica periferica: meccanismi metabolici}
La fatica periferica deriva da una ridotta efficienza della giunzione neuromuscolare o da alterazioni metaboliche all'interno della fibra muscolare. I principali fattori includono l'accumulo di metaboliti (\(\mathrm{ROS}\), \(\mathrm{P_i}\), \(\mathrm{H^+}\), \(\mathrm{K^+}\)), la deplezione delle riserve di glicogeno, l'accumulo di calore e la disidratazione. Questi fattori compromettono i) la produzione di ATP; ii) la formazione/rottura dei ponti trasversi actina-miosina; iii) l'omeostasi ionica.

\subsection{Ruolo dell'ATP, della fosfocreatina e del fosfato inorganico}
Durante la fatica, i livelli di \textbf{ATP} intracellulare non scendono mai sotto il 50--60\% dei valori a riposo, valori ben al di sopra della soglia necessaria per la contrazione. Tuttavia, è probabile che si verifichino \textbf{deplezioni localizzate} di ATP in zone ad alto consumo (es. vicino alle triadi, dove operano \(\mathrm{Na^+/K^+}\)-ATPasi ed enzimi glicolitici). Una riduzione locale di ATP potrebbe innescare meccanismi protettivi che limitano il rilascio di \(\mathrm{Ca^{2+}}\) dal reticolo sarcoplasmatico, preservando la cellula dal rigor e dal danno strutturale.
La \textbf{fosfocreatina (PCr)} diminuisce drasticamente durante sforzi massimali, ma la sua disponibilità non sembra essere il fattore limitante diretto per la fatica. Al contrario, il forte accumulo di \textbf{fosfato inorganico (\(\mathrm{P_i}\))} concorre significativamente all'insorgenza della fatica. Il \(\mathrm{P_i}\) può:
\begin{itemize}
\item Ridurre il rilascio di \(\mathrm{Ca^{2+}}\) dal reticolo sarcoplasmatico;
\item Diminuire la sensibilità delle miofibrille al \(\mathrm{Ca^{2+}}\);
\item Precipitare con il calcio formando fosfato di calcio all'interno del reticolo sarcoplasmatico, riducendo ulteriormente il calcio disponibile per la contrazione.
\end{itemize}
Questi effetti sono particolarmente evidenti nelle fibre rapide di tipo IIx.

\subsection{Acidosi metabolica e ioni idrogeno}
L'accumulo di \textbf{ioni \(\mathrm{H^+}\)} (che abbassa il pH intracellulare fino a \(\approx 6.5\)) è tradizionalmente considerato una causa primaria di fatica. Il pH ridotto inibisce enzimi chiave della produzione anaerobica di energia, come la glicogeno fosforilasi, la creatina chinasi e la PFK-1. Inoltre, gli \(\mathrm{H^+}\) possono interferire con il legame del calcio alla troponina C.
Tuttavia, il ruolo dell'acidosi è fortemente dipendente dalla temperatura e dal tipo di esercizio. In sport intermittenti come il calcio, dove la temperatura muscolare può raggiungere i \SI{40}{\celsius}, l'effetto inibitorio del pH basso è attenuato, suggerendo che l'acidosi non sia il principale meccanismo di fatica in questi contesti.

\subsection{Potassio extracellulare e ROS}
Durante la contrazione, il \(\mathrm{K^+}\) viene rilasciato nello spazio extracellulare. Il suo accumulo depolarizza il sarcolemma, interferendo con la propagazione del potenziale d'azione. Gli individui allenati mostrano un minore accumulo extracellulare di \(\mathrm{K^+}\) grazie a una maggiore espressione e attività della pompa \(\mathrm{Na^+/K^+}\)-ATPasi, indotta dall'allenamento HIIT.
Le \textbf{specie reattive dell'ossigeno (ROS)} interferiscono con il riassorbimento del \(\mathrm{Ca^{2+}}\) nel reticolo sarcoplasmatico e riducono la sensibilità delle miofibrille al calcio. Inoltre, i ROS possono inibire la \(\mathrm{Na^+/K^+}\)-ATPasi, peggiorando l'accumulo extracellulare di \(\mathrm{K^+}\). La somministrazione di antiossidanti (es. N-acetilcisteina) può contrastare questi effetti e ritardare l'insorgenza della fatica.

\section{Fatica nell'esercizio di endurance}
Negli esercizi prolungati a intensità costante, il meccanismo predominante è la \textbf{ridotta disponibilità di carboidrati}. L'esaurimento del glicogeno muscolare e la successiva riduzione del glucosio ematico compromettono il rilascio di \(\mathrm{Ca^{2+}}\) dal reticolo sarcoplasmatico e riducono la capacità contrattile.
\begin{itemize}
\item L'assunzione di carboidrati (pre- o intra-esercizio) mantiene la glicemia, ritarda l'esaurimento delle riserve e migliora le prestazioni.
\item La deplezione glicogenica aumenta l'ossidazione dei BCAA. La conseguente diminuzione dei BCAA ematici favorisce l'ingresso di triptofano nel cervello, aumentando la sintesi di serotonina e promuovendo la \textbf{fatica centrale} (percezione di fatica, ridotta motivazione).
\end{itemize}
Anche la \textbf{disidratazione} gioca un ruolo cruciale: riduce il flusso ematico cutaneo e la sudorazione, compromettendo la dissipazione del calore, aumentando la temperatura corporea e riducendo la disponibilità di \(\mathrm{O_2}\) ai muscoli.

\section{Fatica nell'esercizio intermittente ad alta intensità (HIIE)}
Gli sport di squadra (es. calcio) alternano fasi ad alta e bassa intensità. La fatica in questi contesti si manifesta in due forme:
\begin{itemize}
\item \textbf{Fatica progressiva}: declino generale delle prestazioni verso la fine della gara.
\item \textbf{Fatica temporanea}: riduzione momentanea della capacità di sprint nei minuti immediatamente successivi a uno sforzo massimale.
\end{itemize}
L'analisi delle fibre muscolari post-partita mostra un'eterogeneità nella deplezione glicogenica: circa il \SI{50}{\percent} delle fibre (specialmente tipo IIx) risulta completamente o quasi svuotato, mentre il restante \SI{50}{\percent} conserva riserve. Questo squilibrio spiega la riduzione selettiva delle azioni ad alta intensità nel secondo tempo.
La \textbf{deplezione di PCr} durante gli sprint ripetuti (con recuperi medi di \(\approx \SI{70}{\second}\)) è fortemente associata alla fatica temporanea. La supplementazione con \textbf{creatina} ha dimostrato di attenuare questo fenomeno negli sport intermittenti.

\section{Adattamenti e integrazione}
L'allenamento specifico induce adattamenti che ritardano l'insorgenza della fatica:
\begin{itemize}
\item \textbf{HIIT e trasporto del lattato}: aumenta l'espressione dei trasportatori \textbf{MCT4} (esporto di lattato/\(\mathrm{H^+}\)) e \textbf{MCT1} (importo da fibre ossidative), migliorando la clearance del lattato e mitigando l'acidosi intracellulare. L'aumento di \(\mathrm{H^+}\) extracellulare favorisce inoltre il rilascio di \(\mathrm{O_2}\) dall'emoglobina (effetto Bohr).
\item \textbf{Pompa \(\mathrm{Na^+/K^+}\)-ATPasi}: l'allenamento intermittente ne aumenta la densità e l'attività, migliorando il riassorbimento del \(\mathrm{K^+}\) extracellulare e preservando l'eccitabilità di membrana.
\item \textbf{Sistemi antiossidanti endogeni}: l'esposizione cronica a stress ossidativo moderato potenzia le difese cellulari, riducendo l'interferenza dei ROS sulla contrazione.
\end{itemize}

\section{Sintesi funzionale}
\begin{itemize}
\item La fatica è un fenomeno multifattoriale, classificabile in centrale (SNC) e periferica (muscolo/giunzione neuromuscolare).
\item La fatica periferica è guidata da alterazioni metaboliche: accumulo di \(\mathrm{P_i}\), \(\mathrm{H^+}\), \(\mathrm{K^+}\) e ROS; deplezione di PCr e glicogeno.
\item L'ATP non si esaurisce mai completamente, ma deplezioni localizzate possono innescare meccanismi protettivi che riducono il rilascio di calcio.
\item Negli esercizi di endurance, la carenza di carboidrati e la disidratazione sono i principali limitanti, con ripercussioni sia periferiche che centrali (via serotonina).
\item Nell'HIIE, la fatica è eterogenea (deplezione selettiva delle fibre IIx) e presenta una componente temporanea legata al calo di PCr e all'accumulo di \(\mathrm{K^+}\).
\item L'allenamento migliora la resistenza alla fatica potenziando le pompe ioniche, i trasportatori di lattato e le difese antiossidanti.
\end{itemize}

% =========================================================
% Fine Capitolo 22
% =========================================================

\end{document}